% ===========================================================================
% ĐỀ ÁN TỐT NGHIỆP THẠC SĨ KỸ THUẬT
%
% ỨNG DỤNG MÔ HÌNH HỌC MÁY
% TRONG HỖ TRỢ KIỂM DUYỆT TÍN DỤNG
%
% Học viên   : Đồng Quốc Huy - B25CHHT028
% Chuyên ngành: Hệ thống thông tin - Mã số: 8.48.01.04
% Hướng dẫn  : TS. Đỗ Tiến Dũng
% Cơ sở đào tạo: Học viện Công nghệ Bưu chính Viễn thông
%
% BIÊN DỊCH (bắt buộc dùng XeLaTeX do có tiếng Việt và phông Times New Roman):
%     xelatex main.tex
%     xelatex main.tex      (chạy lần 2 để cập nhật mục lục và tham chiếu)
% hoặc:
%     latexmk -xelatex main.tex
% ===========================================================================

% Lưu ý về cỡ chữ: lớp ``extreport'' chỉ nhận các cỡ 8, 9, 10, 11, 12, 14, 17,
% 20pt -- không có 13pt. Nếu khai báo [13pt] ở đây, LaTeX sẽ bỏ qua tùy chọn
% (chỉ ghi một cảnh báo nhỏ trong log) và dựng toàn bộ tài liệu ở cỡ mặc định
% 10pt. Vì vậy lớp được khai báo ở 12pt, sau đó gói ``fontsize'' trong preamble
% đặt lại cỡ chữ cơ sở đúng 13pt theo quy định trình bày luận văn.
\documentclass[a4paper,12pt,oneside]{extreport}

% ===========================================================================
% Gói lệnh, định dạng và macro dùng chung cho toàn bộ đề án
% Biên dịch bằng XeLaTeX (cần cho tiếng Việt và phông Times New Roman)
% ===========================================================================

% --- Cỡ chữ cơ sở ---------------------------------------------------------
% Quy định trình bày luận văn yêu cầu cỡ chữ 13pt, nhưng lớp extreport không có
% tùy chọn này (xem ghi chú trong main.tex). Gói ``fontsize'' cho phép đặt cỡ
% chữ cơ sở tùy ý và tự tính lại toàn bộ các cỡ dẫn xuất (\large, \small,
% \footnotesize...) cùng khoảng cách dòng tương ứng.
\usepackage[fontsize=13pt]{fontsize}

% --- Ngôn ngữ và phông chữ ------------------------------------------------
\usepackage{fontspec}
% Ngôn ngữ chính là tiếng Việt (khai báo sau cùng), nạp kèm tiếng Anh cho các
% thuật ngữ và tài liệu tham khảo.
%
% Ghi chú: định dạng XeLaTeX trên máy biên dịch không nạp sẵn bộ quy tắc ngắt âm
% tiết, nên toàn bộ tài liệu được dựng ở chế độ không ngắt âm tiết. Với tiếng
% Việt đây là hành vi đúng (tiếng Việt không ngắt âm tiết trong từ). Để bù lại
% khả năng xuống dòng bị hạn chế, tham số \emergencystretch được nới rộng ở phần
% dưới, và danh mục tài liệu tham khảo được căn lề trái.
\usepackage[english,vietnamese]{babel}
\setmainfont{Times New Roman}[
  Ligatures=TeX,
  BoldFont={Times New Roman Bold},
  ItalicFont={Times New Roman Italic},
  BoldItalicFont={Times New Roman Bold Italic},
]
\setsansfont{Arial}[Ligatures=TeX]
\setmonofont{Consolas}[Scale=0.86]


% --- Trang và khổ giấy ----------------------------------------------------
% Lề theo quy định trình bày luận văn: trái 3,5cm (đóng quyển), phải 2cm
% headheight phải đủ chứa dòng đầu trang ở cỡ chữ 13pt; nếu thiếu, fancyhdr sẽ
% cảnh báo và dòng đầu trang bị đẩy chồng lên vùng chữ.
\usepackage[
  a4paper, top=2.5cm, bottom=2.5cm, left=3.5cm, right=2.0cm,
  headheight=18pt, headsep=0.6cm, footskip=1.2cm
]{geometry}
\usepackage{setspace}
\onehalfspacing                       % giãn dòng 1,5
\setlength{\parindent}{1cm}
\setlength{\parskip}{0.2\baselineskip}

% --- Toán học -------------------------------------------------------------
% Thứ tự nạp gói ở đây là bắt buộc: amsmath và mathtools phải được nạp TRƯỚC
% unicode-math, nếu không unicode-math sẽ báo lỗi định nghĩa lại ký hiệu.
\usepackage{amsmath}
\usepackage{mathtools}

% Phông toán tương thích Times.
% Mặc định LaTeX dựng công thức bằng Computer Modern, trong khi phần chữ dùng
% Times New Roman. Hai phông khác nhau rõ về nét và tỷ lệ nên công thức trông
% "lạc" giữa văn bản; ngoài ra Computer Modern chỉ có sẵn ở một số cỡ thiết kế
% cố định nên ở cỡ chữ 13pt LaTeX phải thay bằng cỡ gần nhất.
%
% TeX Gyre Termes Math là bộ ký hiệu toán OpenType dựng theo Times, dùng trực
% tiếp với XeLaTeX qua gói unicode-math. Gói này đã bao gồm toàn bộ ký hiệu của
% amssymb/amsfonts nên KHÔNG nạp thêm hai gói đó (sẽ xung đột định nghĩa).
% Tùy chọn ``math-style=ISO'' đặt chữ cái Latin và Hy Lạp trong công thức ở dạng
% nghiêng theo chuẩn ISO 80000-2.
\usepackage{unicode-math}
\setmathfont{TeX Gyre Termes Math}
\unimathsetup{math-style=ISO, bold-style=ISO}

% unicode-math cung cấp \symbf thay cho gói bm; giữ lại tên \bm cho tương thích
% với phần nội dung đã viết.
\newcommand{\bm}[1]{\symbf{#1}}

\numberwithin{equation}{chapter}

% --- Bảng biểu ------------------------------------------------------------
\usepackage{booktabs}
\usepackage{longtable}
\usepackage{tabularx}
\usepackage{multirow}
\usepackage{array}
\usepackage{makecell}
\usepackage{colortbl}
\renewcommand{\arraystretch}{1.22}
\newcolumntype{L}[1]{>{\raggedright\arraybackslash}p{#1}}
\newcolumntype{C}[1]{>{\centering\arraybackslash}p{#1}}
\newcolumntype{R}[1]{>{\raggedleft\arraybackslash}p{#1}}

% --- Hình ảnh và đồ họa ---------------------------------------------------
\usepackage{graphicx}
\graphicspath{{figures/}{../results/figures/}}
\usepackage{float}
\usepackage{subcaption}
\usepackage{xcolor}
\usepackage{tikz}
\usetikzlibrary{
  shapes.geometric, shapes.misc, arrows.meta, positioning, calc,
  fit, backgrounds, decorations.pathreplacing, patterns, chains, shadows.blur
}
\usepackage{pgfplots}
\pgfplotsset{compat=1.18}

% --------------------------------------------------------------------------
% Tự động thu nhỏ hình/sơ đồ vượt khổ chữ.
%
% Môi trường ``fitfigure'' đo chiều rộng thực tế của nội dung rồi CHỈ thu nhỏ
% khi nội dung rộng hơn khổ chữ. Đây là điểm khác biệt quan trọng so với
% \resizebox{\textwidth}{!}{...}: lệnh đó ép MỌI nội dung về đúng chiều rộng khổ
% chữ, khiến những sơ đồ vốn đã vừa bị phóng to quá mức và cỡ chữ trong hình trở
% nên không đồng nhất giữa các hình.
%
% Hiện thực bằng hộp lưu (savebox) của LaTeX, không cần thêm gói phụ thuộc.
% --------------------------------------------------------------------------
\newsavebox{\cffitbox}
\newenvironment{fitfigure}
  {\begin{lrbox}{\cffitbox}}
  {\end{lrbox}%
   % Ghi chiều rộng tự nhiên của hình ra tệp nhật ký. Nhờ đó có thể rà soát
   % những sơ đồ rộng hơn khổ chữ nhiều lần - đó là các sơ đồ sẽ bị thu nhỏ
   % mạnh và chữ bên trong trở nên khó đọc, cần thiết kế lại cho gọn hơn.
   \typeout{FITFIG \the\wd\cffitbox\space / \the\textwidth}%
   \ifdim\wd\cffitbox>\textwidth
     \resizebox{\textwidth}{!}{\usebox{\cffitbox}}%
   \else
     \usebox{\cffitbox}%
   \fi}

% --- Chú thích hình/bảng --------------------------------------------------
\usepackage[font=small,labelfont=bf,justification=centering,skip=6pt]{caption}

% --- Danh sách ------------------------------------------------------------
\usepackage{enumitem}
\setlist{itemsep=2pt, topsep=4pt, parsep=0pt}

% --- Mã nguồn -------------------------------------------------------------
\usepackage{listings}
\definecolor{codebg}{HTML}{F7F8FA}
\definecolor{codekw}{HTML}{1F4E79}
\definecolor{codestr}{HTML}{548235}
\definecolor{codecmt}{HTML}{8A8A8A}
\definecolor{codenum}{HTML}{A0A0A0}
\lstdefinestyle{cf}{
  backgroundcolor=\color{codebg},
  basicstyle=\ttfamily\scriptsize,
  keywordstyle=\color{codekw}\bfseries,
  stringstyle=\color{codestr},
  commentstyle=\color{codecmt}\itshape,
  numberstyle=\tiny\color{codenum},
  numbers=left, numbersep=7pt,
  breaklines=true, breakatwhitespace=false,
  showstringspaces=false, captionpos=b,
  frame=single, rulecolor=\color{gray!35},
  xleftmargin=14pt, framexleftmargin=14pt,
  tabsize=4, columns=fullflexible, keepspaces=true,
  literate=%
    {á}{{\'a}}1 {à}{{\`a}}1 {ả}{{\h a}}1 {ã}{{\~a}}1 {ạ}{{\d a}}1
    {ă}{{\u a}}1 {â}{{\^a}}1 {đ}{{\dj}}1 {é}{{\'e}}1 {è}{{\`e}}1
    {ê}{{\^e}}1 {í}{{\'i}}1 {ì}{{\`i}}1 {ó}{{\'o}}1 {ò}{{\`o}}1
    {ô}{{\^o}}1 {ơ}{{\o'}}1 {ú}{{\'u}}1 {ù}{{\`u}}1 {ư}{{\u'}}1
    {ý}{{\'y}}1
}
\lstset{style=cf}
\renewcommand{\lstlistingname}{Đoạn mã}

% --- Tham chiếu chéo và siêu liên kết -------------------------------------
\usepackage[
  colorlinks=true, linkcolor=black, citecolor=black,
  urlcolor=blue!55!black, filecolor=black, bookmarksnumbered=true
]{hyperref}
\usepackage{bookmark}
\usepackage{url}
\urlstyle{same}

% --- Tiêu đề chương/mục ---------------------------------------------------
\usepackage{titlesec}
\titleformat{\chapter}[display]
  {\normalfont\LARGE\bfseries\centering}
  {\chaptertitlename\ \thechapter}{14pt}{\LARGE}
\titlespacing*{\chapter}{0pt}{-18pt}{26pt}
\titleformat{\section}{\normalfont\large\bfseries}{\thesection}{0.8em}{}
\titleformat{\subsection}{\normalfont\normalsize\bfseries}{\thesubsection}{0.8em}{}
\titleformat{\subsubsection}{\normalfont\normalsize\itshape}{\thesubsubsection}{0.8em}{}
\setcounter{secnumdepth}{3}
\setcounter{tocdepth}{2}

% --- Đầu trang / chân trang -----------------------------------------------
\usepackage{fancyhdr}
\pagestyle{fancy}
\fancyhf{}
\fancyhead[R]{\small\nouppercase{\leftmark}}
\fancyfoot[C]{\thepage}
\renewcommand{\headrulewidth}{0.4pt}
\fancypagestyle{plain}{\fancyhf{}\fancyfoot[C]{\thepage}%
  \renewcommand{\headrulewidth}{0pt}}

% --- Danh mục viết tắt ----------------------------------------------------
\usepackage{longtable}

% ===========================================================================
% MÀU VÀ KIỂU TIKZ DÙNG CHUNG CHO CÁC SƠ ĐỒ KIẾN TRÚC
% ===========================================================================
\definecolor{cfBlue}{HTML}{1F4E79}
\definecolor{cfBlueL}{HTML}{DCE9F5}
\definecolor{cfOrange}{HTML}{C55A11}
\definecolor{cfOrangeL}{HTML}{FDF0E0}
\definecolor{cfGreen}{HTML}{548235}
\definecolor{cfGreenL}{HTML}{E8F0E0}
\definecolor{cfRed}{HTML}{C0504D}
\definecolor{cfRedL}{HTML}{F7E4E3}
\definecolor{cfGray}{HTML}{5A5A5A}
\definecolor{cfGrayL}{HTML}{F0F0F0}
\definecolor{cfPurple}{HTML}{7030A0}
\definecolor{cfPurpleL}{HTML}{EDE3F5}

\tikzset{
  cfnode/.style={
    draw, rounded corners=3pt, align=center, font=\small,
    minimum height=9mm, minimum width=24mm, inner sep=4pt, line width=0.7pt
  },
  source/.style={cfnode, fill=cfBlueL, draw=cfBlue},
  stream/.style={cfnode, fill=cfOrangeL, draw=cfOrange},
  compute/.style={cfnode, fill=cfGreenL, draw=cfGreen},
  storage/.style={cfnode, fill=cfRedL, draw=cfRed},
  external/.style={cfnode, fill=cfPurpleL, draw=cfPurple},
  ui/.style={cfnode, fill=cfGrayL, draw=cfGray},
  decision/.style={
    draw=cfBlue, fill=cfBlueL, diamond, aspect=2.2, align=center,
    font=\scriptsize, inner sep=2pt, line width=0.7pt
  },
  flowarrow/.style={-{Stealth[length=2.4mm]}, line width=0.75pt, draw=cfGray},
  dataflow/.style={
    -{Stealth[length=2.4mm]}, line width=0.75pt, draw=cfBlue,
    font=\scriptsize
  },
  faultflow/.style={
    -{Stealth[length=2.4mm]}, line width=0.75pt, draw=cfRed, dashed,
    font=\scriptsize
  },
  grouplabel/.style={font=\scriptsize\bfseries, text=cfGray},
  zone/.style={
    draw=gray!45, dashed, rounded corners=5pt, inner sep=7pt, line width=0.6pt
  },
}

% ===========================================================================
% MACRO TIỆN ÍCH
% ===========================================================================
% Định dạng số theo chuẩn Việt Nam (dấu phẩy thập phân)
\usepackage{siunitx}
\sisetup{
  output-decimal-marker={,},
  group-separator={.},
  group-minimum-digits=4,
  detect-all
}

% --------------------------------------------------------------------------
% Chèn hình an toàn: nếu hình chưa được sinh (chưa chạy hết thực nghiệm),
% thay bằng một khung báo thiếu thay vì làm hỏng toàn bộ quá trình biên dịch.
% --------------------------------------------------------------------------
% Tên file chứa dấu gạch dưới nên phải dùng \detokenize để in nguyên văn,
% tránh việc LaTeX hiểu "_" là ký hiệu chỉ số dưới trong chế độ toán.
% Hình được đặt trong \makebox rộng đúng bằng khổ chữ và bắt đầu bằng \noindent:
% nếu không, phần thụt đầu dòng 1cm sẽ cộng vào chiều rộng hình và làm tràn lề
% ngay cả khi hình được khai báo hẹp hơn khổ chữ.
\newcommand{\cffig}[2][\textwidth]{%
  \par\noindent
  \makebox[\linewidth][c]{%
    \IfFileExists{../results/figures/#2.pdf}
      {\includegraphics[width=#1]{#2}}
      {\fbox{\parbox[c][2.8cm][c]{0.8\textwidth}{\centering\small\itshape
        Hình \texttt{\detokenize{#2}} chưa được sinh --
        chạy lại các script thực nghiệm để tạo hình này.}}}%
  }\par
}
% Tên gọi ngắn dùng trong các chương
\providecommand{\fig}{}
\renewcommand{\fig}[2][\textwidth]{\cffig[#1]{#2}}

% Thuật ngữ tiếng Anh in nghiêng khi xuất hiện lần đầu
\newcommand{\eng}[1]{\textit{#1}}
% Tên thành phần hệ thống
\newcommand{\comp}[1]{\textsf{#1}}

% --------------------------------------------------------------------------
% Tên biến / đường dẫn mã nguồn.
%
% Chuỗi kiểu ``scripts/06_verify_consistency.py'' là một khối liền không có
% khoảng trắng nên TeX không thể ngắt dòng, gây tràn lề khi chuỗi rơi vào cuối
% dòng. Định nghĩa dưới đây tạo điểm ngắt hợp lệ ngay sau mỗi dấu gạch dưới
% bằng cách định nghĩa lại ``\_'' trong phạm vi cục bộ của lệnh.
% --------------------------------------------------------------------------
\newcommand{\code}[1]{{%
  \ttfamily\small
  \def\_{\char`\_\allowbreak}%
  #1}}

% --------------------------------------------------------------------------
% Cho phép TeX co giãn khoảng trắng nhiều hơn trước khi chấp nhận tràn lề.
% Cần thiết với văn bản tiếng Việt xen lẫn thuật ngữ tiếng Anh và tên riêng
% nước ngoài: những từ này không có quy tắc ngắt âm tiết tiếng Việt nên trình
% bày dễ bị tràn ở các dòng có từ dài.
% --------------------------------------------------------------------------
\emergencystretch=3em
\hbadness=2500        % bỏ qua cảnh báo dòng thưa ở mức chấp nhận được
\tolerance=1200

% --------------------------------------------------------------------------
% Tham số đặt hình/bảng. Mặc định của LaTeX quá dè dặt: một hình cỡ trung bình
% bị đẩy sang "trang chỉ có hình", để lại hơn nửa trang trắng. Nới rộng tỷ lệ
% hình được phép chiếm ở đầu/cuối trang và siết \floatpagefraction để hình chỉ
% chiếm trọn một trang khi thực sự đủ lớn.
% --------------------------------------------------------------------------
\renewcommand{\topfraction}{0.90}
\renewcommand{\bottomfraction}{0.70}
\renewcommand{\textfraction}{0.08}
\renewcommand{\floatpagefraction}{0.80}
\setcounter{topnumber}{3}
\setcounter{bottomnumber}{2}
\setcounter{totalnumber}{4}

% Ô nhấn mạnh kết quả chính
\usepackage[most]{tcolorbox}
\newtcolorbox{keyresult}[1][]{
  enhanced, breakable,
  colback=cfBlueL!45, colframe=cfBlue, boxrule=0.7pt, arc=3pt,
  left=8pt, right=8pt, top=6pt, bottom=6pt,
  fonttitle=\bfseries\small, title={Kết quả chính}, #1
}
\newtcolorbox{notebox}[1][]{
  enhanced, breakable,
  colback=cfOrangeL!60, colframe=cfOrange, boxrule=0.7pt, arc=3pt,
  left=8pt, right=8pt, top=6pt, bottom=6pt,
  fonttitle=\bfseries\small, title={Lưu ý}, #1
}

% --------------------------------------------------------------------------
% Trích dẫn tài liệu tham khảo
% Danh mục tài liệu được đánh số thủ công bằng môi trường enumerate (phù hợp
% quy định trình bày của cơ sở đào tạo). Lệnh \cite được định nghĩa lại để tham
% chiếu tới nhãn tương ứng, nhờ đó số thứ tự trích dẫn luôn khớp với danh mục
% ngay cả khi bổ sung hoặc sắp xếp lại tài liệu.
% --------------------------------------------------------------------------
% Hỗ trợ trích dẫn nhiều tài liệu cùng lúc: \cite{a,b,c} -> [8, 9, 10]
\makeatletter
\newcommand{\cf@citeone}[1]{\ref{ref:#1}}
\newcommand{\cf@citelist}[1]{%
  \def\cf@sep{}%
  \@for\cf@key:=#1\do{\cf@sep\cf@citeone{\cf@key}\def\cf@sep{, }}%
}
\renewcommand{\cite}[1]{[\cf@citelist{#1}]}
\makeatother

% Ký hiệu toán học hay dùng
\newcommand{\PD}{\mathrm{PD}}
\newcommand{\LGD}{\mathrm{LGD}}
\newcommand{\EAD}{\mathrm{EAD}}
\newcommand{\EL}{\mathrm{EL}}
\newcommand{\WOE}{\mathrm{WOE}}
\newcommand{\IV}{\mathrm{IV}}
\newcommand{\PSI}{\mathrm{PSI}}
\newcommand{\AUC}{\mathrm{AUC}}
\newcommand{\KS}{\mathrm{KS}}
\newcommand{\DTI}{\mathrm{DTI}}
\newcommand{\LTV}{\mathrm{LTV}}
\DeclareMathOperator*{\argmin}{arg\,min}
\DeclareMathOperator*{\argmax}{arg\,max}
\DeclareMathOperator{\sigmoid}{\sigma}
\DeclareMathOperator{\Ex}{\mathbb{E}}

% Tránh chữ mồ côi
\widowpenalty=10000
\clubpenalty=10000
\raggedbottom


% Số liệu thực nghiệm (sinh tự động bởi scripts/05_export_latex_tables.py)
% Macro số liệu thực nghiệm - SINH TỰ ĐỘNG, KHÔNG SỬA TAY
% Sinh bởi scripts/06_export_latex_tables.py

\newcommand{\ApproveRate}{29,18\%}
\newcommand{\ApproveRateDrop}{6,64}
\newcommand{\AutomationRate}{67,26\%}
\newcommand{\BadApprove}{3,43\%}
\newcommand{\BadRateReduction}{2,71}
\newcommand{\BadRateReductionPct}{44,1\%}
\newcommand{\BadReject}{34,45\%}
\newcommand{\BadReview}{10,93\%}
\newcommand{\BayesAUC}{0,9109}
\newcommand{\BestAUC}{0,8614}
\newcommand{\BestApprovalRate}{50,47\%}
\newcommand{\BestBadRateApproved}{3,52\%}
\newcommand{\BestBrier}{0,0976}
\newcommand{\BestECE}{0,0075}
\newcommand{\BestFOne}{0,4731}
\newcommand{\BestGini}{0,7229}
\newcommand{\BestKS}{0,5689}
\newcommand{\BestLatPFive}{10,46}
\newcommand{\BestModel}{CatBoost}
\newcommand{\BestThreshold}{0,0746}
\newcommand{\CatAUC}{0,8614}
\newcommand{\CatLat}{10,46}
\newcommand{\CostOptimalCut}{0,0746}
\newcommand{\CutoffCapBadRate}{4,92\%}
\newcommand{\CutoffCostBadRate}{3,33\%}
\newcommand{\CutoffRateBadRate}{6,01\%}
\newcommand{\DRTest}{17,68\%}
\newcommand{\DRTrain}{14,35\%}
\newcommand{\DTAUC}{0,8206}
\newcommand{\ETAUC}{0,8471}
\newcommand{\GBAUC}{0,8575}
\newcommand{\GapBayes}{0,0495}
\newcommand{\GapBestLR}{0,0035}
\newcommand{\GapBestSVM}{0,0074}
\newcommand{\GapTopDeploy}{0,0006}
\newcommand{\GapTopSecond}{0,0006}
\newcommand{\GenSeconds}{39,5}
\newcommand{\KSPoints}{56,9}
\newcommand{\LGBAUC}{0,8590}
\newcommand{\LRAUC}{0,8579}
\newcommand{\LatRatioTopDeploy}{19,4}
\newcommand{\MLPAUC}{0,8560}
\newcommand{\MaxIV}{1,244}
\newcommand{\MaxMissingPct}{11,65\%}
\newcommand{\MaxOverfitGap}{0,0050}
\newcommand{\MaxPSI}{0,0598}
\newcommand{\MeanOverfitGap}{0,0027}
\newcommand{\NAccounts}{73.023}
\newcommand{\NCategorical}{3}
\newcommand{\NConstraints}{15}
\newcommand{\NCustomers}{40.000}
\newcommand{\NFeatures}{116}
\newcommand{\NFeaturesInitial}{154}
\newcommand{\NFeaturesSelected}{116}
\newcommand{\NMissingVars}{9}
\newcommand{\NModels}{12}
\newcommand{\NNumeric}{113}
\newcommand{\NOutlierVars}{122}
\newcommand{\NPSIAlert}{0}
\newcommand{\NPSIWarn}{0}
\newcommand{\NRules}{21}
\newcommand{\NTest}{9.000}
\newcommand{\NTiedModels}{3}
\newcommand{\NTotal}{60.000}
\newcommand{\NTrain}{42.000}
\newcommand{\NTransactions}{4.538.127}
\newcommand{\NViolations}{0}
\newcommand{\PDApproveMax}{0,0654}
\newcommand{\PDRejectMin}{0,1092}
\newcommand{\PropBadRate}{3,43\%}
\newcommand{\PropReviewRate}{32,74\%}
\newcommand{\RFAUC}{0,8527}
\newcommand{\RejectRate}{38,08\%}
\newcommand{\ReviewRate}{32,74\%}
\newcommand{\ReviewRateRise}{5,52}
\newcommand{\SVMAUC}{0,8540}
\newcommand{\SVMTrainSec}{16,1}
\newcommand{\SecondAUC}{0,8614}
\newcommand{\SecondModel}{CatBoost}
\newcommand{\ShapBehaviourPct}{34,0}
\newcommand{\ShapCICPct}{30,8}
\newcommand{\ShapDerivedPct}{24,0}
\newcommand{\ShapLoanPct}{4,2}
\newcommand{\ShapStaticPct}{7,0}
\newcommand{\StackAUC}{0,8620}
\newcommand{\StackLat}{202,39}
\newcommand{\SysMean}{40,33}
\newcommand{\SysPFifty}{34,93}
\newcommand{\SysPNinetyFive}{74,66}
\newcommand{\SysPNinetyNine}{98,37}
\newcommand{\ThinFileRate}{9,66\%}
\newcommand{\TieTolerance}{0,002}
\newcommand{\TopAUC}{0,8620}
\newcommand{\TopAUCLatPFive}{202,39}
\newcommand{\TopAUCModel}{Kết hợp xếp tầng (Stacking)}
\newcommand{\TopDecileBad}{40,7\%}
\newcommand{\TopThreeDecileBad}{76,8\%}
\newcommand{\TradApproveRate}{35,82\%}
\newcommand{\TradBadRate}{6,14\%}
\newcommand{\TradReviewRate}{27,22\%}
\newcommand{\TrainMinutes}{27,3}
\newcommand{\ValidityDetected}{100\%}
\newcommand{\VintageCoverage}{100\%}
\newcommand{\WOEAUC}{0,8570}
\newcommand{\XGBAUC}{0,8612}
\newcommand{\XGBLat}{18,53}


% --- Thông tin đề án -------------------------------------------------------
\newcommand{\thesisTitle}{ỨNG DỤNG MÔ HÌNH HỌC MÁY TRONG HỖ TRỢ KIỂM DUYỆT TÍN DỤNG}
\newcommand{\thesisAuthor}{Đồng Quốc Huy}
\newcommand{\thesisStudentId}{B25CHHT028}
\newcommand{\thesisMajor}{Hệ thống thông tin}
\newcommand{\thesisMajorCode}{8.48.01.04}
\newcommand{\thesisSupervisor}{TS. Đỗ Tiến Dũng}
\newcommand{\thesisInstitution}{HỌC VIỆN CÔNG NGHỆ BƯU CHÍNH VIỄN THÔNG}
\newcommand{\thesisYear}{2026}
\newcommand{\thesisCity}{HÀ NỘI}

\hypersetup{
  pdftitle={\thesisTitle},
  pdfauthor={\thesisAuthor},
  pdfsubject={Đề án tốt nghiệp Thạc sĩ Kỹ thuật - Hệ thống thông tin},
  pdfkeywords={kiểm duyệt tín dụng, học máy, chấm điểm tín dụng, rủi ro tín dụng,
               hồi quy logistic, rừng ngẫu nhiên, máy vectơ hỗ trợ,
               gradient boosting, mô hình hỗ trợ ra quyết định},
}

\begin{document}

% ===========================================================================
% PHẦN ĐẦU
% ===========================================================================
\pagenumbering{roman}

% ===========================================================================
% Bìa chính và bìa phụ
% ===========================================================================
\begin{titlepage}
\thispagestyle{empty}
\centering
\singlespacing

{\large\bfseries \thesisInstitution\par}
\vspace{2pt}
{\rule{0.55\textwidth}{0.8pt}\par}

\vspace{3.2cm}

{\Large\bfseries \thesisAuthor\par}

\vspace{2.6cm}

{\large\bfseries \thesisTitle\par}

\vspace{2.6cm}

{\large\bfseries ĐỀ ÁN TỐT NGHIỆP THẠC SĨ KỸ THUẬT\par}
\vspace{4pt}
{\normalsize (Theo định hướng ứng dụng)\par}

\vfill

{\large\bfseries \thesisCity\ --\ \thesisYear\par}
\end{titlepage}

% --------------------------------------------------------------------------
% Bìa phụ (có thông tin chuyên ngành và người hướng dẫn)
% --------------------------------------------------------------------------
\begin{titlepage}
\thispagestyle{empty}
\centering
\singlespacing

{\large\bfseries \thesisInstitution\par}
\vspace{2pt}
{\rule{0.55\textwidth}{0.8pt}\par}

\vspace{2.4cm}

{\Large\bfseries \thesisAuthor\par}

\vspace{1.9cm}

{\large\bfseries \thesisTitle\par}

\vspace{1.9cm}

{\normalsize
\begin{tabular}{@{}ll@{}}
\textbf{CHUYÊN NGÀNH} & : \thesisMajor \\[3pt]
\textbf{MÃ SỐ}        & : \thesisMajorCode \\
\end{tabular}\par}

\vspace{1.5cm}

{\large\bfseries ĐỀ ÁN TỐT NGHIỆP THẠC SĨ KỸ THUẬT\par}
\vspace{4pt}
{\normalsize (Theo định hướng ứng dụng)\par}

\vspace{1.8cm}

{\normalsize\bfseries NGƯỜI HƯỚNG DẪN KHOA HỌC: \thesisSupervisor\par}

\vfill

{\large\bfseries \thesisCity\ --\ \thesisYear\par}
\end{titlepage}

% ===========================================================================
% Lời cam đoan
% ===========================================================================
\chapter*{LỜI CAM ĐOAN}
\addcontentsline{toc}{chapter}{LỜI CAM ĐOAN}
\markboth{LỜI CAM ĐOAN}{LỜI CAM ĐOAN}

Tôi xin cam đoan đề án tốt nghiệp với đề tài \textit{``\thesisTitle''} là công
trình nghiên cứu của riêng tôi, được thực hiện dưới sự hướng dẫn khoa học của
\thesisSupervisor.

Các nội dung nghiên cứu, kết quả thực nghiệm trình bày trong đề án là trung thực
và chưa từng được công bố trong bất kỳ công trình nào khác. Toàn bộ mã nguồn của
mô hình nguyên mẫu, quy trình sinh dữ liệu, kịch bản thực nghiệm và các bảng
biểu, đồ thị trong đề án đều do tôi tự xây dựng và thực hiện.

Tôi xin nêu rõ về dữ liệu sử dụng trong đề án: do đặc thù bảo mật của dữ liệu tín
dụng cá nhân và các quy định về bảo vệ dữ liệu khách hàng, đề án \textbf{không sử
dụng dữ liệu thật của bất kỳ tổ chức tín dụng nào}. Toàn bộ dữ liệu phục vụ thực
nghiệm được sinh mô phỏng bằng chương trình do tôi xây dựng, theo một mô hình xác
suất được thiết kế và hiệu chỉnh để tái tạo các đặc điểm thống kê của dữ liệu tín
dụng bán lẻ. Cơ sở và giới hạn của cách tiếp cận này được trình bày minh bạch tại
mục~\ref{sec:data-design} và mục~\ref{subsec:han-che-dinh-huong}.

Các thông tin trích dẫn trong đề án đều được chỉ rõ nguồn gốc và liệt kê đầy đủ
trong danh mục tài liệu tham khảo. Tôi xin chịu hoàn toàn trách nhiệm về nội dung
đề án của mình.

\vspace{1.4cm}

\hfill\begin{minipage}{0.45\textwidth}
\centering
\textit{\thesisCity, tháng 12 năm \thesisYear}\\[4pt]
\textbf{Học viên thực hiện}\\[2.2cm]
\textbf{\thesisAuthor}
\end{minipage}

% ===========================================================================
% Lời cảm ơn
% ===========================================================================
\chapter*{LỜI CẢM ƠN}
\addcontentsline{toc}{chapter}{LỜI CẢM ƠN}
\markboth{LỜI CẢM ƠN}{LỜI CẢM ƠN}

Trước hết, tôi xin bày tỏ lòng biết ơn sâu sắc tới \thesisSupervisor\ -- người đã
trực tiếp hướng dẫn khoa học cho đề án này. Thầy đã định hướng cho tôi cách tiếp
cận vấn đề một cách hệ thống, từ việc xác định phạm vi nghiên cứu, lựa chọn
phương pháp, cho tới cách tổ chức thực nghiệm và trình bày kết quả một cách chặt
chẽ, trung thực.

Tôi xin trân trọng cảm ơn các thầy, cô giáo Khoa Công nghệ thông tin và Khoa Đào
tạo Sau đại học -- \thesisInstitution\ đã tận tình giảng dạy và tạo điều kiện
thuận lợi cho tôi trong suốt quá trình học tập chương trình cao học, cũng như
trong thời gian thực hiện đề án tốt nghiệp.

Tôi cũng xin gửi lời cảm ơn tới các đồng nghiệp trong lĩnh vực công nghệ ngân
hàng đã chia sẻ những kinh nghiệm thực tiễn quý báu về quy trình phê duyệt tín
dụng và vận hành mô hình, giúp đề án bám sát yêu cầu nghiệp vụ thực tế.

Cuối cùng, tôi xin cảm ơn gia đình và bạn bè đã luôn động viên, hỗ trợ tôi trong
suốt quá trình học tập và nghiên cứu.

Do thời gian và năng lực còn hạn chế, đề án chắc chắn không tránh khỏi những
thiếu sót. Tôi rất mong nhận được các ý kiến đóng góp của quý thầy, cô và các
bạn đọc để đề án được hoàn thiện hơn.

\vspace{0.9cm}

\hfill\begin{minipage}{0.45\textwidth}
\centering
\textit{\thesisCity, tháng 12 năm \thesisYear}\\[4pt]
\textbf{Học viên thực hiện}\\[2.2cm]
\textbf{\thesisAuthor}
\end{minipage}


\cleardoublepage
\phantomsection
\addcontentsline{toc}{chapter}{MỤC LỤC}
\renewcommand{\contentsname}{MỤC LỤC}
\tableofcontents

\cleardoublepage
% ===========================================================================
% Danh mục thuật ngữ và chữ viết tắt
% ===========================================================================
\phantomsection
\addcontentsline{toc}{chapter}{DANH MỤC THUẬT NGỮ VÀ CHỮ VIẾT TẮT}
\chapter*{DANH MỤC THUẬT NGỮ VÀ CHỮ VIẾT TẮT}
\markboth{DANH MỤC THUẬT NGỮ VÀ CHỮ VIẾT TẮT}{}

\begingroup
\small
\setlength{\extrarowheight}{1.5pt}
\begin{longtable}{@{}L{2.2cm}L{5.0cm}L{6.6cm}@{}}
\toprule
\textbf{Viết tắt} & \textbf{Thuật ngữ tiếng Anh} & \textbf{Giải nghĩa tiếng Việt} \\
\midrule
\endfirsthead
\multicolumn{3}{@{}l}{\small\textit{(tiếp theo trang trước)}}\\
\toprule
\textbf{Viết tắt} & \textbf{Thuật ngữ tiếng Anh} & \textbf{Giải nghĩa tiếng Việt} \\
\midrule
\endhead
\bottomrule
\multicolumn{3}{r@{}}{\small\textit{(xem tiếp trang sau)}}\\
\endfoot
\bottomrule
\endlastfoot

API   & Application Programming Interface & Giao diện lập trình ứng dụng \\
AUC   & Area Under the ROC Curve & Diện tích dưới đường cong ROC \\
CART  & Classification and Regression Tree & Cây phân loại và hồi quy \\
CIC   & Credit Information Center & Trung tâm Thông tin tín dụng Quốc gia Việt Nam \\
DPD   & Days Past Due & Số ngày quá hạn \\
DTI   & Debt-to-Income ratio & Tỷ lệ nghĩa vụ trả nợ trên thu nhập \\
EAD   & Exposure at Default & Dư nợ tại thời điểm vỡ nợ \\
ECE   & Expected Calibration Error & Sai số hiệu chuẩn kỳ vọng \\
EL    & Expected Loss & Tổn thất kỳ vọng \\
FN    & False Negative & Âm tính giả (bỏ sót hồ sơ xấu) \\
FP    & False Positive & Dương tính giả (từ chối nhầm hồ sơ tốt) \\
GBM   & Gradient Boosting Machine & Máy tăng cường độ dốc \\
IV    & Information Value & Giá trị thông tin của biến \\
JSON  & JavaScript Object Notation & Định dạng trao đổi dữ liệu JSON \\
KS    & Kolmogorov--Smirnov statistic & Thống kê Kolmogorov--Smirnov \\
LGD   & Loss Given Default & Tỷ lệ tổn thất khi vỡ nợ \\
LR    & Logistic Regression & Hồi quy logistic \\
LTV   & Loan-to-Value ratio & Tỷ lệ vay trên giá trị tài sản bảo đảm \\
MLP   & Multi-Layer Perceptron & Mạng nơ-ron truyền thẳng nhiều lớp \\
NHNN  & --- & Ngân hàng Nhà nước Việt Nam \\
NPL   & Non-Performing Loan & Nợ xấu (nợ nhóm 3--5) \\
PAVA  & Pool Adjacent Violators Algorithm & Thuật toán gộp vi phạm liền kề \\
PD    & Probability of Default & Xác suất vỡ nợ \\
PR-AUC & Area Under Precision--Recall Curve & Diện tích dưới đường Precision--Recall \\
PSI   & Population Stability Index & Chỉ số ổn định phân phối \\
PTI   & Payment-to-Income ratio & Tỷ lệ khoản trả góp trên thu nhập \\
RBF   & Radial Basis Function & Hàm cơ sở bán kính (hạt nhân Gauss) \\
REST  & Representational State Transfer & Kiến trúc dịch vụ web REST \\
RF    & Random Forest & Rừng ngẫu nhiên \\
ROC   & Receiver Operating Characteristic & Đường đặc trưng hoạt động của bộ thu nhận \\
SHAP  & SHapley Additive exPlanations & Phương pháp giải thích mô hình dựa trên giá trị Shapley \\
SQL   & Structured Query Language & Ngôn ngữ truy vấn có cấu trúc \\
SVM   & Support Vector Machine & Máy vectơ hỗ trợ \\
TCTD  & --- & Tổ chức tín dụng \\
WOE   & Weight of Evidence & Trọng số bằng chứng \\
XGBoost & eXtreme Gradient Boosting & Thư viện tăng cường độ dốc cực hạn \\
\end{longtable}
\endgroup

\vspace{0.5cm}
\noindent\textbf{Một số thuật ngữ được sử dụng thống nhất trong đề án:}

\begingroup
\small
\begin{longtable}{@{}L{4.4cm}L{9.6cm}@{}}
\toprule
\textbf{Thuật ngữ} & \textbf{Cách hiểu trong đề án} \\
\midrule
\endfirsthead
\toprule
\textbf{Thuật ngữ} & \textbf{Cách hiểu trong đề án} \\
\midrule
\endhead
\bottomrule
\endlastfoot

Kiểm duyệt tín dụng &
Quá trình đánh giá, sàng lọc hồ sơ vay để hình thành đề xuất có căn cứ. Phân biệt
với \textit{phê duyệt} -- hành vi ra quyết định của cấp có thẩm quyền. \\

Mô hình hỗ trợ kiểm duyệt &
Mô hình cung cấp đánh giá rủi ro định lượng kèm lý giải làm đầu vào cho quyết
định của con người, không thay thế thẩm quyền phê duyệt. \\

Tính đúng đắn theo thời điểm \newline (\eng{point-in-time correctness}) &
Nguyên tắc bảo đảm mọi đặc trưng của một hồ sơ chỉ được tính từ dữ liệu đã tồn
tại trước hoặc tại thời điểm nộp hồ sơ. \\

Rò rỉ nhãn \newline (\eng{label leakage}) &
Hiện tượng đặc trưng chứa thông tin phát sinh sau thời điểm nộp hồ sơ, khiến mô
hình cho kết quả cao giả tạo khi đánh giá và sụp đổ khi triển khai. \\

Cửa sổ quan sát \newline (\eng{performance window}) &
Khoảng thời gian kể từ khi cấp tín dụng dùng để xác định hồ sơ có phát sinh nợ
xấu hay không. Đề án sử dụng cửa sổ 12 tháng. \\

Kiểm định ngoài thời gian \newline (\eng{out-of-time validation}) &
Cách đánh giá mô hình trên dữ liệu thuộc giai đoạn sau giai đoạn huấn luyện, phản
ánh đúng điều kiện vận hành thực tế. \\

Trần Bayes &
Giới hạn lý thuyết về chất lượng dự báo mà không mô hình nào vượt qua được, tính
từ xác suất vỡ nợ thật của quá trình sinh dữ liệu. \\

Hiệu chuẩn xác suất \newline (\eng{calibration}) &
Quá trình điều chỉnh để giá trị xác suất mô hình đưa ra trùng khớp với tần suất
xảy ra thực tế -- điều kiện cần để dùng PD trong phân vùng quyết định. \\

Khách hàng thin-file &
Khách hàng có rất ít hoặc không có lịch sử tín dụng tại CIC, khiến việc đánh giá
phải dựa chủ yếu vào dữ liệu hành vi giao dịch. \\

Cắt ngọn \newline (\eng{winsorize}) &
Kỹ thuật giới hạn giá trị cực trị về mức phân vị biên thay vì loại bỏ quan sát,
nhằm giữ nguyên cỡ mẫu. \\

Trôi dạt dữ liệu \newline (\eng{data drift}) &
Hiện tượng phân phối dữ liệu hoặc quan hệ giữa đặc trưng và nhãn thay đổi theo
thời gian, làm mô hình suy giảm chất lượng. \\
\end{longtable}
\endgroup


\cleardoublepage
\phantomsection
\addcontentsline{toc}{chapter}{DANH MỤC BẢNG}
\renewcommand{\listtablename}{DANH MỤC BẢNG}
\listoftables

\cleardoublepage
\phantomsection
\addcontentsline{toc}{chapter}{DANH MỤC HÌNH VẼ}
\renewcommand{\listfigurename}{DANH MỤC HÌNH VẼ}
\listoffigures

% ===========================================================================
% PHẦN NỘI DUNG
% ===========================================================================
\cleardoublepage
\pagenumbering{arabic}
\setcounter{page}{1}

% ===========================================================================
% MỞ ĐẦU
% ===========================================================================
\chapter*{MỞ ĐẦU}
\addcontentsline{toc}{chapter}{MỞ ĐẦU}
\markboth{MỞ ĐẦU}{MỞ ĐẦU}

\section*{1. Lý do chọn đề tài}
\addcontentsline{toc}{section}{1. Lý do chọn đề tài}

Trong những năm gần đây, chuyển đổi số trong lĩnh vực ngân hàng đã thúc đẩy mạnh
mẽ quá trình số hóa các sản phẩm và dịch vụ tín dụng. Cùng với sự phát triển của
ngân hàng số, nhu cầu cấp tín dụng nhanh chóng, thuận tiện và có khả năng xử lý
số lượng lớn hồ sơ ngày càng gia tăng. Đối với các tổ chức tín dụng, kiểm duyệt
tín dụng là một trong những nghiệp vụ quan trọng, có ảnh hưởng trực tiếp đến chất
lượng tài sản, khả năng kiểm soát rủi ro và hiệu quả kinh doanh. Do đó, việc ứng
dụng các mô hình học máy hỗ trợ cán bộ tín dụng trong quá trình đánh giá và ra
quyết định đang trở thành một yêu cầu thực tiễn.

Trong phương thức kiểm duyệt tín dụng truyền thống, thông tin khách hàng thường
được tổng hợp từ nhiều nguồn khác nhau như hồ sơ đăng ký vay, thông tin khách
hàng, lịch sử giao dịch, dữ liệu tín dụng và các thông tin liên quan khác. Quá
trình tổng hợp, kiểm tra và đánh giá dữ liệu có thể mất nhiều thời gian, đặc biệt
khi khối lượng hồ sơ lớn. Bên cạnh đó, việc đánh giá dựa nhiều vào các bộ quy tắc
nghiệp vụ cố định có thể hạn chế khả năng phát hiện các mối quan hệ phức tạp giữa
các đặc trưng của khách hàng và khả năng phát sinh rủi ro tín dụng.

Hạn chế nêu trên không nằm ở năng lực chuyên môn của cán bộ tín dụng mà nằm ở
\textit{bản chất của phương pháp đánh giá}. Một bộ quy tắc, dù được xây dựng công
phu đến đâu, cũng chỉ biểu diễn được quan hệ giữa từng chỉ tiêu riêng lẻ với
ngưỡng cho phép. Trong khi đó, rủi ro tín dụng thường phát sinh từ \textit{tổ
hợp} nhiều yếu tố: một khách hàng có thu nhập cao nhưng dòng tiền không ổn định,
thâm niên công việc ngắn và đang có quan hệ tín dụng với nhiều tổ chức cùng lúc
có thể rủi ro hơn một khách hàng thu nhập thấp hơn nhưng ổn định về mọi mặt. Bộ
quy tắc xét từng điều kiện độc lập khó nắm bắt được kiểu tương tác này.

Sự phát triển của các phương pháp học máy mở ra khả năng khai thác dữ liệu lịch
sử để xây dựng các mô hình dự đoán và hỗ trợ đánh giá rủi ro. Thay vì chỉ dựa
trên các điều kiện kiểm tra cố định, mô hình học máy có thể học các mối quan hệ
giữa đặc điểm khách hàng, hành vi tín dụng và kết quả tín dụng trong quá khứ để
đưa ra dự đoán đối với hồ sơ mới. Khi được tích hợp vào quy trình nghiệp vụ, kết
quả dự đoán có thể được sử dụng như một nguồn thông tin hỗ trợ cán bộ trong quá
trình kiểm duyệt tín dụng.

Tuy nhiên, việc ứng dụng học máy vào kiểm duyệt tín dụng không chỉ là bài toán
xây dựng mô hình dự đoán mà còn liên quan đến quá trình thu thập, chuẩn hóa, biến
đổi dữ liệu, xây dựng đặc trưng, tích hợp mô hình và cung cấp kết quả dự đoán một
cách ổn định. Thực tế cho thấy nhiều mô hình đạt chỉ tiêu tốt trong phòng thí
nghiệm nhưng không áp dụng được, vì thiếu các yếu tố mà môi trường vận hành đòi
hỏi: xác suất dự báo phải được hiệu chuẩn đúng để dùng làm căn cứ định giá, quyết
định phải giải trình được trước cán bộ tín dụng và cơ quan quản lý, và mô hình
phải xử lý được cả những hồ sơ thiếu dữ liệu. Vì vậy, cần có một giải pháp mang
tính hệ thống, trong đó mô hình học máy được kết hợp với quy trình xử lý dữ liệu
và các luật nghiệp vụ nhằm hỗ trợ quá trình kiểm duyệt tín dụng.

Xuất phát từ những yêu cầu trên, đề án lựa chọn đề tài
\textit{``\thesisTitle''} nhằm nghiên cứu, thiết kế và xây dựng thử nghiệm một mô
hình có khả năng tiếp nhận thông tin hồ sơ tín dụng, xử lý dữ liệu, tạo các đặc
trưng đầu vào, sử dụng mô hình học máy để đánh giá khả năng rủi ro và cung cấp
kết quả hỗ trợ kiểm duyệt. Đề tài có ý nghĩa cả về mặt khoa học và thực tiễn,
đồng thời có tính khả thi nhờ sự phát triển của các công nghệ xử lý dữ liệu, học
máy và các nền tảng triển khai hiện nay.

\section*{2. Tổng quan về vấn đề nghiên cứu}
\addcontentsline{toc}{section}{2. Tổng quan về vấn đề nghiên cứu}

Kiểm duyệt tín dụng là một bài toán đã được nghiên cứu rộng rãi trong lĩnh vực
tài chính -- ngân hàng. Các phương pháp truyền thống thường sử dụng chấm điểm tín
dụng và các bộ luật nghiệp vụ dựa trên những tiêu chí như thu nhập, lịch
sử tín dụng, dư nợ, khả năng trả nợ và đặc điểm của khoản vay. Những phương pháp
này có ưu điểm là dễ giải thích và phù hợp với các quy định nghiệp vụ, tuy nhiên
khả năng biểu diễn các mối quan hệ phức tạp trong dữ liệu còn hạn chế.

Cùng với sự gia tăng về quy mô và đa dạng của dữ liệu, nhiều nghiên cứu đã tập
trung áp dụng các thuật toán học máy cho bài toán đánh giá tín dụng. Một số
phương pháp phổ biến gồm hồi quy Logistic, Decision Tree, Random Forest, Support
Vector Machine, Gradient Boosting và các mô hình học sâu. Các nghiên cứu cho thấy
mô hình học máy có khả năng khai thác nhiều đặc trưng đồng thời và cải thiện khả
năng phân loại khách hàng có nguy cơ tín dụng cao hoặc thấp so với một số phương
pháp truyền thống.

Công trình đối sánh quy mô lớn của Lessmann và cộng sự~\cite{lessmann2015} trên
nhiều bộ dữ liệu tín dụng là một mốc tham chiếu quan trọng: các phương pháp học
tập hợp cho kết quả tốt hơn mô hình đơn lẻ, song mức chênh lệch so với hồi quy
logistic là \textit{không lớn} -- một kết luận thường bị bỏ qua khi lựa chọn mô
hình trong thực tiễn. Các tổng quan hệ thống gần đây~\cite{shi2022} tiếp tục
khẳng định xu hướng này, đồng thời chỉ ra rằng yếu tố quyết định chất lượng mô
hình thường nằm ở \textit{chất lượng và độ phong phú của đặc trưng} hơn là ở việc
lựa chọn thuật toán.

Bên cạnh vấn đề lựa chọn thuật toán, các nghiên cứu cũng tập trung vào các bước
xử lý dữ liệu như làm sạch dữ liệu, xử lý giá trị thiếu, lựa chọn đặc trưng, cân
bằng dữ liệu và đánh giá mô hình. Các chỉ số như Accuracy, Precision, Recall,
F1-score, ROC-AUC và Gini thường được sử dụng để đánh giá hiệu quả của mô hình
phân loại rủi ro tín dụng.

Tuy nhiên, một hạn chế thường gặp là nhiều nghiên cứu tập trung chủ yếu vào hiệu
quả dự đoán của mô hình, trong khi chưa giải quyết đầy đủ bài toán tích hợp mô
hình vào một quy trình hỗ trợ kiểm duyệt tín dụng hoàn chỉnh. Trong môi trường
thực tế, dữ liệu tín dụng có thể đến từ nhiều nguồn, có cấu trúc và chất lượng
khác nhau, đồng thời yêu cầu xử lý dữ liệu, kiểm tra điều kiện nghiệp vụ và cung
cấp kết quả trong thời gian phù hợp với quy trình kinh doanh.

Do đó, đề án tập trung nghiên cứu theo hướng kết hợp giữa xử lý dữ liệu, xây dựng
đặc trưng, mô hình học máy và luật nghiệp vụ để hình thành một mô hình hỗ trợ
kiểm duyệt tín dụng. Đề án không hướng tới việc thay thế hoàn toàn quyết định của
cán bộ hoặc tổ chức tín dụng mà tập trung xây dựng một công cụ hỗ trợ, cung cấp
kết quả đánh giá rủi ro và thông tin tham khảo nhằm nâng cao hiệu quả và tính
nhất quán của quá trình kiểm duyệt.

\section*{3. Mục tiêu nghiên cứu}
\addcontentsline{toc}{section}{3. Mục tiêu nghiên cứu}

\subsection*{3.1. Mục tiêu chung}

Đề án nhằm nghiên cứu, thiết kế và xây dựng thử nghiệm mô hình học máy hỗ trợ
kiểm duyệt tín dụng, trong đó khai thác dữ liệu khách hàng và thông tin tín dụng
để xây dựng các đặc trưng đầu vào, huấn luyện mô hình học máy và đưa ra kết quả
dự đoán mức độ rủi ro. Mô hình được xây dựng theo định hướng hỗ trợ quá trình
kiểm duyệt tín dụng, giúp cung cấp thêm cơ sở định lượng cho việc đánh giá hồ sơ
thay vì chỉ dựa trên các tiêu chí và quy tắc nghiệp vụ cố định.

\subsection*{3.2. Mục tiêu về mặt lý luận}

Về mặt lý luận, đề án tập trung hệ thống hóa và làm rõ cơ sở lý thuyết về kiểm
duyệt tín dụng, đánh giá rủi ro tín dụng và khả năng ứng dụng các phương pháp học
máy trong bài toán hỗ trợ ra quyết định tín dụng. Đề án nghiên cứu quy trình từ
thu thập và tiền xử lý dữ liệu, xây dựng đặc trưng, lựa chọn mô hình, huấn luyện,
đánh giá đến tích hợp mô hình. Trên cơ sở đó, đề án phân tích và
đánh giá mức độ phù hợp của các mô hình học máy đối với bài toán phân loại rủi ro
tín dụng thông qua các tiêu chí như Accuracy, Precision, Recall, F1-score và
ROC-AUC. Kết quả nghiên cứu góp phần làm rõ mối quan hệ giữa chất lượng dữ liệu,
đặc trưng đầu vào, lựa chọn mô hình và hiệu quả dự đoán trong bài toán tín dụng,
đồng thời bổ sung cơ sở tham khảo cho việc nghiên cứu ứng dụng học máy vào việc
hỗ trợ ra quyết định trong lĩnh vực tài chính -- ngân hàng.

\subsection*{3.3. Mục tiêu về mặt thực tiễn}

Về mặt thực tiễn, đề án hướng tới xây dựng một mô hình học máy thử nghiệm có khả
năng tiếp nhận thông tin hồ sơ tín dụng, thực hiện xử lý dữ liệu, tạo đặc trưng
để đánh giá rủi ro và trả về kết quả hỗ trợ kiểm duyệt. Việc ứng dụng mô hình học
máy giúp khai thác hiệu quả hơn dữ liệu lịch sử, hỗ trợ nhận diện các hồ sơ có
mức độ rủi ro khác nhau và cung cấp kết quả đánh giá một cách nhất quán. Qua đó,
có thể góp phần rút ngắn thời gian đánh giá hồ sơ, giảm khối lượng công việc thủ
công, nâng cao khả năng phát hiện rủi ro và hỗ trợ cán bộ tín dụng trong quá
trình ra quyết định. Kết quả của đề án cũng tạo cơ sở thực nghiệm để đánh giá khả
năng áp dụng mô hình học máy vào quy trình kiểm duyệt tín dụng và làm tiền đề cho
việc phát triển có quy mô lớn hơn trong môi trường thực tế.

\subsection*{3.4. Các mục tiêu cụ thể}

Từ hai định hướng trên, đề án đặt ra các mục tiêu cụ thể sau:

\begin{enumerate}[label=(\arabic*)]
  \item Hệ thống hóa cơ sở lý luận về tín dụng, kiểm duyệt tín dụng, rủi ro tín
  dụng và các phương pháp học máy áp dụng cho bài toán đánh giá tín dụng.

  \item Phân tích bài toán kiểm duyệt tín dụng: xác định các mức độ rủi ro cần dự
  báo, đặc trưng đầu vào, kết quả đầu ra và các yêu cầu kỹ thuật, pháp lý đối với
  mô hình.

  \item Xây dựng quy trình thu thập, tổng hợp, kiểm tra chất lượng và tiền xử lý
  dữ liệu hồ sơ tín dụng.

  \item Trích xuất và lựa chọn tập đặc trưng cho mô hình từ bốn nhóm nguồn dữ
  liệu: thông tin cá nhân và tài chính, tham số khoản vay, lịch sử tín dụng và
  quan hệ CIC, hành vi giao dịch tài khoản.

  \item Xây dựng nhãn rủi ro theo tiêu chuẩn nợ xấu và cửa sổ quan sát phù hợp;
  xác định điểm cắt phê duyệt theo các tiêu chí nghiệp vụ khác nhau.

  \item Thiết kế, huấn luyện, hiệu chuẩn và đánh giá so sánh các mô hình học máy
  theo bộ chỉ tiêu đầy đủ, bao gồm cả chỉ tiêu đặc thù của ngành tín dụng (KS,
  Gini, IV, PSI) và thời gian chấm điểm.

  \item Thiết kế và hiện thực bộ máy hỗ trợ ra quyết định kết hợp luật nghiệp vụ
  với mô hình học máy, có cơ chế giải thích kết quả cho từng hồ sơ.

  \item Mô phỏng luồng xử lý hồ sơ, thử nghiệm tích hợp qua dịch vụ REST và đối
  chiếu kết quả với quy trình thẩm định truyền thống chỉ dựa trên luật nghiệp vụ.
\end{enumerate}

\section*{4. Đối tượng và phạm vi nghiên cứu}
\addcontentsline{toc}{section}{4. Đối tượng và phạm vi nghiên cứu}

\subsection*{4.1. Đối tượng nghiên cứu}

Đối tượng nghiên cứu của đề án là mô hình học máy hỗ trợ kiểm duyệt tín dụng, tập
trung vào dữ liệu và quy trình đánh giá hồ sơ tín dụng. Cụ thể,
đề án nghiên cứu các dữ liệu liên quan đến khách hàng và khoản vay như thông tin
cá nhân, thông tin khoản vay, thu nhập, lịch sử tín dụng, tình trạng dư nợ và các
đặc trưng có liên quan đến khả năng phát sinh rủi ro tín dụng. Trên cơ sở các dữ
liệu này, đề án nghiên cứu phương pháp xử lý dữ liệu, xây dựng đặc trưng, huấn
luyện mô hình học máy và sử dụng kết quả dự đoán để hỗ trợ quá trình kiểm duyệt
tín dụng. Đối tượng nghiên cứu bao gồm mô hình học máy và kết quả đánh giá rủi ro
thành một quy trình hỗ trợ kiểm duyệt thống nhất.

\subsection*{4.2. Phạm vi nghiên cứu}

\textbf{Về nội dung:} Đề án giới hạn ở việc xây dựng và thử nghiệm mô hình hỗ trợ
đánh giá rủi ro tín dụng, bao gồm các bước chính: tiếp nhận dữ liệu hồ sơ,
tiền xử lý dữ liệu, xây dựng đặc trưng, áp dụng mô hình học máy, đánh giá kết quả
dự đoán và cung cấp kết quả hỗ trợ kiểm duyệt. Đề án không đi sâu nghiên cứu toàn
bộ quy trình nghiệp vụ tín dụng của ngân hàng, không xây dựng hệ thống quản lý
tín dụng hoàn chỉnh và không thay thế các quy định, chính sách hoặc quyết định
phê duyệt chính thức của tổ chức tín dụng. Các nội dung sau \textit{không} thuộc
phạm vi: định giá và thẩm định tài sản bảo đảm, quy trình giải ngân và thu hồi
nợ, mô hình ước lượng tổn thất khi vỡ nợ (LGD) và dư nợ tại thời điểm vỡ nợ
(EAD).

\textbf{Về không gian:} Nghiên cứu được thực hiện trong phạm vi bài toán kiểm
duyệt tín dụng tại tổ chức tín dụng/ngân hàng, tập trung vào việc xây dựng mô
hình và thử nghiệm trong môi trường nghiên cứu. Đề án không giới hạn vào
một địa bàn cụ thể và không thực hiện triển khai chính thức trên toàn bộ hệ thống
nghiệp vụ của một ngân hàng.

\textbf{Về dữ liệu:} Do đặc thù bảo mật của dữ liệu tín dụng cá nhân, đề án sử
dụng \textbf{dữ liệu mô phỏng} được sinh bằng chương trình tự xây dựng theo một
mô hình xác suất có cấu trúc. Thiết kế, cơ sở hiệu chỉnh và giới hạn của cách
tiếp cận này được trình bày minh bạch tại mục~\ref{sec:data-design}. Trong trường
hợp áp dụng trên dữ liệu thực tế, dữ liệu cần được ẩn danh và giới hạn quyền truy
cập nhằm bảo đảm yêu cầu về bảo mật thông tin.

\textbf{Về thời gian:} Đề án được thực hiện trong thời gian triển khai đề án tốt
nghiệp. Bộ dữ liệu sử dụng để huấn luyện và đánh giá mô hình được giới hạn trong
khoảng thời gian phù hợp với dữ liệu nghiên cứu và được phân chia thành các giai
đoạn phục vụ huấn luyện, kiểm thử và đánh giá mô hình. Các kết quả nghiên cứu
phản ánh phạm vi dữ liệu và điều kiện thực nghiệm trong thời gian thực hiện đề
án, chưa nhằm đánh giá khả năng vận hành của mô hình trong thời gian dài hoặc
trên quy mô dữ liệu sản xuất thực tế.

\textbf{Về kết quả:} Đề án tập trung đánh giá hai khía cạnh chính là hiệu quả của
mô hình học máy trong dự đoán rủi ro tín dụng và khả năng hỗ trợ quá trình kiểm
duyệt. Kết quả đầu ra của mô hình là thông tin đánh giá hoặc phân loại mức độ rủi
ro nhằm hỗ trợ cán bộ tín dụng và quy trình nghiệp vụ; quyết định
tín dụng cuối cùng vẫn thuộc về người hoặc bộ phận có thẩm quyền theo quy định
của tổ chức tín dụng.

\section*{5. Phương pháp nghiên cứu}
\addcontentsline{toc}{section}{5. Phương pháp nghiên cứu}

Đề án sử dụng kết hợp phương pháp nghiên cứu lý thuyết và phương pháp nghiên cứu
thực nghiệm nhằm nghiên cứu, xây dựng và đánh giá mô hình học máy trong hỗ trợ
kiểm duyệt tín dụng.

\textbf{Phương pháp nghiên cứu lý thuyết.} Đề án tiến hành thu thập, tổng hợp và
phân tích các tài liệu, giáo trình, công trình nghiên cứu và các tài liệu chuyên
ngành liên quan đến tín dụng, kiểm duyệt tín dụng, đánh giá rủi ro tín dụng, chấm
điểm tín dụng và ứng dụng học máy trong lĩnh vực tài chính -- ngân hàng. Trên cơ
sở nghiên cứu lý thuyết, đề án làm rõ các khái niệm, đặc điểm của bài toán đánh
giá tín dụng, các yếu tố ảnh hưởng đến rủi ro tín dụng và nguyên lý hoạt động của
các mô hình học máy phù hợp. Đồng thời, nghiên cứu các phương pháp tiền xử lý dữ
liệu, lựa chọn và xây dựng đặc trưng, huấn luyện và đánh giá mô hình. Kết quả
nghiên cứu lý thuyết được sử dụng làm cơ sở để lựa chọn mô hình học máy, thiết kế
quy trình xử lý dữ liệu và xây dựng mô hình học máy hỗ trợ kiểm duyệt tín dụng.

\textbf{Phương pháp nghiên cứu thực nghiệm.} Đề án tiến hành xây dựng môi trường
thực nghiệm để kiểm chứng các nội dung nghiên cứu lý thuyết trên bộ dữ liệu tín
dụng phục vụ nghiên cứu. Dữ liệu được phân tích và xử lý thông qua các bước làm
sạch, xử lý dữ liệu thiếu, mã hóa dữ liệu, xây dựng và lựa chọn đặc trưng trước
khi đưa vào mô hình. Một số mô hình học máy phù hợp với bài toán phân loại rủi ro
tín dụng được xây dựng, huấn luyện và kiểm thử trên dữ liệu thực nghiệm. Cụ thể,
phương pháp thực nghiệm được triển khai qua bốn nhóm công việc:

\begin{itemize}
  \item \textbf{Phân tích dữ liệu:} thống kê mô tả tập dữ liệu, kiểm tra tính hợp
  lệ theo ràng buộc nghiệp vụ, phân tích giá trị khuyết thiếu và giá trị ngoại
  lệ; đánh giá năng lực dự báo của từng đặc trưng thông qua chỉ số Information
  Value.

  \item \textbf{Thực nghiệm học máy:} huấn luyện, hiệu chuẩn và so sánh các mô
  hình theo phương pháp kiểm định ngoài thời gian (\eng{out-of-time validation});
  tinh chỉnh siêu tham số bằng kiểm định chéo theo chuỗi thời gian; giải thích mô
  hình bằng phương pháp SHAP và phép hoán vị.

  \item \textbf{Mô phỏng quy trình kiểm duyệt:} xây dựng bộ máy hỗ trợ ra quyết
  định, mô phỏng luồng xử lý hồ sơ, xác định ngưỡng phê duyệt theo ba tiêu chí
  nghiệp vụ và thử nghiệm tích hợp qua dịch vụ REST.

  \item \textbf{Đánh giá đối chiếu:} so sánh kết quả của mô hình đề xuất với quy
  trình thẩm định truyền thống chỉ dựa trên luật nghiệp vụ, và đối chiếu với các
  công trình nghiên cứu liên quan.
\end{itemize}

\section*{6. Kết cấu của đề án}
\addcontentsline{toc}{section}{6. Kết cấu của đề án}

Ngoài phần Mở đầu, Kết luận, Danh mục tài liệu tham khảo và Phụ lục, nội dung đề
án được trình bày trong ba chương:

\begin{itemize}
  \item \textbf{Chương 1 -- Cơ sở lý luận:} trình bày tổng quan về tín dụng và
  kiểm duyệt tín dụng, tổng quan về học máy trong đánh giá tín dụng, dữ liệu và
  đặc trưng trong bài toán tín dụng, các mô hình học máy áp dụng trong đánh giá
  tín dụng cùng cơ sở so sánh và lựa chọn mô hình phù hợp.

  \item \textbf{Chương 2 -- Xây dựng mô hình học máy trong hỗ trợ kiểm duyệt tín
  dụng:} phân tích bài toán kiểm duyệt tín dụng; trình bày quy trình thu thập
  và tổng hợp dữ liệu, tiền xử lý và làm sạch dữ liệu, trích xuất và lựa chọn đặc trưng, xây
  dựng nhãn và xác định quy tắc phê duyệt; thiết kế mô hình học máy và quy trình
  huấn luyện, tối ưu mô hình.

  \item \textbf{Chương 3 -- Triển khai thực nghiệm và đánh giá kết quả:} trình
  bày kịch bản và môi trường thực nghiệm, kết quả huấn luyện các mô hình, đánh
  giá hiệu năng và so sánh giữa các thuật toán, mô phỏng và giải thích mô hình;
  cuối cùng là phần đánh giá chung, đối chiếu với quy trình thẩm định
  truyền thống và với các công trình nghiên cứu liên quan.
\end{itemize}

% ===========================================================================
% CHƯƠNG 1: CƠ SỞ LÝ LUẬN
% ===========================================================================
\chapter{CƠ SỞ LÝ LUẬN}
% Tiêu đề chạy rút gọn: tiêu đề chương đầy đủ dài hơn một dòng đầu trang,
% khiến fancyhdr phải nới \headheight và làm lệch khung chữ của cả chương.
\markboth{Chương 1. Cơ sở lý luận}{}
\label{chap:co-so-ly-luan}

\noindent\textit{Chương 1 hệ thống hóa cơ sở lý thuyết của đề án. Nội dung chương
lần lượt trình bày tổng quan về tín dụng và hoạt động kiểm duyệt tín dụng; tổng
quan về học máy và vai trò của học máy trong đánh giá tín dụng; đặc điểm dữ liệu
và phương pháp xây dựng đặc trưng cho bài toán tín dụng; và cuối cùng là các mô
hình học máy được nghiên cứu trong đề án cùng cơ sở so sánh, lựa chọn mô hình phù
hợp. Đây là nền tảng để Chương 2 xây dựng mô hình và Chương 3 tiến hành thực
nghiệm.}

\section{Tổng quan về tín dụng và kiểm duyệt tín dụng}
\label{sec:tong-quan-tin-dung}

\subsection{Khái niệm tín dụng}
\label{subsec:khai-niem-tin-dung}

Tín dụng là quan hệ chuyển nhượng quyền sử dụng một lượng giá trị -- thường dưới
hình thái tiền tệ -- từ bên cho vay sang bên đi vay trong một khoảng thời gian
xác định, với cam kết hoàn trả cả gốc và lãi khi đến hạn. Ba đặc trưng cơ bản của
quan hệ tín dụng gồm: \textit{tính hoàn trả} (khoản vốn phải được trả lại, khác
với cấp phát hay tài trợ không hoàn lại), \textit{tính thời hạn} (quyền sử dụng
vốn chỉ được chuyển giao trong một thời gian nhất định) và \textit{tính sinh lời}
(bên cho vay nhận được khoản lãi như phần bù cho việc từ bỏ quyền sử dụng vốn và
cho rủi ro phải gánh chịu).

Đặc trưng thứ ba dẫn tới một hệ quả quan trọng về mặt kinh tế: lãi suất không chỉ
phản ánh giá của vốn theo thời gian mà còn phải bao hàm \textit{phần bù rủi ro}.
Một tổ chức tín dụng cho vay với cùng một mức lãi suất đối với mọi khách hàng sẽ
gặp hiện tượng lựa chọn ngược: khách hàng có rủi ro thấp thấy mức lãi suất đó
không hấp dẫn và tìm nguồn vốn khác, còn khách hàng có rủi ro cao sẵn sàng chấp
nhận. Kết quả là danh mục tín dụng dần tập trung vào nhóm rủi ro cao trong khi
lãi suất không được điều chỉnh tương xứng. Đây chính là lý do nền tảng khiến việc
\textit{phân biệt được mức độ rủi ro của từng khách hàng} trở thành năng lực cốt
lõi của hoạt động tín dụng, chứ không đơn thuần là một yêu cầu về tuân thủ.

Trong phạm vi đề án, đối tượng nghiên cứu là tín dụng bán lẻ đối với khách hàng
cá nhân: các khoản vay tiêu dùng, vay mua ô tô, vay thế chấp bất động sản và thấu
chi. Nhóm sản phẩm này có hai đặc điểm khiến chúng đặc biệt phù hợp với phương
pháp mô hình hóa thống kê. Thứ nhất, số lượng hồ sơ lớn nhưng giá trị mỗi khoản
vay tương đối nhỏ, nên chi phí thẩm định thủ công cho từng hồ sơ khó có thể tương
xứng với quy mô khoản vay. Thứ hai, các khoản vay trong cùng một sản phẩm có mức
độ đồng nhất cao, cho phép rút ra quy luật thống kê từ dữ liệu lịch sử -- điều
khó thực hiện với tín dụng doanh nghiệp lớn, nơi mỗi khoản vay gần như là một
trường hợp riêng biệt.

\subsection{Quy trình cấp và kiểm duyệt tín dụng}
\label{subsec:quy-trinh-kiem-duyet}

Quy trình cấp tín dụng tại các tổ chức tín dụng thường được tổ chức thành năm
bước chính:

\begin{enumerate}[label=(\arabic*)]
  \item \textbf{Tiếp nhận hồ sơ.} Khách hàng nộp đề nghị vay vốn kèm các tài liệu
  chứng minh nhân thân, thu nhập, mục đích sử dụng vốn và tài sản bảo đảm (nếu
  có). Hồ sơ có thể được tiếp nhận tại quầy, qua kênh trực tuyến hoặc qua đối tác
  phân phối.

  \item \textbf{Thu thập và kiểm tra thông tin.} Cán bộ tín dụng đối chiếu thông
  tin khách hàng khai báo với dữ liệu nội bộ của ngân hàng (lịch sử quan hệ, số
  dư tài khoản, lịch sử giao dịch) và với dữ liệu bên ngoài, trong đó quan trọng
  nhất là bản tin tín dụng từ Trung tâm Thông tin tín dụng Quốc gia Việt Nam
  (CIC)~\cite{tt15}.

  \item \textbf{Thẩm định và kiểm duyệt.} Đây là bước trọng tâm của quy trình:
  đánh giá năng lực tài chính, khả năng trả nợ, mức độ rủi ro của khách hàng và
  sự phù hợp của khoản vay với chính sách tín dụng hiện hành. Kết quả của bước
  này là đề xuất phê duyệt, từ chối hoặc yêu cầu bổ sung thông tin.

  \item \textbf{Phê duyệt.} Cấp có thẩm quyền -- tùy theo hạn mức được phân cấp
  -- ra quyết định cuối cùng trên cơ sở hồ sơ và đề xuất thẩm định.

  \item \textbf{Giải ngân và giám sát sau vay.} Sau khi hoàn tất thủ tục pháp lý,
  khoản vay được giải ngân và chuyển sang giai đoạn theo dõi việc trả nợ, phân
  loại nợ và trích lập dự phòng rủi ro theo quy định~\cite{tt11}.
\end{enumerate}

Phạm vi của đề án tập trung vào bước thứ ba -- \textbf{kiểm duyệt tín dụng}. Cần
phân biệt rõ hai khái niệm dễ nhầm lẫn: \textit{kiểm duyệt} là quá trình đánh giá
và sàng lọc hồ sơ để hình thành đề xuất có căn cứ, còn \textit{phê duyệt} là hành
vi ra quyết định của cấp có thẩm quyền. Hệ thống được xây dựng trong đề án hỗ trợ
bước kiểm duyệt: nó cung cấp đánh giá rủi ro định lượng kèm lý giải, làm đầu vào
cho quyết định của con người, chứ không thay thế thẩm quyền phê duyệt.

Sự phân biệt này không chỉ mang tính thuật ngữ mà có ý nghĩa thiết kế trực tiếp.
Một mô hình \textit{thay thế} con người buộc phải đưa ra quyết định nhị phân cho
mọi hồ sơ, kể cả những hồ sơ mà bản thân nó không đủ chắc chắn. Một mô hình
\textit{hỗ trợ} có thể -- và nên -- thừa nhận vùng không chắc chắn của mình và
chuyển những hồ sơ đó cho người có thẩm quyền xem xét. Toàn bộ thiết kế ba vùng
quyết định được trình bày ở Chương~\ref{chap:xay-dung-he-thong} bắt nguồn từ nhận
định này.

\subsection{Các yếu tố ảnh hưởng đến quyết định tín dụng}
\label{subsec:yeu-to-anh-huong}

Thực tiễn thẩm định tín dụng từ lâu đã tổng kết các yếu tố cần xem xét thành
khung phân tích \textbf{5C}~\cite{thomas2017}:

\begin{itemize}
  \item \textbf{Character} (Uy tín): thái độ và lịch sử thực hiện nghĩa vụ trả nợ
  của khách hàng. Trong dữ liệu, yếu tố này được phản ánh qua lịch sử tín dụng
  tại CIC: nhóm nợ, số ngày quá hạn, tiền sử nợ đã xử lý rủi ro.

  \item \textbf{Capacity} (Năng lực trả nợ): khả năng tạo ra dòng tiền để trả nợ,
  đo bằng thu nhập, tính ổn định của thu nhập và các chỉ tiêu như tỷ lệ nợ trên
  thu nhập ($\DTI$) hay tỷ lệ khoản trả góp trên thu nhập.

  \item \textbf{Capital} (Vốn tự có): phần vốn khách hàng tự bỏ ra trong phương
  án vay, thể hiện mức độ ``cùng chịu rủi ro''. Với khoản vay có tài sản bảo đảm,
  yếu tố này gắn với tỷ lệ vay trên giá trị tài sản ($\LTV$).

  \item \textbf{Collateral} (Tài sản bảo đảm): nguồn thu hồi nợ thứ cấp khi nguồn
  trả nợ chính không thực hiện được.

  \item \textbf{Conditions} (Điều kiện): bối cảnh kinh tế vĩ mô, đặc điểm ngành
  nghề và mục đích sử dụng vốn.
\end{itemize}

Khung 5C có giá trị định hướng nhưng bộc lộ một hạn chế khi chuyển thành quy tắc
vận hành: nó liệt kê các yếu tố cần xem xét mà không quy định cách \textit{kết
hợp} chúng thành một đánh giá tổng thể. Trong thực tế, mức độ quan trọng tương
đối của từng yếu tố thay đổi theo phân khúc khách hàng và theo thời gian; hơn nữa
các yếu tố tương tác với nhau theo cách phi tuyến. Một khách hàng có $\DTI$ ở mức
0,45 có thể chấp nhận được nếu thu nhập ổn định và lịch sử tín dụng tốt, nhưng
lại rủi ro cao nếu thâm niên công việc ngắn và đang có quan hệ tín dụng với nhiều
tổ chức. Chính khoảng trống này -- xác định trọng số và dạng tương tác giữa các
yếu tố từ dữ liệu thay vì từ phán đoán chủ quan -- là chỗ mà phương pháp học máy
đóng góp giá trị.

\subsection{Rủi ro tín dụng và đánh giá rủi ro tín dụng}
\label{subsec:rui-ro-tin-dung}

\textbf{Rủi ro tín dụng} là khả năng bên đi vay không thực hiện được nghĩa vụ trả
nợ theo cam kết, dẫn tới tổn thất cho bên cho vay. Theo khung Basel
II~\cite{bcbs2006}, tổn thất kỳ vọng của một khoản tín dụng được phân rã thành
tích của ba thành phần:

\begin{equation}
\EL = \PD \times \LGD \times \EAD,
\label{eq:expected-loss}
\end{equation}

trong đó $\PD$ (\eng{Probability of Default}) là xác suất vỡ nợ trong một khoảng
thời gian xác định -- thường là 12 tháng; $\LGD$ (\eng{Loss Given Default}) là tỷ
lệ tổn thất khi vỡ nợ xảy ra; và $\EAD$ (\eng{Exposure at Default}) là dư nợ tại
thời điểm vỡ nợ.

Phân rã này cho thấy vị trí của bài toán mà đề án giải quyết. Trong ba thành
phần, $\PD$ là đại lượng gắn trực tiếp với \textit{đặc điểm của khách hàng} và
được quyết định tại thời điểm cấp tín dụng; $\LGD$ và $\EAD$ phụ thuộc chủ yếu
vào cấu trúc khoản vay, tài sản bảo đảm và quá trình thu hồi nợ sau này. Vì vậy
$\PD$ là đại lượng mà hoạt động kiểm duyệt tín dụng có thể tác động và cũng là
đại lượng mà mô hình học máy có thể ước lượng từ dữ liệu lịch sử. Đề án tập trung
vào việc ước lượng $\PD$; việc mô hình hóa $\LGD$ và $\EAD$ không thuộc phạm vi
nghiên cứu.

\paragraph{Định nghĩa vỡ nợ.} Để ước lượng $\PD$, trước hết phải định nghĩa được
thế nào là ``vỡ nợ''. Thông tư 11/2021/TT-NHNN~\cite{tt11} quy định việc phân
loại nợ thành năm nhóm theo mức độ rủi ro, trong đó nợ thuộc nhóm 3 (Nợ dưới tiêu
chuẩn), nhóm 4 (Nợ nghi ngờ) và nhóm 5 (Nợ có khả năng mất vốn) được xếp vào
\textbf{nợ xấu}. Về mặt số ngày quá hạn, ranh giới của nợ nhóm 3 là 90 ngày. Định
nghĩa này trùng với ngưỡng mà Basel sử dụng cho khái niệm \eng{default}, nhờ đó
kết quả nghiên cứu có thể đối chiếu với các công trình quốc tế. Đề án sử dụng
thống nhất định nghĩa: một hồ sơ được gán nhãn \textit{xấu} nếu phát sinh tình
trạng quá hạn từ 90 ngày trở lên (tương đương nợ nhóm 3--5) trong cửa sổ quan sát
kể từ thời điểm cấp tín dụng.

\paragraph{Ba thế hệ phương pháp đánh giá.} Công tác đánh giá rủi ro tín dụng đã
trải qua ba giai đoạn phát triển. Giai đoạn đầu dựa trên \textit{phán đoán chuyên
gia}: cán bộ tín dụng vận dụng kinh nghiệm để đưa ra nhận định, linh hoạt với
trường hợp đặc thù nhưng thiếu nhất quán và khó mở rộng theo quy mô. Từ thập niên
1960, \textit{thẻ điểm thống kê} (\eng{scorecard}) dựa trên hồi quy logistic ra
đời, quy đổi mỗi đặc điểm khách hàng thành một số điểm; phương pháp này bảo đảm
tính nhất quán, cho phép kiểm định thống kê và dễ giải thích cho cơ quan quản lý.
Từ thập niên 2000, \textit{mô hình học máy} được áp dụng ngày càng rộng rãi nhờ
khả năng nắm bắt quan hệ phi tuyến và tương tác bậc cao mà không cần chỉ định
trước dạng hàm, đổi lại khó giải thích hơn và đòi hỏi kiểm định chặt chẽ hơn.

Ba thế hệ này không thay thế nhau hoàn toàn. Trong thực tiễn hiện nay, các tổ
chức tín dụng thường vận hành đồng thời cả ba: mô hình học máy cung cấp điểm rủi
ro, bộ luật nghiệp vụ bảo đảm tuân thủ chính sách và ngưỡng pháp lý, còn phán
đoán chuyên gia xử lý các hồ sơ nằm ngoài vùng mà mô hình đủ tin cậy. Hệ thống
được xây dựng trong đề án đi theo mô hình kết hợp này.

\subsection{Những hạn chế của phương pháp kiểm duyệt tín dụng truyền thống}
\label{subsec:han-che-truyen-thong}

Phương pháp kiểm duyệt truyền thống -- dựa trên bộ quy tắc nghiệp vụ cố định kết
hợp với phán đoán của cán bộ tín dụng -- bộc lộ bốn hạn chế chính.

\paragraph{Thứ nhất, không biểu diễn được tương tác giữa các yếu tố.} Một bộ quy
tắc có dạng ``nếu $\DTI > 0{,}6$ thì từ chối'' xét từng chỉ tiêu độc lập với các
chỉ tiêu còn lại. Cách làm này ngầm giả định rằng ảnh hưởng của mỗi yếu tố tới
rủi ro là cộng tính và không phụ thuộc giá trị của các yếu tố khác. Giả định đó
thường bị vi phạm trong thực tế: tác động của tỷ lệ $\DTI$ cao lên xác suất vỡ nợ
mạnh hơn nhiều ở nhóm khách hàng có thu nhập không ổn định so với nhóm có thu
nhập ổn định. Về mặt hình thức, bộ quy tắc chỉ tạo ra được các mặt phân chia song
song với trục tọa độ trong không gian đặc trưng, trong khi ranh giới rủi ro thực
tế có dạng phức tạp hơn nhiều.

\paragraph{Thứ hai, ngưỡng được đặt theo phán đoán chứ không theo dữ liệu.} Các
ngưỡng như ``$\DTI$ tối đa 60\%'' hay ``thâm niên công việc tối thiểu 6 tháng''
thường được kế thừa từ thông lệ hoặc từ kinh nghiệm của người xây dựng chính
sách. Chúng hiếm khi được hiệu chỉnh lại theo dữ liệu thực tế của chính tổ chức,
và càng hiếm khi được kiểm định xem mức ngưỡng đó có thực sự tối ưu về mặt đánh
đổi giữa quy mô cấp tín dụng và chất lượng danh mục hay không.

\paragraph{Thứ ba, kết quả đầu ra mang tính nhị phân, không có thang đo mức độ.}
Một bộ quy tắc cho biết hồ sơ ``đạt'' hay ``không đạt'' nhưng không cho biết hồ
sơ đạt ở mức nào. Điều này khiến tổ chức tín dụng không thể thực hiện hai nghiệp
vụ quan trọng: định giá theo rủi ro (khách hàng rủi ro thấp được hưởng lãi suất
ưu đãi hơn) và phân bổ nguồn lực thẩm định (tập trung nhân lực vào nhóm hồ sơ
thực sự cần xem xét kỹ).

\paragraph{Thứ tư, tính nhất quán và tốc độ xử lý bị giới hạn bởi yếu tố con
người.} Cùng một hồ sơ có thể nhận được kết luận khác nhau từ hai cán bộ khác
nhau, hoặc từ cùng một cán bộ ở hai thời điểm khác nhau. Khi khối lượng hồ sơ
tăng, thời gian xử lý kéo dài và áp lực công việc làm gia tăng nguy cơ sai sót.

Bốn hạn chế trên đều có chung một đặc điểm: chúng không phải là vấn đề về năng
lực chuyên môn mà là giới hạn cấu trúc của phương pháp. Do đó, hướng khắc phục
không phải là tăng cường đào tạo hay bổ sung thêm quy tắc, mà là bổ sung một
thành phần có khả năng học quan hệ từ dữ liệu và cho ra kết quả trên thang đo
liên tục. Đó chính là vai trò của mô hình học máy mà đề án xây
dựng.

\section{Tổng quan về học máy trong đánh giá tín dụng}
\label{sec:tong-quan-hoc-may}

\subsection{Khái niệm học máy}
\label{subsec:khai-niem-hoc-may}

Học máy (\eng{machine learning}) là lĩnh vực nghiên cứu các thuật toán cho phép
hệ thống máy tính cải thiện hiệu quả thực hiện một nhiệm vụ thông qua kinh
nghiệm, thay vì thông qua việc được lập trình tường minh từng bước. Định nghĩa
được trích dẫn rộng rãi của Mitchell~\cite{mitchell1997} phát biểu: một chương trình máy tính được
gọi là \textit{học} từ kinh nghiệm $E$ đối với lớp nhiệm vụ $T$ và độ đo hiệu quả
$P$, nếu hiệu quả của nó trên nhiệm vụ $T$, đo bằng $P$, cải thiện theo kinh
nghiệm $E$.

Áp dụng vào bài toán của đề án, ba thành phần này được xác định cụ thể như sau:

\begin{itemize}
  \item \textbf{Nhiệm vụ $T$}: dự báo khả năng một hồ sơ vay sẽ phát sinh nợ xấu
  trong 12 tháng kể từ thời điểm cấp tín dụng.
  \item \textbf{Kinh nghiệm $E$}: tập hồ sơ vay trong quá khứ, mỗi hồ sơ gồm các
  đặc trưng được quan sát tại thời điểm nộp và kết quả trả nợ đã biết.
  \item \textbf{Độ đo hiệu quả $P$}: khả năng phân biệt giữa hồ sơ tốt và hồ sơ
  xấu, đo bằng các chỉ tiêu như $\AUC$, hệ số Gini và thống kê $\KS$.
\end{itemize}

Điểm khác biệt căn bản giữa học máy và lập trình theo quy tắc nằm ở chiều của
quá trình suy luận. Lập trình theo quy tắc đi từ quy luật do con người phát biểu
tới kết quả; học máy đi ngược lại -- từ các cặp (dữ liệu đầu vào, kết quả đã
biết) để suy ra quy luật. Trong bài toán tín dụng, chiều ngược này có ý nghĩa
thực tiễn rõ rệt: thay vì yêu cầu chuyên gia phát biểu chính xác mối quan hệ giữa
hàng trăm đặc trưng và xác suất vỡ nợ (điều vượt quá khả năng của phán đoán chủ
quan), thuật toán rút ra quan hệ đó từ chính dữ liệu lịch sử của tổ chức.

\subsection{Phân loại các phương pháp học máy}
\label{subsec:phan-loai-hoc-may}

Các phương pháp học máy thường được phân loại theo dạng thông tin giám sát có
trong dữ liệu huấn luyện.

\paragraph{Học có giám sát (\eng{supervised learning}).} Dữ liệu huấn luyện gồm
các cặp $(\bm{x}_i, y_i)$, trong đó $\bm{x}_i$ là vectơ đặc trưng và $y_i$ là
nhãn đã biết. Mục tiêu là học một ánh xạ $f: \mathcal{X} \to \mathcal{Y}$ sao cho
$f(\bm{x})$ xấp xỉ tốt nhãn thật trên dữ liệu mới. Khi $\mathcal{Y}$ là tập hữu
hạn các lớp, bài toán là \textit{phân loại}; khi $\mathcal{Y}$ liên tục, bài toán
là \textit{hồi quy}.

\paragraph{Học không giám sát (\eng{unsupervised learning}).} Dữ liệu chỉ gồm
$\bm{x}_i$ mà không có nhãn. Thuật toán tìm cấu trúc tiềm ẩn trong dữ liệu: phân
cụm khách hàng thành các nhóm tương đồng, giảm chiều dữ liệu, hoặc phát hiện bất
thường. Trong lĩnh vực tín dụng, nhóm phương pháp này thường được dùng cho phân
khúc khách hàng và phát hiện gian lận.

\paragraph{Học bán giám sát (\eng{semi-supervised learning}).} Kết hợp một lượng
nhỏ dữ liệu có nhãn với lượng lớn dữ liệu không nhãn. Trong bài toán tín dụng,
tình huống này xuất hiện tự nhiên: các hồ sơ \textit{bị từ chối} không có kết quả
trả nợ nên không có nhãn, trong khi chúng chiếm một phần đáng kể tổng số hồ sơ
nộp vào. Vấn đề này được biết đến với tên gọi \eng{reject inference}.

\paragraph{Học tăng cường (\eng{reinforcement learning}).} Tác tử học chính sách
hành động thông qua tương tác với môi trường và tín hiệu thưởng/phạt. Phương pháp
này ít được dùng cho chấm điểm tín dụng do khó thiết lập môi trường thử nghiệm an
toàn, nhưng có ứng dụng trong tối ưu hạn mức và chiến lược thu hồi nợ.

Bài toán của đề án -- dự báo nhãn nhị phân ``tốt/xấu'' từ dữ liệu lịch sử đã có
kết quả -- thuộc nhóm \textbf{học có giám sát, bài toán phân loại nhị phân}. Tuy
nhiên, cần lưu ý một đặc thù: kết quả cần có không phải là nhãn dự báo
mà là \textit{xác suất} vỡ nợ được hiệu chuẩn đúng, vì xác suất này được dùng
trực tiếp để phân vùng quyết định và tính tổn thất kỳ vọng theo
công thức~\eqref{eq:expected-loss}.

\subsection{Học máy có giám sát trong bài toán tín dụng}
\label{subsec:hoc-co-giam-sat}

Phát biểu hình thức bài toán như sau. Cho tập huấn luyện
$\mathcal{D} = \{(\bm{x}_i, y_i)\}_{i=1}^{n}$ với $\bm{x}_i \in \mathbb{R}^d$ là
vectơ đặc trưng của hồ sơ thứ $i$ và $y_i \in \{0, 1\}$ là nhãn ($y_i = 1$ nếu hồ
sơ phát sinh nợ xấu). Cần tìm hàm $f$ ước lượng xác suất có điều kiện

\begin{equation}
\PD(\bm{x}) = \Pr(y = 1 \mid \bm{X} = \bm{x}).
\label{eq:pd-definition}
\end{equation}

Hàm $f$ được chọn bằng cách tối thiểu hóa một hàm mất mát trên tập huấn luyện, có
bổ sung thành phần điều chuẩn để kiểm soát độ phức tạp:

\begin{equation}
\hat{f} = \argmin_{f \in \mathcal{F}}
\left[ \frac{1}{n} \sum_{i=1}^{n} \ell\big(y_i, f(\bm{x}_i)\big)
       + \lambda \, \Omega(f) \right].
\label{eq:erm}
\end{equation}

Hàm mất mát chuẩn cho bài toán phân loại nhị phân có đầu ra xác suất là hàm
\textbf{entropy chéo nhị phân} (\eng{binary cross-entropy}, còn gọi là log-loss):

\begin{equation}
\ell(y, p) = -\big[\, y \log p + (1 - y) \log (1 - p) \,\big].
\label{eq:logloss}
\end{equation}

Log-loss thuộc lớp \textit{quy tắc chấm điểm chính đáng} (\eng{strictly proper
scoring rule}): kỳ vọng của nó đạt cực tiểu khi và chỉ khi $p$ bằng đúng xác suất
thật. Tính chất này có ý nghĩa quyết định đối với bài toán tín dụng, nơi giá trị
xác suất được sử dụng trực tiếp chứ không chỉ dùng để xếp hạng. Nếu thay bằng một
hàm mất mát không chính đáng, chẳng hạn tối ưu trực tiếp độ chính xác phân loại,
mô hình có thể xếp hạng tốt nhưng cho ra giá trị xác suất sai lệch hệ thống.

\paragraph{Ba đặc thù của bài toán tín dụng.} So với bài toán phân loại nhị phân
thông thường, bài toán chấm điểm tín dụng có ba đặc thù cần được xử lý riêng:

\begin{enumerate}[label=(\arabic*)]
  \item \textbf{Dữ liệu mất cân bằng.} Tỷ lệ hồ sơ xấu thường ở mức 3--20\%, tùy
  sản phẩm và phân khúc. Mức mất cân bằng này khiến chỉ tiêu độ chính xác
  (\eng{accuracy}) mất ý nghĩa: một mô hình dự báo ``mọi hồ sơ đều tốt'' vẫn đạt
  độ chính xác 85\% mà không có giá trị sử dụng nào.

  \item \textbf{Chi phí hai loại sai lầm rất khác nhau.} Bỏ sót một hồ sơ xấu
  (\eng{false negative}) gây tổn thất bằng một phần đáng kể dư nợ, trong khi từ
  chối sai một hồ sơ tốt (\eng{false positive}) chỉ làm mất phần lãi biên. Tỷ lệ
  chi phí thường được ước lượng trong khoảng 10:1 đến 15:1. Hệ quả là ngưỡng
  quyết định phải được chọn theo hàm chi phí~\cite{elkan2001} chứ không phải theo
  giá trị mặc định 0,5.

  \item \textbf{Trôi dạt theo thời gian.} Phân phối đặc trưng và quan hệ giữa đặc
  trưng với nhãn thay đổi theo chu kỳ kinh tế, theo chính sách sản phẩm và theo
  thành phần khách hàng~\cite{gama2014}. Do đó việc đánh giá mô hình phải được
  thực hiện trên dữ liệu của giai đoạn \textit{sau} giai đoạn huấn luyện, chứ
  không phải trên mẫu ngẫu nhiên trộn lẫn các giai đoạn.
\end{enumerate}

\subsection{Quy trình xây dựng mô hình học máy}
\label{subsec:quy-trinh-xay-dung-mo-hinh}

Quy trình xây dựng một mô hình học máy phục vụ nghiệp vụ tín dụng gồm bảy bước,
tổ chức thành ba khối như minh họa ở Hình~\ref{fig:ml-pipeline}: khối
\textit{chuẩn bị dữ liệu} (bước 1--3), khối \textit{xây dựng mô hình} (bước 4--6)
và khối \textit{vận hành} (bước 7).

\begin{figure}[htbp]
\centering
\begin{tikzpicture}[
  font=\small,
  module/.style={draw=blue!45!black, fill=blue!9, rounded corners=3pt,
                 line width=0.6pt, blur shadow={shadow blur steps=5}},
  panel/.style={draw=black!55, rounded corners=3pt, line width=0.5pt},
  inner/.style={draw=black!60, fill=white, rounded corners=2pt,
                align=center, font=\footnotesize, inner sep=2pt,
                minimum height=0.95cm, text width=2.35cm},
  flow/.style={-{Stealth[length=2.2mm]}, line width=0.8pt, draw=black!70},
  up/.style={-{Triangle[length=2.2mm, width=4.2mm]}, line width=2.6pt,
             draw=orange!75!yellow},
]
% ---------------- khối trung tâm: xây dựng mô hình -----------------------
\draw[module] (0,0) rectangle (8.6,7.2);
\node[font=\small\bfseries] at (4.3,6.88) {Khối xây dựng mô hình học máy};

% tầng trên: bước (4)
\draw[panel, fill=yellow!28] (0.3,4.55) rectangle (8.3,6.55);
\node[font=\small] at (4.3,6.25) {(4) Trích xuất và lựa chọn đặc trưng};
\node[inner] at (1.60,5.22) {Đặc trưng\\theo cửa sổ};
\node[inner] at (4.30,5.22) {Tỷ số và\\tương tác};
\node[inner] at (7.00,5.22) {Lọc IV, PSI,\\tương quan};

% tầng giữa: bước (5)
\node[draw=black!55, fill=orange!18, rounded corners=3pt, align=center,
      minimum width=7.0cm, minimum height=1.0cm] (train) at (4.3,3.55)
      {(5) Huấn luyện và tối ưu mô hình\\[-1pt]
       {\footnotesize trọng số lớp · dừng sớm · tìm kiếm ngẫu nhiên}};

% tầng dưới: bước (6)
\draw[panel, fill=green!22] (0.3,0.3) rectangle (8.3,2.55);
\node[font=\small] at (4.3,2.25) {(6) Đánh giá và hiệu chuẩn};
\node[inner] at (1.60,1.15) {AUC, Gini,\\KS, ECE};
\node[inner] at (4.30,1.15) {Hiệu chuẩn\\xác suất};
\node[inner] at (7.00,1.15) {Chọn ngưỡng\\theo chi phí};

% mũi tên giữa các tầng
\draw[up] (4.3,4.55) -- (4.3,4.09);
\draw[up] (4.3,3.03) -- (4.3,2.58);

% ---------------- khối trái: chuẩn bị dữ liệu -----------------------------
\draw[module] (-2.3,0) rectangle (-0.55,7.2);
\node[rotate=90, align=center] at (-1.425,3.6)
  {\textbf{Khối chuẩn bị dữ liệu}\\[2pt]
   {\footnotesize (1) Xác định bài toán và biến mục tiêu}\\[1pt]
   {\footnotesize (2) Thu thập \,·\, (3) Tiền xử lý, làm sạch}};

% ---------------- khối phải: vận hành -------------------------------------
\draw[module] (9.15,0) rectangle (10.9,7.2);
\node[rotate=90, align=center] at (10.025,3.6)
  {\textbf{Khối vận hành}\\[2pt]
   {\footnotesize (7) Triển khai và giám sát}\\[1pt]
   {\footnotesize dịch vụ REST \,·\, giám sát PSI}};

% ---------------- đầu vào / đầu ra ---------------------------------------
\node[rotate=90, font=\footnotesize] at (-3.0,3.6) {Dữ liệu hồ sơ tín dụng};
\node[rotate=90, font=\footnotesize] at (11.6,3.6) {PD, khuyến nghị và lý do};
\draw[flow] (-2.77,3.6) -- (-2.3,3.6);
\draw[flow] (-0.55,3.6) -- (0,3.6);
\draw[flow] (8.6,3.6) -- (9.15,3.6);
\draw[flow] (10.9,3.6) -- (11.37,3.6);

% ---------------- vòng phản hồi -------------------------------------------
\draw[flow, dashed, draw=black!55] (10.025,0) -- (10.025,-0.55)
  -- (-1.425,-0.55) -- (-1.425,0);
\node[font=\footnotesize\itshape, fill=white, inner sep=2pt] at (4.3,-0.55)
  {cập nhật dữ liệu và huấn luyện lại khi phát hiện suy giảm};
\end{tikzpicture}
\caption{Quy trình xây dựng và vận hành mô hình học máy cho bài toán tín dụng}
\label{fig:ml-pipeline}
\end{figure}

Bước (1) định nghĩa thế nào là hồ sơ ``xấu'', xác định cửa sổ quan sát và phạm vi
khách hàng áp dụng; đây là bước có ảnh hưởng lớn nhất tới toàn bộ kết quả, vì một
định nghĩa nhãn không phù hợp sẽ làm sai lệch mọi bước tiếp theo. Các bước (2)
đến (6) lần lượt tổng hợp dữ liệu về cùng một đơn vị phân tích, làm sạch, tạo đặc
trưng, huấn luyện và đánh giá mô hình. Bước (7) theo dõi mô hình trong vận hành để
xác định thời điểm huấn luyện lại. Nội dung chi tiết của từng bước, gắn với bài
toán cụ thể của đề án, được trình bày ở Chương~\ref{chap:xay-dung-he-thong}.

Quy trình này có \textit{tính lặp}: kết quả ở bước sau thường buộc phải quay lại
điều chỉnh các bước trước. Chẳng hạn, một đặc trưng có chỉ số Information Value
cao bất thường ở bước (4) có thể là dấu hiệu biến đó chứa thông tin của tương
lai, buộc phải xem lại cách tổng hợp dữ liệu ở bước (2).

\subsection{Vai trò của học máy trong hỗ trợ kiểm duyệt tín dụng}
\label{subsec:vai-tro-hoc-may}

Trong quy trình kiểm duyệt mà đề án xây dựng, mô hình học máy đảm nhiệm bốn vai
trò. Thứ
nhất, \textit{lượng hóa rủi ro trên thang đo liên tục}: thay vì kết luận nhị phân
``đạt/không đạt'', mô hình cho ra giá trị $\PD \in (0, 1)$, nhờ đó có thể phân
hạng rủi ro, định giá theo rủi ro và phân bổ nguồn lực thẩm định theo mức độ cần
thiết. Thứ hai, \textit{phát hiện quan hệ mà bộ quy tắc bỏ sót}, kể cả những
tương tác mà người xây dựng chính sách không dự liệu trước. Thứ ba, \textit{bảo
đảm tính nhất quán}: cùng một hồ sơ luôn nhận được cùng một kết quả đánh giá,
không phụ thuộc người xử lý hay thời điểm xử lý. Thứ tư, \textit{giảm khối lượng
công việc thủ công} bằng cách tự động xử lý các hồ sơ nằm ở hai vùng biên, để cán
bộ tín dụng tập trung vào nhóm hồ sơ thực sự cần cân nhắc.

Đi kèm là bốn giới hạn. Mô hình chỉ học được quy luật có trong dữ liệu lịch sử,
nên những tình huống chưa từng xuất hiện nằm ngoài khả năng dự báo. Chất lượng mô
hình phụ thuộc trực tiếp vào chất lượng dữ liệu. Các mô hình phức tạp khó giải
thích hơn thẻ điểm truyền thống, đòi hỏi bổ sung công cụ giải thích. Cuối cùng,
mô hình suy giảm theo thời gian và cần được giám sát, hiệu chỉnh định
kỳ~\cite{sculley2015}. Bốn giới hạn này là cơ sở để đề án xác định vai trò của mô
hình là \textit{hỗ trợ} chứ không phải \textit{thay thế} quá trình kiểm duyệt của
con người.

\section{Dữ liệu và đặc trưng trong bài toán tín dụng}
\label{sec:du-lieu-dac-trung}

\subsection{Đặc điểm dữ liệu tín dụng}
\label{subsec:dac-diem-du-lieu}

Dữ liệu phục vụ bài toán chấm điểm tín dụng có năm đặc điểm gây khó khăn riêng so
với dữ liệu của nhiều lĩnh vực khác.

\paragraph{Tính đúng đắn theo thời điểm.} Mọi đặc trưng của một hồ sơ chỉ được
phép tính từ dữ liệu đã tồn tại \textit{trước hoặc tại} thời điểm nộp hồ sơ. Vi
phạm nguyên tắc này dẫn tới hiện tượng rò rỉ nhãn (\eng{label leakage}): mô hình
sử dụng thông tin của tương lai, cho kết quả đánh giá cao giả tạo trong môi
trường thử nghiệm nhưng suy giảm nghiêm trọng khi triển khai. Đây là lỗi khó phát
hiện, vì nó không sinh ra thông báo lỗi mà chỉ làm các chỉ tiêu chất lượng cao
lên bất thường.

\paragraph{Độ trễ của nhãn.} Kết quả trả nợ chỉ được biết sau cửa sổ quan sát,
thường là 12 tháng. Do đó tại mọi thời điểm luôn tồn tại một khối hồ sơ gần đây
chưa có nhãn. Điều này giới hạn lượng dữ liệu huấn luyện khả dụng và buộc phải
cân nhắc đánh đổi: cửa sổ quan sát dài hơn cho nhãn đáng tin cậy hơn nhưng làm dữ
liệu huấn luyện cũ hơn.

\paragraph{Mất cân bằng lớp.} Tỷ lệ hồ sơ xấu thường thấp, khiến số quan sát của
lớp thiểu số -- vốn là lớp cần quan tâm -- bị hạn chế.

\paragraph{Giá trị khuyết mang thông tin.} Trong dữ liệu tín dụng, việc một
trường bị trống thường không phải lỗi thu thập mà tự nó là một tín hiệu. Khách
hàng không có bất kỳ thông tin nào tại CIC (\eng{thin-file}) nghĩa là chưa từng
có quan hệ tín dụng -- đây là một trạng thái có ý nghĩa rủi ro riêng, khác hẳn
với việc không truy vấn được dữ liệu. Cách xử lý đúng là phân biệt hai
trường hợp và bổ sung biến chỉ báo khuyết thay vì đơn thuần điền giá trị thống
kê.

\paragraph{Trôi dạt theo thời gian.} Cả phân phối biên của đặc trưng lẫn quan hệ
giữa đặc trưng và nhãn đều thay đổi theo thời gian, do biến động kinh tế vĩ mô,
thay đổi chính sách sản phẩm hoặc dịch chuyển cơ cấu khách hàng.

\subsection{Các nhóm dữ liệu sử dụng trong đánh giá tín dụng}
\label{subsec:nhom-du-lieu}

Dữ liệu phục vụ đánh giá tín dụng khách hàng cá nhân được phân thành bốn nhóm
nguồn, tổng hợp tại Bảng~\ref{tab:nhom-du-lieu}.

\begin{table}[htbp]
\centering
\caption{Các nhóm nguồn dữ liệu sử dụng trong đánh giá tín dụng cá nhân}
\label{tab:nhom-du-lieu}
\small
\begin{tabular}{@{}L{3.0cm}L{4.6cm}L{4.6cm}@{}}
\toprule
\textbf{Nhóm} & \textbf{Nội dung chính} & \textbf{Đặc điểm} \\
\midrule
Thông tin cá nhân và tài chính &
Tuổi, giới tính, tình trạng hôn nhân, số người phụ thuộc, trình độ, nghề nghiệp,
thâm niên, thu nhập, nơi cư trú &
Khách hàng tự khai báo; cần đối chiếu với chứng từ. Thay đổi chậm theo thời gian
\\[2pt]
Tham số khoản vay &
Sản phẩm, số tiền, kỳ hạn, lãi suất, mục đích vay, kênh tiếp nhận, tài sản bảo
đảm &
Xác định tại thời điểm nộp hồ sơ; phản ánh cả nhu cầu lẫn mức độ chấp nhận rủi ro
của khách hàng \\[2pt]
Lịch sử tín dụng và quan hệ CIC &
Điểm CIC, nhóm nợ, số ngày quá hạn, số tổ chức có quan hệ, dư nợ, tỷ lệ sử dụng
hạn mức, số lần truy vấn &
Nguồn thông tin khách quan, có năng lực dự báo cao nhất; phụ thuộc dịch vụ bên
ngoài và không có với khách hàng thin-file \\[2pt]
Hành vi giao dịch tài khoản &
Số lượng và giá trị giao dịch, dòng tiền vào/ra, tính đều đặn của thu nhập, cơ
cấu chi tiêu theo nhóm &
Chỉ có với khách hàng đã có quan hệ tài khoản; cập nhật liên tục và phản ánh tình
trạng tài chính hiện tại \\
\bottomrule
\end{tabular}
\end{table}

Trong bốn nhóm trên, nhóm dữ liệu hành vi giao dịch có ba đặc điểm riêng: tổ chức
tín dụng \textit{tự sở hữu} nên không phụ thuộc bên thứ ba, phản ánh tình trạng
tài chính \textit{hiện tại} thay vì lịch sử, và luôn có sẵn kể cả với khách hàng
thiếu lịch sử tín dụng. Onay và Öztürk~\cite{onay2018} nhấn mạnh vai trò ngày
càng lớn của nhóm dữ liệu này trong bối cảnh dữ liệu lớn, đặc biệt với phân khúc
khách hàng mà các nguồn dữ liệu truyền thống gần như không cung cấp thông tin.
Một trong những nội dung thực nghiệm của đề án là định lượng mức đóng góp của
nhóm dữ liệu này vào chất lượng dự báo.

\subsection{Xây dựng và lựa chọn đặc trưng}
\label{subsec:xay-dung-dac-trung}

\paragraph{Xây dựng đặc trưng.} Dữ liệu thô hiếm khi ở dạng phù hợp trực tiếp cho
mô hình. Quá trình xây dựng đặc trưng (\eng{feature engineering}) biến đổi dữ
liệu thô thành các biến mang ý nghĩa dự báo, thông qua bốn kỹ thuật chính:

\begin{itemize}
  \item \textbf{Tổng hợp theo cửa sổ thời gian}: từ danh sách giao dịch, tính
  tổng dòng tiền vào, số lần giao dịch, giá trị trung bình trong các cửa sổ 7,
  30, 90 và 180 ngày trước thời điểm nộp hồ sơ.

  \item \textbf{Tạo tỷ số}: các chỉ tiêu nghiệp vụ như $\DTI$, $\LTV$, tỷ lệ thu
  nhập còn lại sau trả nợ. Tỷ số có ưu điểm là chuẩn hóa được sự khác biệt về quy
  mô giữa các khách hàng.

  \item \textbf{Tạo biến xu thế}: so sánh giá trị ở cửa sổ ngắn với cửa sổ dài để
  nắm bắt chiều hướng thay đổi, chẳng hạn dòng tiền vào 30 ngày gần nhất so với
  mức trung bình 90 ngày.

  \item \textbf{Tạo biến tương tác}: kết hợp hai hoặc nhiều biến để biểu diễn
  quan hệ phi cộng tính, chẳng hạn tích của số ngày quá hạn với số tổ chức tín
  dụng có quan hệ.
\end{itemize}

\paragraph{Weight of Evidence và Information Value.} Trong ngành tín dụng, hai
công cụ được sử dụng phổ biến để đánh giá và biến đổi đặc trưng là Weight of
Evidence ($\WOE$) và Information Value ($\IV$)~\cite{siddiqi2017}. Với một biến
được rời rạc hóa thành $K$ khoảng, $\WOE$ của khoảng thứ $k$ được định nghĩa:

\begin{equation}
\WOE_k = \ln \frac{\text{tỷ lệ hồ sơ tốt trong khoảng } k}
                  {\text{tỷ lệ hồ sơ xấu trong khoảng } k}
       = \ln \frac{g_k / G}{b_k / B},
\label{eq:woe}
\end{equation}

trong đó $g_k, b_k$ lần lượt là số hồ sơ tốt và xấu trong khoảng $k$; $G, B$ là
tổng số hồ sơ tốt và xấu. Information Value tổng hợp năng lực phân biệt của toàn
bộ biến:

\begin{equation}
\IV = \sum_{k=1}^{K} \left( \frac{g_k}{G} - \frac{b_k}{B} \right) \WOE_k.
\label{eq:iv}
\end{equation}

Thực hành ngành sử dụng thang diễn giải: $\IV < 0{,}02$ -- biến không có khả năng
phân biệt; $0{,}02 \le \IV < 0{,}1$ -- yếu; $0{,}1 \le \IV < 0{,}3$ -- trung
bình; $0{,}3 \le \IV < 0{,}5$ -- mạnh; $\IV \ge 0{,}5$ -- rất mạnh, cần kiểm tra
khả năng rò rỉ nhãn.

Phép biến đổi $\WOE$ có ba ưu điểm đối với bài toán tín dụng: nó đưa mọi biến về
cùng một thang đo theo log-odds, xử lý được giá trị khuyết như một khoảng riêng,
và cho phép áp đặt ràng buộc đơn điệu để bảo đảm mô hình phù hợp với trực giác
nghiệp vụ (điểm tín dụng càng cao thì rủi ro càng thấp).

\paragraph{Lựa chọn đặc trưng.} Khi tập đặc trưng ban đầu lớn, việc giữ lại toàn
bộ gây ba vấn đề: hệ số của mô hình tuyến tính mất ổn định do đa cộng tuyến, chi
phí tính toán tăng, và khả năng diễn giải giảm. Các phương pháp lựa chọn đặc
trưng được chia thành ba nhóm:

\begin{itemize}
  \item \textbf{Phương pháp lọc} (\eng{filter}): đánh giá từng đặc trưng độc lập
  với mô hình, dựa trên $\IV$, hệ số tương quan hoặc kiểm định thống kê. Chi phí
  tính toán thấp, phù hợp để sàng lọc sơ bộ.

  \item \textbf{Phương pháp bao} (\eng{wrapper}): thử các tập con đặc trưng và
  đánh giá bằng chính mô hình đích. Cho kết quả tốt hơn nhưng chi phí tính toán
  cao.

  \item \textbf{Phương pháp nhúng} (\eng{embedded}): việc lựa chọn diễn ra trong
  quá trình huấn luyện, chẳng hạn điều chuẩn $L_1$ hoặc tầm quan trọng đặc trưng
  của mô hình cây.
\end{itemize}

Đề án sử dụng quy trình lọc ba bước nối tiếp -- theo năng lực dự báo ($\IV$),
theo mức tương quan, và theo độ ổn định theo thời gian ($\PSI$) -- được trình bày
chi tiết ở mục~\ref{sec:trich-xuat-dac-trung}.

\section{Các mô hình học máy áp dụng trong đánh giá tín dụng}
\label{sec:cac-mo-hinh}

Mục này trình bày nguyên lý của năm nhóm thuật toán được đề án nghiên cứu và thực
nghiệm, theo đúng thứ tự từ đơn giản đến phức tạp.

\subsection{Hồi quy Logistic}
\label{subsec:hoi-quy-logistic}

Hồi quy logistic là mô hình tuyến tính tổng quát dùng cho biến phụ thuộc nhị
phân. Mô hình giả định log-odds của xác suất vỡ nợ là hàm tuyến tính của các đặc
trưng:

\begin{equation}
\ln \frac{\PD(\bm{x})}{1 - \PD(\bm{x})} = \beta_0 + \sum_{j=1}^{d} \beta_j x_j,
\qquad\text{tương đương}\qquad
\PD(\bm{x}) = \frac{1}{1 + e^{-(\beta_0 + \bm{\beta}^{\!\top}\bm{x})}}.
\label{eq:logistic}
\end{equation}
\vspace{-0.6\baselineskip}

Tham số $\bm{\beta}$ được ước lượng bằng phương pháp hợp lý cực đại, tương đương
với việc tối thiểu hóa hàm log-loss~\eqref{eq:logloss}. Để kiểm soát quá khớp,
thường bổ sung thành phần điều chuẩn $L_2$:

\begin{equation}
\hat{\bm{\beta}} = \argmin_{\bm{\beta}}
\left[ -\sum_{i=1}^{n} \Big( y_i \log p_i + (1-y_i)\log(1-p_i) \Big)
       + \frac{1}{2C} \lVert \bm{\beta} \rVert_2^2 \right].
\label{eq:logistic-l2}
\end{equation}

\paragraph{Ưu điểm.} Hồi quy logistic có ba ưu thế khiến nó vẫn là mô hình chuẩn
mực trong ngành tín dụng sau hơn nửa thế kỷ. Thứ nhất, \textit{khả năng giải
thích}: mỗi hệ số $\beta_j$ cho biết trực tiếp mức thay đổi của log-odds khi đặc
trưng $x_j$ tăng một đơn vị, và dấu của hệ số có thể được đối chiếu với kỳ vọng
nghiệp vụ. Thứ hai, \textit{xác suất đầu ra được hiệu chuẩn tốt} do hàm mất mát
chính là log-loss. Thứ ba, \textit{tính ổn định}: mô hình có ít tham số nên ít
nhạy cảm với nhiễu trong dữ liệu huấn luyện.

\paragraph{Hạn chế.} Giả định tuyến tính trên thang log-odds là hạn chế cơ bản:
mô hình không tự nắm bắt được quan hệ phi tuyến hay tương tác giữa các biến, trừ
khi những quan hệ này được đưa vào tường minh dưới dạng biến phái sinh.

\paragraph{Thẻ điểm dựa trên biến đổi WOE.} Trong thực tiễn ngân hàng, hồi quy
logistic thường được áp dụng trên các đặc trưng đã qua biến đổi
$\WOE$~\eqref{eq:woe} thay vì trên giá trị gốc. Cách làm này khắc phục một phần
hạn chế tuyến tính -- vì việc rời rạc hóa cho phép biểu diễn quan hệ phi tuyến
theo từng khoảng -- đồng thời cho phép quy đổi kết quả thành thang điểm nguyên
dễ vận hành:

\begin{equation}
\text{Điểm} = \text{Offset} + \text{Factor} \times \ln(\text{odds}),
\qquad
\text{Factor} = \frac{\text{PDO}}{\ln 2},
\label{eq:scorecard}
\end{equation}

trong đó PDO (\eng{Points to Double the Odds}) là số điểm tương ứng với việc tỷ
số odds tăng gấp đôi. Thẻ điểm xây dựng theo cách này có thể in ra thành bảng
tra, giúp cán bộ tín dụng hiểu chính xác vì sao một hồ sơ đạt mức điểm cụ thể.

\subsection{Decision Tree}
\label{subsec:decision-tree}

Cây quyết định~\cite{breiman1984} phân chia đệ quy không gian đặc trưng thành các
vùng con, tại mỗi
bước chọn cặp (đặc trưng, ngưỡng) làm giảm nhiều nhất độ không thuần khiết của
nút. Với bài toán phân loại, độ không thuần khiết thường được đo bằng chỉ số
Gini:

\begin{equation}
G(t) = 1 - \sum_{c \in \{0,1\}} p_c(t)^2,
\label{eq:gini-impurity}
\end{equation}

trong đó $p_c(t)$ là tỷ lệ quan sát thuộc lớp $c$ tại nút $t$. Phép chia tốt nhất
là phép chia làm cực đại mức giảm độ không thuần khiết có trọng số:

\begin{equation}
\Delta G = G(t) - \frac{n_L}{n_t} G(t_L) - \frac{n_R}{n_t} G(t_R).
\label{eq:gini-gain}
\end{equation}

\paragraph{Ưu điểm.} Cây quyết định tự động nắm bắt quan hệ phi tuyến và tương
tác giữa các biến, không yêu cầu chuẩn hóa thang đo, xử lý được cả biến định
lượng lẫn định tính, và -- với cây nông -- có thể diễn giải trực tiếp thành chuỗi
điều kiện gần với ngôn ngữ nghiệp vụ.

\paragraph{Hạn chế.} Cây đơn có phương sai cao: một thay đổi nhỏ trong dữ liệu
huấn luyện có thể dẫn tới cấu trúc cây hoàn toàn khác. Cây sâu dễ quá khớp, trong
khi cây nông lại quá đơn giản để nắm bắt quan hệ phức tạp. Để kiểm soát, đề án sử
dụng kỹ thuật tỉa nhánh theo độ phức tạp chi phí (\eng{cost-complexity pruning})
với tiêu chí:

\begin{equation}
R_\alpha(T) = R(T) + \alpha \, |\tilde{T}|,
\label{eq:ccp}
\end{equation}

trong đó $R(T)$ là sai số của cây $T$, $|\tilde{T}|$ là số nút lá và $\alpha$ là
tham số phạt độ phức tạp.

Chính hạn chế về phương sai cao của cây đơn là động lực dẫn tới hai nhóm phương
pháp học tập hợp trình bày ở hai mục tiếp theo.

\subsection{Random Forest}
\label{subsec:random-forest}

Rừng ngẫu nhiên~\cite{breiman2001} xây dựng nhiều cây quyết định độc lập và tổng
hợp kết quả bằng cách lấy trung bình xác suất dự báo:

\begin{equation}
\PD_{\text{RF}}(\bm{x}) = \frac{1}{M} \sum_{m=1}^{M} \PD_m(\bm{x}).
\label{eq:rf}
\end{equation}

Tính đa dạng giữa các cây -- yếu tố quyết định hiệu quả của phương pháp -- được
tạo ra bằng hai cơ chế ngẫu nhiên hóa. Thứ nhất, mỗi cây được huấn luyện trên một
mẫu bootstrap lấy có hoàn lại từ tập huấn luyện. Thứ hai, tại mỗi nút, chỉ một
tập con ngẫu nhiên gồm $\sqrt{d}$ đặc trưng được xem xét để tìm phép chia.

Cơ sở lý thuyết của phương pháp~\cite{hastie2009} nằm ở công thức phương sai của trung bình $M$
biến ngẫu nhiên có cùng phương sai $\sigma^2$ và tương quan cặp $\rho$:

\begin{equation}
\operatorname{Var}\!\left( \frac{1}{M}\sum_{m=1}^{M} Z_m \right)
= \rho \sigma^2 + \frac{1 - \rho}{M}\sigma^2.
\label{eq:rf-variance}
\end{equation}

Số hạng thứ hai giảm về 0 khi $M$ tăng, nhưng số hạng thứ nhất chỉ giảm được bằng
cách làm giảm tương quan $\rho$ giữa các cây. Đây chính là lý do của cơ chế lấy
mẫu ngẫu nhiên đặc trưng tại mỗi nút: nếu mọi cây đều được phép dùng đặc trưng
mạnh nhất, chúng sẽ giống nhau và việc lấy trung bình mất tác dụng.

\paragraph{Ưu điểm.} Rừng ngẫu nhiên giữ được khả năng biểu diễn phi tuyến của
cây nhưng có phương sai thấp hơn nhiều, ít nhạy cảm với siêu tham số, và cung cấp
sẵn thước đo tầm quan trọng của đặc trưng.

\paragraph{Hạn chế.} Mô hình mất khả năng diễn giải trực tiếp; chi phí bộ nhớ và
thời gian suy luận tăng theo số cây; và xác suất đầu ra thường bị lệch về phía
trung tâm, cần hiệu chuẩn lại trước khi dùng làm $\PD$.

\subsection{Support Vector Machine}
\label{subsec:svm}

Máy vectơ hỗ trợ (SVM)~\cite{cortes1995,vapnik1998} tiếp cận bài toán phân loại
từ góc nhìn hình học: tìm siêu
phẳng phân tách hai lớp sao cho \textit{lề} -- khoảng cách từ siêu phẳng tới các
điểm gần nhất của mỗi lớp -- là lớn nhất. Với dữ liệu không phân tách tuyến tính
hoàn hảo, bài toán được phát biểu ở dạng lề mềm:

\begin{equation}
\begin{aligned}
&\min_{\bm{w}, b, \bm{\xi}} \quad
  \frac{1}{2}\lVert \bm{w} \rVert^2 + C \sum_{i=1}^{n} \xi_i \\[2pt]
&\text{với ràng buộc} \quad
  y_i(\bm{w}^{\!\top}\bm{x}_i + b) \ge 1 - \xi_i, \quad \xi_i \ge 0,
\end{aligned}
\label{eq:svm-primal}
\end{equation}

trong đó $\xi_i$ là biến bù cho phép một số điểm nằm sai phía, còn $C$ điều tiết
đánh đổi giữa độ rộng lề và mức chấp nhận vi phạm.

\paragraph{Thủ thuật hạt nhân.} Điểm mạnh đặc trưng của SVM nằm ở chỗ bài toán
đối ngẫu chỉ phụ thuộc vào tích vô hướng giữa các cặp quan sát. Nhờ đó có thể
thay tích vô hướng bằng một hàm hạt nhân $K(\bm{x}_i, \bm{x}_j)$ tương ứng với
tích vô hướng trong một không gian đặc trưng có số chiều cao hơn, mà không cần
tính toán tường minh phép ánh xạ. Hạt nhân được dùng phổ biến nhất là hạt nhân
Gauss (RBF):

\begin{equation}
K(\bm{x}_i, \bm{x}_j) = \exp\!\big( -\gamma \lVert \bm{x}_i - \bm{x}_j \rVert^2 \big).
\label{eq:rbf}
\end{equation}

\paragraph{Ưu điểm.} SVM có nền tảng lý thuyết chặt chẽ dựa trên nguyên lý cực
tiểu hóa rủi ro cấu trúc, hoạt động tốt trong không gian nhiều chiều, và nhờ thủ
thuật hạt nhân có thể biểu diễn ranh giới quyết định rất linh hoạt.

\paragraph{Hạn chế trong bài toán tín dụng.} SVM gặp ba trở ngại thực tế đáng kể
khi áp dụng cho chấm điểm tín dụng ở quy mô lớn:

\begin{enumerate}[label=(\arabic*)]
  \item \textbf{Chi phí tính toán.} Việc giải bài toán đối ngẫu có độ phức tạp từ
  $O(n^2)$ đến $O(n^3)$ theo số quan sát, khiến SVM hạt nhân đầy đủ không khả thi
  với tập dữ liệu hàng chục nghìn hồ sơ trở lên.

  \item \textbf{Đầu ra không phải xác suất.} SVM cho ra giá trị hàm quyết định --
  khoảng cách có dấu tới siêu phẳng -- chứ không phải xác suất. Để dùng được
  trong hệ thống tín dụng, giá trị này phải được quy đổi bằng phương pháp
  Platt~\cite{platt1999}: khớp một hồi quy logistic một biến trên giá trị hàm
  quyết định.

  \item \textbf{Khó giải thích.} Ranh giới quyết định trong không gian hạt nhân
  không có biểu diễn tường minh theo các đặc trưng gốc.
\end{enumerate}

Để khắc phục trở ngại thứ nhất, đề án sử dụng phép xấp xỉ
Nystroem~\cite{williams2001}: thay vì tính
ma trận hạt nhân đầy đủ kích thước $n \times n$, phương pháp chọn $m \ll n$ điểm
mốc và xây dựng một ánh xạ tường minh $\phi: \mathbb{R}^d \to \mathbb{R}^m$ sao
cho $\phi(\bm{x}_i)^{\!\top}\phi(\bm{x}_j) \approx K(\bm{x}_i, \bm{x}_j)$. Sau
phép ánh xạ này, một SVM tuyến tính -- có chi phí tuyến tính theo $n$ -- cho
nghiệm xấp xỉ nghiệm của SVM hạt nhân đầy đủ.

\subsection{Gradient Boosting}
\label{subsec:gradient-boosting}

Khác với rừng ngẫu nhiên xây dựng các cây \textit{song song và độc lập}, gradient
boosting~\cite{friedman2001} xây dựng các cây \textit{tuần tự}, trong đó mỗi cây
mới học phần sai số còn lại của tổ hợp các cây trước đó. Mô hình được xây dựng
theo sơ đồ cộng tính:

\begin{equation}
F_m(\bm{x}) = F_{m-1}(\bm{x}) + \eta \, h_m(\bm{x}),
\label{eq:boosting}
\end{equation}

trong đó $h_m$ là cây thứ $m$ được khớp với gradient âm của hàm mất mát tại dự
báo hiện tại, còn $\eta \in (0, 1]$ là tốc độ học đóng vai trò điều chuẩn: giá
trị $\eta$ nhỏ buộc mô hình cải thiện từ từ, giảm nguy cơ quá khớp nhưng đòi hỏi
nhiều cây hơn.

\paragraph{Các biến thể hiện đại.} Ba thư viện được sử dụng rộng rãi hiện nay cải
tiến ý tưởng gốc theo những hướng khác nhau:

\begin{itemize}
  \item \textbf{XGBoost}~\cite{chen2016} bổ sung thành phần điều chuẩn tường minh
  vào hàm mục tiêu và sử dụng khai triển Taylor bậc hai để xác định phép chia tối
  ưu, thay vì chỉ dùng gradient bậc nhất.

  \item \textbf{LightGBM}~\cite{ke2017} sinh cây theo chiều lá
  (\eng{leaf-wise}) -- mở rộng lá có mức giảm mất mát lớn nhất -- thay vì theo
  chiều sâu, kết hợp với kỹ thuật rời rạc hóa đặc trưng thành histogram để tăng
  tốc độ huấn luyện.

  \item \textbf{CatBoost}~\cite{prokhorenkova2018} xử lý đặc trưng định tính bằng
  thống kê mục tiêu có thứ tự (\eng{ordered target statistics}), khắc phục hiện
  tượng rò rỉ mục tiêu vốn xảy ra khi mã hóa biến định tính theo giá trị trung
  bình của nhãn.
\end{itemize}

\paragraph{Ưu điểm.} Gradient boosting thường cho chất lượng dự báo cao nhất
trong các bài toán dữ liệu dạng bảng, nhờ khả năng nắm bắt tương tác bậc cao và
cơ chế tập trung vào các quan sát khó phân loại.

\paragraph{Hạn chế.} Mô hình nhạy cảm với siêu tham số hơn rừng ngẫu nhiên, dễ
quá khớp nếu không dừng sớm, khó diễn giải trực tiếp, và xác suất đầu ra cần được
hiệu chuẩn -- đặc biệt khi sử dụng trọng số lớp để bù mất cân bằng dữ liệu.

\subsection{So sánh và lựa chọn mô hình phù hợp}
\label{subsec:so-sanh-mo-hinh}

\subsubsection{Hệ thống chỉ tiêu đánh giá}

Việc so sánh các mô hình đòi hỏi một bộ chỉ tiêu phản ánh đầy đủ các khía
cạnh mà bài toán quan tâm. Đề án sử dụng bốn nhóm chỉ tiêu.

\paragraph{Nhóm 1 -- Chỉ tiêu phân loại phổ quát.} Xuất phát từ ma trận nhầm lẫn
với bốn ô TP, FP, TN, FN:

\begin{equation}
\text{Precision} = \frac{TP}{TP + FP}, \qquad
\text{Recall} = \frac{TP}{TP + FN},
\label{eq:prf}
\end{equation}
\begin{equation}
F_1 = \frac{2 \cdot \text{Precision} \cdot \text{Recall}}
           {\text{Precision} + \text{Recall}}.
\label{eq:f1}
\end{equation}

Độ chính xác (\eng{accuracy}) được báo cáo nhưng không dùng làm tiêu chí chính,
vì lý do đã nêu ở mục~\ref{subsec:hoc-co-giam-sat}.

\paragraph{Nhóm 2 -- Chỉ tiêu đặc thù ngành tín dụng.} Ba chỉ tiêu dưới đây được
sử dụng phổ biến trong kiểm định mô hình tín dụng~\cite{baesens2016}:

\begin{itemize}
  \item \textbf{$\AUC$} -- diện tích dưới đường cong ROC, bằng xác suất một hồ sơ
  xấu chọn ngẫu nhiên nhận điểm rủi ro cao hơn một hồ sơ tốt chọn ngẫu nhiên. Chỉ
  tiêu này không phụ thuộc ngưỡng quyết định, nên phù hợp để so sánh mô hình.

  \item \textbf{Hệ số Gini} $= 2 \cdot \AUC - 1$, là cách biểu diễn quen thuộc
  hơn trong ngành ngân hàng, với thang giá trị từ 0 (không phân biệt) đến 1 (phân
  biệt hoàn hảo).

  \item \textbf{Thống kê $\KS$} -- khoảng cách lớn nhất giữa hàm phân phối tích
  lũy của điểm số ở nhóm tốt và nhóm xấu:
  \begin{equation}
  \KS = \max_{s} \big| F_{\text{tốt}}(s) - F_{\text{xấu}}(s) \big|.
  \label{eq:ks}
  \end{equation}
  Ý nghĩa vận hành của $\KS$ là nó chỉ ra điểm cắt tại đó mô hình phân tách hai
  nhóm mạnh nhất.
\end{itemize}

\paragraph{Nhóm 3 -- Chất lượng hiệu chuẩn xác suất.} Một mô hình có thể xếp hạng
rất tốt nhưng cho giá trị xác suất sai lệch. Vì đề án sử dụng $\PD$ trực tiếp
để phân vùng quyết định, chất lượng hiệu chuẩn phải được đo tường minh bằng điểm
Brier và sai số hiệu chuẩn kỳ vọng:

\begin{equation}
\text{Brier} = \frac{1}{n}\sum_{i=1}^{n} (p_i - y_i)^2,
\qquad
\text{ECE} = \sum_{b=1}^{B} \frac{n_b}{n}
             \big| \bar{y}_b - \bar{p}_b \big|.
\label{eq:brier-ece}
\end{equation}

Hai kỹ thuật hiệu chuẩn phổ biến là hồi quy đẳng hướng (\eng{isotonic
regression})~\cite{zadrozny2002} và phương pháp Platt. Nghiên cứu của
Niculescu-Mizil và Caruana~\cite{niculescu2005} cho thấy các mô hình học tập hợp
thường cho xác suất bị nén về phía trung tâm và cần hiệu chuẩn lại.

\paragraph{Nhóm 4 -- Chỉ tiêu vận hành.} Bao gồm thời gian huấn luyện, thời gian
chấm điểm một hồ sơ và mức độ giải thích được. Đây là những chỉ tiêu quyết định
khả năng triển khai thực tế nhưng thường bị bỏ qua trong các so sánh học thuật.

\subsubsection{Cơ sở lựa chọn mô hình}

Bảng~\ref{tab:so-sanh-mo-hinh} tổng hợp đặc điểm của năm nhóm mô hình theo các
tiêu chí quan trọng đối với bài toán kiểm duyệt tín dụng.

\begin{table}[htbp]
\centering
\caption{So sánh định tính các nhóm mô hình theo tiêu chí của bài toán tín dụng}
\label{tab:so-sanh-mo-hinh}
\small
\begin{tabular}{@{}L{2.6cm}ccccc@{}}
\toprule
\textbf{Mô hình} & \textbf{Năng lực} & \textbf{Giải} & \textbf{Hiệu} &
\textbf{Chi phí} & \textbf{Độ ổn} \\
 & \textbf{dự báo} & \textbf{thích} & \textbf{chuẩn} & \textbf{huấn luyện} &
 \textbf{định} \\
\midrule
Hồi quy logistic      & Trung bình & Rất cao    & Rất tốt    & Rất thấp  & Cao \\
Thẻ điểm WOE          & Trung bình & Rất cao    & Rất tốt    & Thấp      & Cao \\
Cây quyết định        & Thấp       & Cao        & Trung bình & Rất thấp  & Thấp \\
Rừng ngẫu nhiên       & Cao        & Trung bình & Cần hiệu chuẩn & Trung bình & Cao \\
Máy vectơ hỗ trợ      & Cao        & Thấp       & Cần hiệu chuẩn & Cao       & Trung bình \\
Gradient boosting     & Rất cao    & Thấp       & Cần hiệu chuẩn & Trung bình & Trung bình \\
\bottomrule
\end{tabular}
\end{table}

Bảng trên cho thấy không tồn tại một mô hình tối ưu trên mọi tiêu chí: các mô
hình có năng lực dự báo cao nhất lại là những mô hình khó giải thích nhất và cần
hiệu chuẩn xác suất. Do vậy việc lựa chọn phải dựa trên đánh đổi phù hợp với bối
cảnh sử dụng cụ thể.

Kết quả đối sánh quy mô lớn của Lessmann và cộng sự~\cite{lessmann2015} cung cấp
căn cứ định lượng cho đánh đổi này: trên nhiều bộ dữ liệu tín dụng, khoảng cách
$\AUC$ giữa mô hình học tập hợp tốt nhất và hồi quy logistic thường chỉ ở mức
0,01--0,03. Mức chênh lệch khiêm tốn này đặt ra câu hỏi liệu việc chấp nhận mất
khả năng giải thích và tăng chi phí vận hành có tương xứng với mức cải thiện đó
hay không. Câu trả lời phụ thuộc quy mô danh mục: với danh mục lớn, mức cải thiện
0,01 $\AUC$ có thể tương ứng khoản tiết kiệm đáng kể về tổn thất tín dụng; với
danh mục nhỏ, yêu cầu về tính minh bạch và chi phí vận hành thường được ưu tiên
hơn.

Trên cơ sở phân tích đó, đề án xác định nguyên tắc lựa chọn gồm bốn tiêu chí,
theo thứ tự ưu tiên: (i) năng lực phân biệt đo bằng $\AUC$ và $\KS$ trên tập kiểm
thử ngoài thời gian; (ii) chất lượng hiệu chuẩn xác suất sau khi đã hiệu chuẩn;
(iii) thời gian chấm điểm một hồ sơ trong điều kiện xử lý trực tuyến; và (iv)
khả năng giải thích kết quả cho từng hồ sơ. Việc so sánh thực nghiệm theo bốn
tiêu chí này được trình bày ở Chương~\ref{chap:thuc-nghiem}.

\subsubsection{Xử lý mất cân bằng dữ liệu}

Ba nhóm kỹ thuật thường được dùng để xử lý mất cân bằng lớp: \textit{lấy mẫu lại}
-- tăng mẫu lớp thiểu số (chẳng hạn SMOTE~\cite{chawla2002}) hoặc giảm mẫu lớp đa
số; \textit{trọng số lớp} -- gán trọng số lớn hơn cho quan sát thuộc lớp thiểu số
trong hàm mất mát; và \textit{học nhạy cảm chi phí} -- đưa ma trận chi phí vào
quá trình huấn luyện hoặc vào việc chọn ngưỡng~\cite{elkan2001}.

Đề án sử dụng \textbf{trọng số lớp} kết hợp \textbf{lựa chọn ngưỡng theo hàm chi
phí}, không sử dụng kỹ thuật lấy mẫu lại. Ở mức mất cân bằng vừa phải (khoảng
15\%), việc lấy mẫu lại làm thay đổi phân phối tiên nghiệm của dữ liệu, khiến xác
suất đầu ra bị lệch hệ thống và phải điều chỉnh lại, trong khi lợi ích về năng
lực phân biệt không rõ ràng. Trọng số lớp giữ nguyên phân phối gốc, còn phần lệch
xác suất do trọng số gây ra được xử lý ở bước hiệu chuẩn.

\subsubsection{Giải thích mô hình}

Với mô hình sử dụng trong hoạt động cấp tín dụng, khả năng giải thích là điều
kiện bắt buộc: cán bộ tín dụng cần hiểu căn cứ của khuyến nghị để quyết định chấp
nhận hay bác bỏ, còn khách hàng có quyền được biết lý do hồ sơ bị từ chối.

Phương pháp SHAP (\eng{SHapley Additive
exPlanations})~\cite{lundberg2017,lundberg2020} dựa trên giá trị Shapley trong lý
thuyết trò chơi hợp tác. Với mỗi hồ sơ, SHAP phân rã chênh lệch giữa dự báo của
mô hình và giá trị dự báo trung bình thành tổng đóng góp của từng đặc trưng:

\begin{equation}
f(\bm{x}) = \phi_0 + \sum_{j=1}^{d} \phi_j,
\label{eq:shap}
\end{equation}

trong đó $\phi_0$ là giá trị kỳ vọng của dự báo và $\phi_j$ là đóng góp của đặc
trưng thứ $j$. Tính cộng tính này cho phép chuyển kết quả SHAP thành danh sách lý
do (\eng{reason codes}) trình bày được cho người sử dụng.

\subsubsection{Giám sát trôi dạt}

Sau khi triển khai, mô hình cần được giám sát liên tục. Chỉ số được dùng phổ biến
nhất là Population Stability Index, đo mức chênh lệch phân phối của một biến giữa
giai đoạn tham chiếu và giai đoạn hiện tại:

\begin{equation}
\PSI = \sum_{k=1}^{K}
\big( p_k^{\text{hiện tại}} - p_k^{\text{tham chiếu}} \big)
\ln \frac{p_k^{\text{hiện tại}}}{p_k^{\text{tham chiếu}}}.
\label{eq:psi}
\end{equation}

Thực hành ngành sử dụng ngưỡng: $\PSI < 0{,}10$ -- phân phối ổn định;
$0{,}10 \le \PSI < 0{,}25$ -- cần theo dõi; $\PSI \ge 0{,}25$ -- thay đổi đáng
kể, cần xem xét huấn luyện lại mô hình.

\section*{Kết luận chương 1}
\addcontentsline{toc}{section}{Kết luận chương 1}

Chương 1 đã hệ thống hóa cơ sở lý luận làm nền tảng cho đề án, với bốn kết quả
chính.

Thứ nhất, chương phân biệt rành mạch giữa kiểm duyệt (quá trình đánh giá, sàng
lọc hồ sơ) với phê duyệt (hành vi ra quyết định của cấp có thẩm quyền). Sự phân
biệt này xác định vai trò \textit{hỗ trợ} của mô hình và là cơ sở cho thiết kế
ba vùng quyết định ở chương sau. Chương cũng chỉ ra bốn hạn chế mang tính cấu
trúc của phương pháp kiểm duyệt truyền thống: không biểu diễn được tương tác giữa
các yếu tố, ngưỡng đặt theo phán đoán, kết quả nhị phân thiếu thang đo mức độ, và
giới hạn về tính nhất quán cùng tốc độ xử lý.

Thứ hai, trong phân rã $\EL = \PD \times \LGD \times \EAD$, đề án xác định phạm
vi là ước lượng thành phần $\PD$, với định nghĩa vỡ nợ thống nhất theo quy định
phân loại nợ hiện hành: quá hạn từ 90 ngày trở lên, tương ứng nợ nhóm 3--5.

Thứ ba, chương phân tích ba đặc thù của bài toán tín dụng (dữ liệu mất cân bằng,
chi phí hai loại sai lầm chênh lệch lớn, trôi dạt theo thời gian) và năm đặc điểm
của dữ liệu tín dụng, trong đó quan trọng nhất là yêu cầu về tính đúng đắn theo
thời điểm. Những đặc thù này quy định ba lựa chọn phương pháp ở các chương sau:
dùng trọng số lớp thay vì lấy mẫu lại, chọn ngưỡng theo hàm chi phí thay vì giá
trị mặc định, và chia dữ liệu theo thời gian thay vì chia ngẫu nhiên.

Thứ tư, chương trình bày nguyên lý, ưu điểm và hạn chế của năm nhóm thuật toán
được nghiên cứu cùng bốn nhóm chỉ tiêu đánh giá. Phân tích cho thấy
không tồn tại mô hình tối ưu trên mọi tiêu chí, nên việc lựa chọn phải dựa trên
đánh đổi giữa năng lực dự báo, khả năng giải thích, chất lượng hiệu chuẩn và chi
phí vận hành.

% ===========================================================================
% CHƯƠNG 2: XÂY DỰNG MÔ HÌNH HỌC MÁY TRONG HỖ TRỢ KIỂM DUYỆT TÍN DỤNG
% ===========================================================================
\chapter{XÂY DỰNG MÔ HÌNH HỌC MÁY TRONG HỖ TRỢ KIỂM DUYỆT TÍN DỤNG}
% Tiêu đề chạy rút gọn: tiêu đề chương đầy đủ dài hơn một dòng đầu trang,
% khiến fancyhdr phải nới \headheight và làm lệch khung chữ của cả chương.
\markboth{Chương 2. Xây dựng mô hình học máy hỗ trợ kiểm duyệt tín dụng}{}
\label{chap:xay-dung-he-thong}

\noindent\textit{Chương 2 trình bày thực trạng quy trình thẩm định tín dụng, xác
định bài toán học máy và phương pháp thu thập, tiền xử lý dữ liệu hồ sơ vay. Nội
dung chương tập trung vào kỹ thuật trích xuất đặc trưng, gán nhãn rủi ro và thiết
lập quy trình huấn luyện, tối ưu mô hình làm tiền đề cho các thực nghiệm trình
bày ở Chương~\ref{chap:thuc-nghiem}.}

\section{Phân tích bài toán kiểm duyệt tín dụng}
\label{sec:phan-tich-bai-toan}

\subsection{Xác định các mức độ rủi ro và trạng thái khoản vay cần dự báo}
\label{subsec:muc-do-rui-ro}

\paragraph{Biến mục tiêu.} Đối tượng dự báo của mô hình là xác suất một hồ sơ
vay phát sinh nợ xấu trong cửa sổ quan sát kể từ thời điểm cấp tín dụng. Biến mục
tiêu được định nghĩa nhị phân:

\begin{equation}
y_i =
\begin{cases}
1, & \text{nếu hồ sơ } i \text{ phát sinh nợ nhóm 3--5 (quá hạn} \ge 90 \text{ ngày)} \\
   & \text{trong vòng 12 tháng kể từ ngày cấp tín dụng,} \\[4pt]
0, & \text{trong trường hợp còn lại.}
\end{cases}
\label{eq:target}
\end{equation}

\paragraph{Ba trạng thái thay vì hai.} Mặc dù biến mục tiêu là nhị phân, quá
trình gán nhãn phải phân biệt \textbf{ba} trạng thái. Ngoài hai trạng thái
\textit{Tốt} và \textit{Xấu}, cần nhận diện trạng thái \textit{Chưa xác định}
dành cho các hồ sơ chưa đi hết cửa sổ quan sát tại thời điểm trích xuất dữ liệu.
Nếu bỏ qua và mặc định xếp nhóm này vào loại Tốt, dữ liệu huấn luyện sẽ chứa nhãn
sai một cách có hệ thống ở chính những hồ sơ mới nhất -- nhóm gần nhất với điều
kiện hiện tại. Hệ quả là mô hình đánh giá thấp rủi ro của các hồ sơ gần đây. Đề
án loại nhóm Chưa xác định khỏi tập huấn luyện và tập đánh giá.

\paragraph{Thang phân hạng rủi ro.} Kết quả đầu ra không dừng ở một giá trị xác
suất. Để phục vụ nghiệp vụ, thang $\PD$ liên tục được ánh xạ sang bảy hạng rủi ro
theo thông lệ xếp hạng tín dụng, trình bày tại Bảng~\ref{tab:risk-bands-design}.

\begin{table}[htbp]
\centering
\caption{Thang phân hạng rủi ro theo xác suất vỡ nợ}
\label{tab:risk-bands-design}
\small
\begin{tabular}{@{}cccl@{}}
\toprule
\textbf{Hạng} & \textbf{Khoảng $\PD$} & \textbf{Mức rủi ro} &
\textbf{Định hướng xử lý} \\
\midrule
AA  & $[0{,}00;\ 0{,}02)$ & Rất thấp    & Ưu tiên duyệt, có thể ưu đãi lãi suất \\
A   & $[0{,}02;\ 0{,}05)$ & Thấp        & Duyệt theo điều kiện chuẩn \\
BBB & $[0{,}05;\ 0{,}08)$ & Trung bình thấp & Duyệt, có thể yêu cầu bổ sung điều kiện \\
BB  & $[0{,}08;\ 0{,}15)$ & Trung bình  & Chuyển thẩm định thủ công \\
B   & $[0{,}15;\ 0{,}25)$ & Trung bình cao & Chuyển thẩm định thủ công \\
CCC & $[0{,}25;\ 0{,}40)$ & Cao         & Thẩm định kỹ hoặc từ chối \\
D   & $[0{,}40;\ 1{,}00]$ & Rất cao     & Từ chối \\
\bottomrule
\end{tabular}
\end{table}

\subsection{Đặc trưng đầu vào và đầu ra của mô hình kiểm duyệt tín dụng}
\label{subsec:dau-vao-dau-ra}

\paragraph{Đầu vào.} Vectơ đặc trưng của mỗi hồ sơ được tổng hợp từ bốn nhóm
nguồn dữ liệu đã phân tích ở mục~\ref{subsec:nhom-du-lieu}, bổ sung thêm nhóm thứ
năm gồm các biến tổ hợp và biến tương tác. Nguyên tắc bắt buộc đối với mọi đặc
trưng là \textbf{tính đúng đắn theo thời điểm}: giá trị của mọi biến chỉ được
tính từ dữ liệu tồn tại trước hoặc tại thời điểm nộp hồ sơ.

\paragraph{Đầu ra.} Khác với bài toán phân loại thông thường chỉ trả về một nhãn,
kết quả đầu ra của mô hình gồm năm thành phần, phản ánh vai trò hỗ trợ chứ không
thay thế con người:

\begin{enumerate}[label=(\arabic*)]
  \item \textbf{Xác suất vỡ nợ} $\PD \in (0, 1)$ đã được hiệu chuẩn;
  \item \textbf{Hạng rủi ro} theo thang bảy mức ở Bảng~\ref{tab:risk-bands-design};
  \item \textbf{Khuyến nghị} thuộc một trong ba giá trị: \code{APPROVE} (khuyến
  nghị duyệt), \code{REVIEW} (chuyển thẩm định thủ công), \code{REJECT} (khuyến
  nghị từ chối);
  \item \textbf{Danh sách lý do} (\eng{reason codes}) gồm các luật nghiệp vụ bị
  vi phạm và các chỉ tiêu định lượng có ảnh hưởng lớn nhất tới kết quả;
  \item \textbf{Đề xuất điều kiện cấp tín dụng} gồm hạn mức, kỳ hạn và mức lãi
  suất theo rủi ro, trong trường hợp khuyến nghị duyệt.
\end{enumerate}

Thành phần thứ tư (danh sách lý do) là điểm phân biệt giữa mô hình đóng vai trò
hỗ trợ và mô hình ra quyết định tự động. Thiếu thành phần này, cán bộ tín dụng
chỉ nhận được
một con số mà không có căn cứ để chấp nhận hay bác bỏ, dẫn tới hai thái cực đều
không mong muốn: chấp nhận kết quả một cách thụ động hoặc bỏ qua hoàn toàn khuyến
nghị của mô hình.

\subsection{Yêu cầu kỹ thuật và pháp lý đối với mô hình kiểm duyệt tín dụng}
\label{subsec:yeu-cau}

\subsubsection{Yêu cầu chức năng}

Quy trình kiểm duyệt phải thực hiện được tám chức năng, tương ứng các chặng của
một lượt xử lý hồ sơ. Trên \textit{đường xử lý chính}, quy trình
\textbf{tiếp nhận hồ sơ} qua
giao diện dịch vụ và kiểm tra tính đầy đủ, hợp lệ của dữ liệu đầu vào;
\textbf{tổng hợp dữ liệu} khách hàng, tài khoản, giao dịch và bản tin CIC về cấp
độ hồ sơ; \textbf{xây dựng đặc trưng} theo đúng định nghĩa đã dùng khi huấn
luyện, bảo đảm tính đúng đắn theo thời điểm; \textbf{chấm điểm rủi ro} bằng mô
hình đã hiệu chuẩn; \textbf{áp dụng bộ luật nghiệp vụ}; và \textbf{đưa ra khuyến
nghị} ba mức theo thứ tự ưu tiên xác định, kèm \textbf{danh sách lý do} và
\textbf{đề xuất điều kiện cấp tín dụng} (hạn mức, kỳ hạn, lãi suất theo rủi ro).

Ngoài đường xử lý chính, quy trình cần hỗ trợ \textbf{chấm điểm theo lô} phục vụ
tái đánh giá danh mục hiện hữu, và \textbf{giám sát mô hình} thông qua theo dõi
trôi dạt dữ liệu cùng chất lượng dự báo theo thời gian.

\subsubsection{Yêu cầu phi chức năng}

\begin{itemize}
  \item \textbf{Tính nhất quán:} cùng một hồ sơ với cùng dữ liệu đầu vào phải
  luôn cho cùng kết quả. Yêu cầu này đòi hỏi định nghĩa đặc trưng khi huấn luyện
  và khi chấm điểm phải trùng khớp tuyệt đối.

  \item \textbf{Thời gian phản hồi:} thời gian chấm điểm một hồ sơ phải đủ nhỏ để
  tích hợp được vào giao diện làm việc trực tuyến của cán bộ tín dụng. Đề án đặt
  mục tiêu thời gian xử lý một hồ sơ ở mức dưới 100 ms.

  \item \textbf{Khả năng hoạt động khi thiếu dữ liệu:} mô hình phải trả về kết
  quả cả khi không truy vấn được bản tin CIC. Trong trường hợp đó, kết quả là
  \code{REVIEW} kèm lý do, chứ không phải một thông báo lỗi.

  \item \textbf{Khả năng tái lập:} toàn bộ quy trình từ dữ liệu tới kết quả phải
  tái lập được với cùng một hạt giống ngẫu nhiên.
\end{itemize}

\subsubsection{Yêu cầu pháp lý và tuân thủ}

Mô hình sử dụng trong hoạt động cấp tín dụng chịu ràng buộc bởi các quy định pháp
lý, trong đó năm nhóm yêu cầu sau ảnh hưởng trực tiếp tới thiết kế.

\textit{Giải trình được:} mọi khuyến nghị phải kèm lý do có thể trình bày cho
khách hàng và cơ quan quản lý. Yêu cầu này loại trừ khả năng sử dụng mô hình như
một hộp đen, và là lý do đề án tích hợp cơ chế giải thích SHAP ngay trong thiết
kế chứ không coi đó là phần bổ sung.

\textit{Không phân biệt đối xử:} mô hình không được sử dụng các biến bị cấm hoặc
biến đại diện cho chúng làm căn cứ phân biệt. Trong tập đặc trưng của đề án, biến
giới tính được giữ lại phục vụ phân tích thống kê mô tả nhưng cần cân nhắc loại
bỏ khi triển khai thực tế, tùy quy định áp dụng.

\textit{Tuân thủ định nghĩa phân loại nợ:} nhãn rủi ro phải được xây dựng theo
đúng quy định hiện hành~\cite{tt11}, bảo đảm kết quả của mô hình tương thích với
các báo cáo quản trị rủi ro và trích lập dự phòng.

\textit{Bảo mật dữ liệu cá nhân:} đề án sử dụng dữ liệu mô phỏng; khi áp dụng
trên dữ liệu thực, dữ liệu cần được ẩn danh và giới hạn quyền truy
cập~\cite{tt15}.

\textit{Thẩm quyền quyết định thuộc về con người:} đây là ràng buộc bao trùm --
mô hình cung cấp thông tin hỗ trợ, quyết định cấp tín dụng cuối cùng thuộc về
người hoặc bộ phận có thẩm quyền theo phân cấp của tổ chức tín dụng.

\section{Thu thập và tổng hợp dữ liệu}
\label{sec:thu-thap-du-lieu}

\subsection{Nguồn dữ liệu tín dụng truyền thống và dữ liệu bổ sung}
\label{subsec:nguon-du-lieu}

Hệ thống sử dụng năm bảng dữ liệu, chia thành hai nhóm.

\paragraph{Dữ liệu tín dụng truyền thống.} Gồm ba nguồn đã được sử dụng từ lâu
trong thẩm định tín dụng: (i) \textit{hồ sơ khách hàng} chứa thông tin nhân khẩu
học, nghề nghiệp và thu nhập khai báo; (ii) \textit{đơn đề nghị vay} chứa tham số
khoản vay; và (iii) \textit{bản tin CIC} chứa lịch sử quan hệ tín dụng của khách
hàng với toàn hệ thống tổ chức tín dụng.

\paragraph{Dữ liệu bổ sung.} Gồm hai nguồn mà tổ chức tín dụng tự sở hữu nhưng
thường chưa được khai thác đầy đủ trong thẩm định: (iv) \textit{thông tin tài
khoản} và (v) \textit{lịch sử giao dịch}. Nhóm này có ba ưu thế: không phụ thuộc
dịch vụ bên ngoài, phản ánh tình trạng tài chính hiện tại thay vì lịch sử, và đặc
biệt có giá trị với nhóm khách hàng thiếu lịch sử tín dụng -- nhóm mà bản tin CIC
gần như không cung cấp thông tin.

Bảng~\ref{tab:data-schema} mô tả lược đồ của năm nguồn dữ liệu.

\begin{table}[htbp]
\centering
\caption{Lược đồ các bảng dữ liệu sử dụng trong đề án}
\label{tab:data-schema}
\small
\begin{tabular}{@{}L{2.4cm}L{2.7cm}L{7.0cm}@{}}
\toprule
\textbf{Bảng} & \textbf{Khóa chính} & \textbf{Nội dung} \\
\midrule
\code{customers} & \code{customer\_id} &
Nhân khẩu học, nghề nghiệp, thu nhập khai báo, nơi cư trú và thời gian quan hệ
với ngân hàng \\[2pt]
\code{accounts} & \code{account\_id} &
Loại tài khoản, ngày mở, số dư bình quân ba tháng, số dư hiện tại, trạng thái
\\[2pt]
\code{transactions} & \code{txn\_id} &
Thời điểm, chiều giao dịch, số tiền, nhóm giao dịch và kênh thực hiện \\[2pt]
\code{applications} & \code{application\_id} &
Thời điểm nộp, sản phẩm, số tiền đề nghị, kỳ hạn, lãi suất, mục đích vay, kênh
tiếp nhận, tài sản bảo đảm và khoản trả góp \\[2pt]
\code{cic\_reports} & \code{report\_id} &
Thời điểm truy vấn, điểm CIC, nhóm nợ, số ngày quá hạn, số tổ chức có quan hệ, dư
nợ, hạn mức và số lần truy vấn tín dụng \\
\bottomrule
\end{tabular}
\end{table}

\subsection{Thu thập và trích xuất dữ liệu khách hàng vay vốn}
\label{subsec:trich-xuat-du-lieu}

\paragraph{Đơn vị phân tích.} Toàn bộ dữ liệu được quy về đơn vị phân tích là
\textit{hồ sơ vay}: mỗi dòng trong ma trận đặc trưng cuối cùng tương ứng một hồ
sơ, được xác định bởi cặp (mã hồ sơ, thời điểm nộp). Việc chọn hồ sơ -- thay vì
khách hàng -- làm đơn vị phân tích là cần thiết vì một khách hàng có thể nộp
nhiều hồ sơ ở các thời điểm khác nhau, và tình trạng tài chính tại mỗi thời điểm
là khác nhau.

\paragraph{Quy tắc ghép nối theo thời điểm.} Đây là điểm kỹ thuật quan trọng nhất
của bước thu thập dữ liệu. Với mỗi hồ sơ nộp tại thời điểm $t$, việc ghép nối các
nguồn phải tuân thủ:

\begin{itemize}
  \item Thông tin khách hàng và tài khoản: lấy trạng thái tại thời điểm $t$;
  \item Bản tin CIC: chỉ sử dụng bản tin được truy vấn tại thời điểm không muộn
  hơn $t$;
  \item Giao dịch: chỉ tổng hợp các giao dịch có thời điểm phát sinh nhỏ hơn hoặc
  bằng $t$.
\end{itemize}

Vi phạm bất kỳ quy tắc nào trong ba quy tắc trên đều dẫn tới rò rỉ thông tin từ
tương lai. Đề án đưa quy tắc thứ nhất và thứ ba vào chính thuật toán tính đặc
trưng (mục~\ref{subsec:dac-trung-hanh-vi}), còn quy tắc thứ hai được kiểm tra
tường minh trong bước kiểm tra tính hợp lệ dữ liệu.

\subsection{Phân tích thống kê mô tả tập dữ liệu tín dụng}
\label{subsec:thong-ke-mo-ta}

Bước thống kê mô tả lập bảng đặc tả cho từng biến, gồm: kiểu dữ liệu, số quan sát
hợp lệ, tỷ lệ giá trị khuyết, số giá trị phân biệt, và với biến định lượng là các
phân vị 1\%, 25\%, 50\%, 75\%, 99\% cùng trung bình và độ lệch chuẩn.

Bảng mô tả phục vụ ba quyết định ở các bước sau. Tỷ lệ khuyết của từng biến quyết
định chiến lược điền khuyết và xác định biến nào cần loại bỏ do độ phủ quá thấp.
Khoảng cách giữa phân vị 99\% và giá trị lớn nhất cho biết mức độ nặng của đuôi
phân phối, làm căn cứ quyết định có cần cắt ngọn hay không. Số giá trị phân biệt
của biến định tính quyết định cách mã hóa: biến có ít hạng mục phù hợp với mã hóa
one-hot, biến có nhiều hạng mục cần cơ chế xử lý riêng.

\subsection{Phát hiện và xử lý dữ liệu ngoại lệ, dữ liệu không hợp lệ}
\label{subsec:du-lieu-khong-hop-le}

Hai khái niệm dưới đây thường bị gộp chung nhưng có bản chất khác nhau và đòi
hỏi cách xử lý khác nhau.

\paragraph{Dữ liệu không hợp lệ} là dữ liệu \textit{sai} -- vi phạm ràng buộc
nghiệp vụ hoặc ràng buộc logic. Ví dụ: tuổi khách hàng bằng 15, thu nhập âm, nhóm
nợ xấu nhất trong 12 tháng lại tốt hơn nhóm nợ hiện tại, hoặc bản tin CIC được
truy vấn sau thời điểm nộp hồ sơ. Đây là lỗi dữ liệu và phải được loại bỏ hoặc
sửa chữa.

\paragraph{Giá trị ngoại lệ} là giá trị nằm xa phần lớn quan sát nhưng \textit{có
thể đúng}. Một khách hàng có thu nhập gấp hai mươi lần trung vị hoặc một giao
dịch trị giá vài tỷ đồng là bất thường về mặt thống kê nhưng hoàn toàn có thể có
thật. Loại bỏ những quan sát này đồng nghĩa với việc loại bỏ chính nhóm khách
hàng mà tổ chức tín dụng quan tâm.

Đề án xây dựng bộ mười lăm ràng buộc nghiệp vụ để phát hiện dữ liệu không hợp lệ,
chia theo ba bảng dữ liệu chính. Bảng~\ref{tab:validity-rules} trình bày các ràng
buộc này.

\begin{table}[htbp]
\centering
\caption{Bộ ràng buộc kiểm tra tính hợp lệ của dữ liệu}
\label{tab:validity-rules}
\small
\setlength{\tabcolsep}{5pt}
\renewcommand{\arraystretch}{1.15}
\begin{tabular}{@{}L{2.5cm}L{4.8cm}L{6.6cm}@{}}
\toprule
\textbf{Bảng} & \textbf{Mã ràng buộc} & \textbf{Nội dung} \\
\midrule
\code{customers}
  & \code{CUS\_AGE\_RANGE}        & Tuổi khách hàng trong khoảng 18--75 \\
  & \code{CUS\_INCOME\_POSITIVE}  & Thu nhập khai báo lớn hơn 0 \\
  & \code{CUS\_TENURE\_NONNEG}    & Thâm niên công việc không âm \\
  & \code{CUS\_DEPENDENTS\_RANGE} & Số người phụ thuộc trong khoảng 0--10 \\
\midrule
\code{applications}
  & \code{APP\_AMOUNT\_POSITIVE}     & Số tiền đề nghị vay lớn hơn 0 \\
  & \code{APP\_TENOR\_RANGE}         & Kỳ hạn vay trong khoảng 3--360 tháng \\
  & \code{APP\_RATE\_RANGE}          & Lãi suất trong khoảng 0--40\%/năm \\
  & \code{APP\_INSTALLMENT\_POSITIVE} & Khoản trả góp hàng tháng lớn hơn 0 \\
  & \code{APP\_LTV\_RANGE}           & Tỷ lệ vay trên giá trị tài sản bảo đảm
  không vượt 1,2 \\
\midrule
\code{cic\_reports}
  & \code{CIC\_SCORE\_RANGE}       & Điểm CIC trong khoảng 300--850 \\
  & \code{CIC\_GROUP\_RANGE}       & Nhóm nợ trong khoảng 1--5 \\
  & \code{CIC\_GROUP\_ORDER}       & Nhóm nợ xấu nhất 12 tháng không tốt hơn
  nhóm nợ hiện tại \\
  & \code{CIC\_DPD\_NONNEG}        & Số ngày quá hạn không âm \\
  & \code{CIC\_UTILIZATION\_RANGE} & Tỷ lệ sử dụng hạn mức trong khoảng 0--2 \\
  & \code{CIC\_QUERY\_BEFORE\_APPLY} & Bản tin CIC được truy vấn không muộn hơn
  thời điểm nộp hồ sơ \\
\bottomrule
\end{tabular}
\end{table}

\paragraph{Kiểm chứng cơ chế phát hiện.} Một bảng kết quả kiểm tra không ghi nhận
vi phạm nào có thể tương ứng hai khả năng: dữ liệu thực sự sạch, hoặc cơ chế kiểm
tra không hoạt động. Để phân biệt, đề án thực hiện phép \textbf{tiêm lỗi}: tạo
một bản sao dữ liệu, cố ý đưa vào một tỷ lệ lỗi đã biết trước ứng với từng ràng
buộc, rồi chạy lại cơ chế kiểm tra và đối chiếu số lỗi phát hiện được với số lỗi
đã tiêm. Kết quả kiểm chứng được trình bày ở
mục~\ref{subsec:kq-chat-luong-du-lieu}.

\paragraph{Xử lý giá trị ngoại lệ.} Đối với giá trị ngoại lệ, đề án không loại bỏ
quan sát mà áp dụng kỹ thuật \textbf{cắt ngọn} (\eng{winsorize}) tại phân vị 1\%
và 99\% học từ tập huấn luyện. Cách này giữ nguyên số quan sát -- do đó không làm
mất nhóm khách hàng đặc thù -- đồng thời hạn chế ảnh hưởng quá mức của vài giá
trị cực trị lên các mô hình nhạy cảm với thang đo như hồi quy logistic, SVM và
mạng nơ-ron. Việc phát hiện ngoại lệ sử dụng quy tắc $k \times \text{IQR}$ với
$k = 3$ thay vì giá trị thông thường $k = 1{,}5$, bởi các biến tiền tệ trong dữ
liệu tín dụng có phân phối lệch phải mạnh; ngưỡng $1{,}5 \times \text{IQR}$ sẽ
đánh dấu một tỷ lệ rất lớn quan sát là ngoại lệ mà không mang lại thông tin hữu
ích.

\subsection{Phân chia dữ liệu theo khoảng thời gian và nhóm khách hàng}
\label{subsec:phan-chia-du-lieu}

Đề án sử dụng hai chiến lược phân chia với hai mục đích khác nhau.

\paragraph{Phân chia theo thời gian (chính).} Dữ liệu được sắp xếp theo thời điểm
nộp hồ sơ và chia thành ba đoạn liên tiếp: 70\% đầu làm tập huấn luyện, 15\% tiếp
theo làm tập kiểm định, 15\% cuối làm tập kiểm thử. Chiến lược này mô phỏng đúng
điều kiện vận hành thực tế -- mô hình luôn được huấn luyện trên quá khứ và áp
dụng cho tương lai -- nên phản ánh cả ảnh hưởng của trôi dạt dữ liệu.

\paragraph{Phân chia theo nhóm khách hàng (phương án bổ sung).} Toàn bộ hồ sơ của
một khách hàng được đặt vào cùng một tập. Chiến lược này loại trừ rò rỉ thông tin
theo khách hàng nhưng không phản ánh trôi dạt theo thời gian, nên chỉ phù hợp khi
tập dữ liệu có nhiều khách hàng nộp lặp nhiều hồ sơ. Hệ thống hiện thực sẵn
phương án này để có thể đối chiếu khi áp dụng trên dữ liệu thực; toàn bộ kết quả
thực nghiệm trình bày ở Chương~\ref{chap:thuc-nghiem} sử dụng cách chia theo thời
gian.

\paragraph{Tách riêng vai trò trong tập huấn luyện.} Với các mô hình có cơ chế
dừng sớm (gradient boosting), cần một tập dữ liệu độc lập để đánh giá tiêu chí
dừng. Nếu lấy tập này từ tập kiểm định, tập kiểm định sẽ đồng thời phục vụ hai
mục đích -- dừng sớm và hiệu chuẩn xác suất -- làm suy giảm tính độc lập của cả
hai. Đề án tách 15\% cuối của tập huấn luyện thành tập dừng sớm riêng biệt, cho
cấu trúc bốn phần:

\begin{center}
\begin{tabular}{@{}cccc@{}}
\toprule
\textbf{Tập} & \textbf{Tỷ trọng} & \textbf{Vai trò} & \textbf{Sử dụng} \\
\midrule
\code{train\_fit} & 59,5\% & Huấn luyện & Mọi mô hình \\
\code{train\_es}  & 10,5\% & Dừng sớm   & Mô hình boosting \\
\code{valid}      & 15,0\% & Hiệu chuẩn xác suất, chọn ngưỡng & Mọi mô hình \\
\code{test}       & 15,0\% & Đánh giá cuối cùng & Chỉ dùng một lần \\
\bottomrule
\end{tabular}
\end{center}

Nguyên tắc quan trọng nhất là tập \code{test} chỉ được sử dụng \textit{một lần
duy nhất} ở bước báo cáo kết quả cuối cùng. Mọi quyết định về lựa chọn mô hình,
siêu tham số và ngưỡng đều dựa trên tập \code{valid}.

\section{Tiền xử lý và làm sạch dữ liệu}
\label{sec:tien-xu-ly}

Nguyên tắc xuyên suốt của toàn bộ khối tiền xử lý: \textbf{mọi tham số biến đổi
chỉ được ước lượng trên tập huấn luyện} rồi áp dụng cho các tập còn lại. Nếu ước
lượng trung vị, trung bình hay phân vị trên toàn bộ dữ liệu, thông tin từ tập
kiểm thử rò rỉ vào quá trình huấn luyện và kết quả đánh giá sẽ lạc quan hơn thực
tế.

\subsection{Xử lý giá trị khuyết thiếu}
\label{subsec:xu-ly-khuyet}

Trong dữ liệu tín dụng, giá trị khuyết không đồng nghĩa với lỗi thu thập. Đề án
phân biệt hai loại và xử lý theo hai cách khác nhau.

\paragraph{Khuyết mang thông tin.} Khách hàng chưa từng có quan hệ tín dụng
(\eng{thin-file}) sẽ có toàn bộ trường CIC trống. Bản thân trạng thái ``không có
lịch sử tín dụng'' là một tín hiệu rủi ro riêng biệt. Cách xử lý đúng gồm hai
phần: điền giá trị theo \textit{quy ước nghiệp vụ} và bổ sung \textit{biến chỉ
báo khuyết} để mô hình tự học mức rủi ro thực tế của nhóm này. Quy ước điền được
chọn theo hướng trung tính -- không giả định khách hàng tốt cũng không giả định
khách hàng xấu: điểm CIC đặt bằng 600 (mức giữa của thang 300--850), nhóm nợ đặt
bằng 1 (chưa phát sinh nợ), các chỉ tiêu đếm và dư nợ đặt bằng 0.

\paragraph{Khuyết ngẫu nhiên.} Với các trường như thu nhập khai báo hay thâm niên
công việc, giá trị khuyết phản ánh việc khách hàng không cung cấp thông tin. Đề
án điền bằng \textit{trung vị của tập huấn luyện} kèm biến chỉ báo. Trung vị được
chọn thay vì trung bình vì ít nhạy cảm với đuôi phân phối lệch phải điển hình của
các biến tiền tệ.

Việc điền theo thống kê được thực hiện \textit{bên trong} pipeline huấn luyện của
scikit-learn, nhờ đó tham số điền được ước lượng lại ở mỗi lần khớp mô hình trong
kiểm định chéo, loại bỏ hoàn toàn nguy cơ rò rỉ.

\subsection{Phát hiện và xử lý dữ liệu ngoại lệ}
\label{subsec:xu-ly-ngoai-le}

Quy trình xử lý ngoại lệ gồm hai bước. Bước thứ nhất -- thống kê -- sử dụng quy
tắc $3 \times \text{IQR}$ để lập bảng các biến có giá trị nằm ngoài ngưỡng, kèm
tỷ lệ và giá trị các phân vị biên. Bảng này là cơ sở để đánh giá mức độ lệch của
từng biến.

Bước thứ hai -- xử lý -- áp dụng phép cắt ngọn tại phân vị 1\% và 99\%. Ngưỡng
cắt được học từ tập huấn luyện và áp dụng nguyên vẹn cho tập kiểm định và tập
kiểm thử; nếu học lại trên từng tập, phép biến đổi sẽ khác nhau giữa các tập và
làm sai lệch kết quả đánh giá.

Phép cắt ngọn chỉ cần thiết với nhóm mô hình nhạy cảm với thang đo. Các mô hình dựa trên cây chia theo ngưỡng nên bất biến với phép biến
đổi đơn điệu, do đó không chịu ảnh hưởng của giá trị cực trị theo cách tương tự.

\subsection{Chuẩn hóa và mã hóa biến định tính/định lượng}
\label{subsec:chuan-hoa-ma-hoa}

Mỗi họ mô hình có yêu cầu khác nhau về dạng đầu vào, nên đề án sử dụng bốn phương
án biến đổi song song thay vì một phương án chung
(Bảng~\ref{tab:encoding-schemes}).

\begin{table}[htbp]
\centering
\caption{Bốn phương án biến đổi dữ liệu theo yêu cầu của từng họ mô hình}
\label{tab:encoding-schemes}
\small
\begin{tabular}{@{}L{2.3cm}L{4.4cm}L{5.5cm}@{}}
\toprule
\textbf{Phương án} & \textbf{Nội dung biến đổi} & \textbf{Mô hình áp dụng và lý
do} \\
\midrule
\code{numeric\_ohe} &
Điền khuyết bằng trung vị, chuẩn hóa z-score cho biến định lượng; mã hóa one-hot
cho biến định tính &
Hồi quy logistic, SVM, mạng nơ-ron. Các mô hình dựa trên khoảng cách hoặc
gradient nhạy cảm với chênh lệch thang đo giữa các biến \\[2pt]
\code{ordinal} &
Điền khuyết; mã hóa thứ tự cho biến định tính &
Cây quyết định, rừng ngẫu nhiên, Extra Trees, Hist Gradient Boosting. Cây chia
theo ngưỡng nên bất biến với phép biến đổi đơn điệu; chuẩn hóa không cần thiết và
làm mất khả năng diễn giải ngưỡng \\[2pt]
\code{native\_cat} &
Giữ nguyên biến định tính ở dạng hạng mục &
XGBoost, LightGBM, CatBoost. Các thư viện này có cơ chế xử lý biến định tính
riêng, hiệu quả hơn mã hóa bên ngoài \\[2pt]
\code{woe} &
Rời rạc hóa và thay mỗi khoảng bằng giá trị $\WOE$ &
Mô hình thẻ điểm. Đưa mọi biến về thang log-odds và cho phép áp ràng buộc đơn
điệu \\
\bottomrule
\end{tabular}
\end{table}

Toàn bộ phép biến đổi được gói trong cấu trúc \code{ColumnTransformer} của
scikit-learn, bảo đảm tham số chỉ được ước lượng trên dữ liệu huấn luyện của từng
lần khớp mô hình.

\paragraph{Xử lý hạng mục chưa từng xuất hiện.} Khi chấm điểm một hồ sơ mới, biến
định tính có thể nhận giá trị chưa từng xuất hiện trong tập huấn luyện. Đề án
``đóng băng'' danh mục hạng mục theo tập huấn luyện và lưu kèm mô hình; hạng mục
lạ được quy về giá trị khuyết. Đây là hành vi đúng khi triển khai: mô hình không
có cơ sở nào để đánh giá một hạng mục chưa từng quan sát.

\subsection{Phân chia tập dữ liệu huấn luyện, kiểm thử}
\label{subsec:chia-tap}

Cấu trúc bốn tập dữ liệu đã được trình bày ở mục~\ref{subsec:phan-chia-du-lieu}.
Nội dung cần làm rõ ở đây là \textbf{thứ tự thao tác} giữa việc phân chia và các
phép tiền xử lý:

\begin{center}
gán nhãn $\rightarrow$ loại hồ sơ chưa đủ cửa sổ quan sát $\rightarrow$ chia tập
$\rightarrow$ học tham số biến đổi trên tập huấn luyện $\rightarrow$ áp dụng cho
các tập còn lại.
\end{center}

Việc phân chia phải diễn ra \textit{trước} mọi phép biến đổi có ước lượng tham
số. Nếu đảo ngược, chẳng hạn chuẩn hóa trên toàn bộ dữ liệu rồi mới chia tập,
trung bình và độ lệch chuẩn dùng để chuẩn hóa sẽ chứa thông tin của tập kiểm thử;
mô hình khi đó được đánh giá trên dữ liệu mà nó đã gián tiếp tiếp xúc, cho kết
quả lạc quan hơn thực tế.

\section{Trích xuất và lựa chọn đặc trưng}
\label{sec:trich-xuat-dac-trung}

\subsection{Trích xuất nhóm đặc trưng thông tin cá nhân và tài chính}
\label{subsec:dac-trung-ca-nhan}

Nhóm này gồm mười lăm đặc trưng lấy trực tiếp từ hồ sơ khách hàng và thông tin
tài khoản, chia thành ba mảng: \textit{nhân khẩu học} (tuổi, giới tính, tình
trạng hôn nhân, số người phụ thuộc, trình độ học vấn, tỉnh/thành cư trú, hình
thức nhà ở); \textit{nghề nghiệp và thu nhập} (nghề nghiệp, loại hình tổ chức làm
việc, thâm niên công việc, thu nhập khai báo); và \textit{quan hệ với ngân hàng}
(thời gian quan hệ, việc có tài khoản nhận lương, số tài khoản đang mở, tổng số
dư bình quân ba tháng). Danh mục đầy đủ kèm mô tả từng đặc trưng được liệt kê tại
Phụ lục~\ref{app:features}.

Mảng thứ ba phản ánh \textit{mức độ gắn bó} của khách hàng với tổ chức tín dụng,
yếu tố có giá trị dự báo nhưng không xuất hiện trong bộ quy tắc thẩm định truyền
thống.

\subsection{Trích xuất nhóm đặc trưng lịch sử tín dụng và quan hệ CIC}
\label{subsec:dac-trung-cic}

Nhóm này gồm mười sáu đặc trưng lấy từ bản tin CIC, phản ánh bốn khía cạnh của
quan hệ tín dụng: \textit{mức độ tín nhiệm} (điểm CIC), \textit{lịch sử trả nợ}
(nhóm nợ hiện tại, nhóm nợ xấu nhất 12 tháng, số ngày quá hạn tối đa, nợ đã xử lý
rủi ro, số tổ chức từng phát sinh nợ xấu), \textit{mức độ và cơ cấu dư nợ} (số tổ
chức đang có quan hệ, số hợp đồng còn hiệu lực và đã tất toán, tổng dư nợ, tỷ lệ
sử dụng hạn mức, nghĩa vụ trả nợ hàng tháng, độ dài lịch sử tín dụng), và
\textit{mức độ tìm kiếm tín dụng} (số lần truy vấn trong 3 tháng).

Bổ sung vào nhóm là hai biến chỉ báo tạo ra ở bước xử lý khuyết, phản ánh hai
tình huống khác nhau về bản chất: khách hàng \textit{không có} lịch sử tín dụng
và trường hợp \textit{không truy vấn được} dữ liệu.

\subsection{Trích xuất nhóm đặc trưng hành vi và biến tương tác}
\label{subsec:dac-trung-hanh-vi}

\subsubsection{Đặc trưng hành vi theo cửa sổ thời gian}

Từ lịch sử giao dịch, đề án tính hai mươi ba độ đo trên mỗi cửa sổ thời gian
$W \in \{7, 30, 90, 180\}$ ngày, thuộc bốn mảng: \textit{quy mô hoạt động} (số
giao dịch, dòng tiền vào, dòng tiền ra, dòng tiền thuần, giá trị giao dịch trung
bình); \textit{mức biến động} (độ lệch chuẩn và hệ số biến thiên của giá trị giao
dịch); \textit{cơ cấu dòng tiền theo mục đích} (nhận lương, trả nợ, chi tiêu nhóm
rủi ro, rút tiền mặt, thương mại điện tử, hóa đơn tiện ích); và \textit{dấu hiệu
hành vi} (giao dịch ban đêm, giao dịch giá trị lớn, giao dịch qua kênh di động),
cùng năm tỷ số phái sinh từ các độ đo trên.

Bổ sung thêm bảy đặc trưng không phụ thuộc cửa sổ: số ngày kể từ giao dịch gần
nhất, ba biến xu thế so sánh cửa sổ ngắn với cửa sổ dài, tỷ lệ dòng tiền vào thực
tế so với thu nhập khai báo, mức độ đều đặn của việc nhận lương và độ ổn định thu
nhập.

Trong số này, biến \textit{tỷ lệ dòng tiền vào thực tế so với thu nhập khai báo}
có ý nghĩa nghiệp vụ riêng: nó cung cấp một phép kiểm chứng độc lập cho thông tin
thu nhập khách hàng tự khai, công việc mà quy trình thẩm định truyền thống phải
thực hiện thủ công qua đối chiếu chứng từ.

\subsubsection{Thuật toán bảo đảm tính đúng đắn theo thời điểm}
\label{subsubsec:point-in-time}

Yêu cầu tính đúng đắn theo thời điểm đặt ra một thách thức về tính toán. Cách làm
trực tiếp, tức với mỗi hồ sơ lọc toàn bộ giao dịch của khách hàng trong cửa sổ rồi
tổng hợp -- có độ phức tạp $O(N \times |W| \times \bar{m})$ với $N$ là số hồ sơ,
$|W|$ là số cửa sổ và $\bar{m}$ là số giao dịch trung bình mỗi khách hàng. Với
hàng chục nghìn hồ sơ và hàng triệu giao dịch, chi phí này không chấp nhận được.

Đề án sử dụng kỹ thuật \textbf{luỹ kế theo nhóm kết hợp ghép gần nhất theo thời
gian}. Với mỗi độ đo $m$, gọi $C_m(i, t)$ là tổng luỹ kế của độ đo đó cho khách
hàng $i$ tính đến thời điểm $t$. Khi đó tổng của độ đo trong cửa sổ $(t - W, t]$
được tính bằng hiệu hai giá trị luỹ kế:

\begin{equation}
S_m(i, t, W) = C_m(i, t) - C_m(i, t - W).
\label{eq:cumsum-window}
\end{equation}

Hai giá trị luỹ kế được truy xuất bằng phép ghép gần nhất hướng lùi
(\code{merge\_asof}), tương đương thao tác ``tìm bản ghi gần nhất không vượt quá
mốc thời gian'' với độ phức tạp tuyến tính sau khi dữ liệu đã được sắp xếp. Nhờ
đó, toàn bộ đặc trưng cửa sổ cho $N$ hồ sơ được tính bằng $2 \times |W|$ phép
ghép vectơ hóa thay vì $N \times |W|$ truy vấn riêng lẻ.

Phương pháp này còn có một ưu điểm về mặt đúng đắn: vì công thức chỉ sử dụng giá
trị luỹ kế tại hai mốc thời gian, nó \textit{không thể} vô tình bao gồm giao dịch
phát sinh sau thời điểm nộp hồ sơ. Tính đúng đắn theo thời điểm được bảo đảm bởi
chính cấu trúc thuật toán chứ không phụ thuộc vào việc kiểm tra bổ sung.

\subsubsection{Biến tổ hợp và biến tương tác}

Mười ba biến phái sinh được xây dựng nhằm biểu diễn các chỉ tiêu nghiệp vụ then
chốt mà không biến gốc nào thể hiện được:

\begin{table}[htbp]
\centering
\caption{Các biến tổ hợp và biến tương tác}
\label{tab:derived-features}
\small
\setlength{\tabcolsep}{4pt}
\begin{tabular}{@{}L{3.3cm}C{5.3cm}L{6.1cm}@{}}
\toprule
\textbf{Biến} & \textbf{Công thức} & \textbf{Ý nghĩa nghiệp vụ} \\
\midrule
\code{dti} &
$\dfrac{\text{Trả góp mới} + \text{Nợ hiện hữu}}{\text{Thu nhập}}$ &
Tổng gánh nặng nợ trên thu nhập, chỉ tiêu trung tâm của đánh giá khả năng trả nợ
\\[6pt]
\code{pti} &
$\dfrac{\text{trả góp khoản vay mới}}{\text{thu nhập}}$ &
Gánh nặng riêng của khoản vay đang xét \\[6pt]
\code{disposable\_income} &
thu nhập $-$ trả nợ $-$ chi phí sinh hoạt &
Thu nhập còn lại thực tế sau khi trừ nghĩa vụ và chi phí tối thiểu \\[6pt]
\code{credit\_history\_ratio} &
$\dfrac{\text{độ dài lịch sử tín dụng}}{\text{số tháng từ tuổi 18}}$ &
Mức độ tham gia thị trường tín dụng so với độ tuổi \\[6pt]
\code{dpd\_x\_institutions} &
số ngày quá hạn $\times$ số tổ chức &
Biến tương tác: quá hạn tại nhiều tổ chức nghiêm trọng hơn quá hạn tại một tổ
chức \\[6pt]
\code{high\_dti\_short\_tenure} &
$\mathbb{1}[\DTI > 0{,}5 \wedge \text{thâm niên} < 12]$ &
Cờ rủi ro kết hợp: gánh nặng nợ cao đi kèm công việc chưa ổn định \\[6pt]
\code{income\_verification\_ratio} &
$\dfrac{\text{dòng tiền vào 30 ngày}}{\text{thu nhập khai báo}}$ &
Kiểm chứng độc lập thông tin thu nhập tự khai \\
\bottomrule
\end{tabular}
\end{table}

Biến \code{disposable\_income} sử dụng mô hình chi phí sinh hoạt gồm hai thành
phần: một mức sàn cố định phụ thuộc số người phụ thuộc, và một phần tỷ lệ thuận
với thu nhập. Cách mô hình hóa này khắc phục nhược điểm của việc dùng một mức cố
định duy nhất, vốn quá khắt khe với nhóm thu nhập thấp và quá nới lỏng với nhóm
thu nhập cao.

\subsection{Đánh giá mức độ quan trọng của các đặc trưng}
\label{subsec:danh-gia-dac-trung}

Đề án sử dụng ba phương pháp bổ trợ nhau để đánh giá tầm quan trọng của đặc
trưng, mỗi phương pháp trả lời một câu hỏi khác nhau.

\paragraph{Information Value -- năng lực dự báo đơn biến.} Trả lời câu hỏi: mỗi
đặc trưng, xét riêng lẻ, phân biệt hồ sơ tốt và xấu tốt đến mức nào? Ưu điểm là
chi phí tính toán thấp và không phụ thuộc mô hình. Hạn chế là không tính đến
tương tác: một biến có $\IV$ thấp vẫn có thể có giá trị khi kết hợp với biến
khác.

\paragraph{SHAP -- đóng góp trong mô hình đã huấn luyện.} Trả lời câu hỏi: trong
mô hình cụ thể đã được huấn luyện, đặc trưng nào thực sự chi phối kết quả dự báo?
Phương pháp này tính đến tương tác và cho kết quả ở cả cấp độ toàn cục lẫn cấp độ
từng hồ sơ.

\paragraph{Phép hoán vị -- kiểm chứng chéo.} Trả lời câu hỏi: nếu phá vỡ quan hệ
giữa một đặc trưng và nhãn bằng cách xáo trộn ngẫu nhiên giá trị của nó, chất
lượng dự báo suy giảm bao nhiêu? Phương pháp này độc lập với cấu trúc nội bộ của
mô hình, nên dùng làm phép kiểm chứng chéo cho kết quả SHAP.

\subsection{Lựa chọn tập đặc trưng tối ưu cho mô hình}
\label{subsec:lua-chon-dac-trung}

Tập đặc trưng ban đầu gồm \NFeaturesInitial{} biến, trong đó nhiều biến được sinh
theo cùng một mẫu trên bốn cửa sổ thời gian khác nhau nên tương quan mạnh với
nhau. Đề án áp dụng quy trình lọc ba bước nối tiếp, sắp xếp từ rẻ đến đắt về chi
phí tính toán (Hình~\ref{fig:feature-selection}).

\begin{figure}[htbp]
\centering
\begin{tikzpicture}[
  node distance=0.62cm,
  fstep/.style={draw, rounded corners=2pt, align=center, minimum height=1.0cm,
               text width=3.05cm, font=\small, fill=blue!4},
  arrow/.style={-{Stealth[length=2mm]}, thick, draw=black!65},
  note/.style={font=\footnotesize\itshape, align=center, text width=3.05cm},
]
\node[fstep] (s0) {Tập đặc trưng\\ban đầu};
\node[fstep, right=of s0] (s1) {Bước 1\\Lọc theo $\IV$};
\node[fstep, right=of s1] (s2) {Bước 2\\Lọc theo tương quan};
\node[fstep, right=of s2] (s3) {Bước 3\\Lọc theo $\PSI$};

\draw[arrow] (s0) -- (s1);
\draw[arrow] (s1) -- (s2);
\draw[arrow] (s2) -- (s3);

\node[note, below=0.32cm of s1] {loại biến\\không phân biệt được};
\node[note, below=0.32cm of s2] {loại biến\\dư thừa};
\node[note, below=0.32cm of s3] {loại biến\\không ổn định};
\end{tikzpicture}
\caption{Quy trình lựa chọn tập đặc trưng qua ba bộ lọc nối tiếp}
\label{fig:feature-selection}
\end{figure}

\paragraph{Bước 1 -- lọc theo năng lực dự báo.} Loại các đặc trưng có
$\IV < 0{,}02$, ngưỡng mà theo thực hành ngành tương ứng với ``không có khả năng
phân biệt''. Ngưỡng trên được đặt ở mức $\IV = 2{,}0$ để phát hiện trường hợp rò
rỉ nhãn. Lưu ý rằng nhóm đặc trưng CIC vốn có $\IV$ cao (khoảng 0,9--1,3) do phản
ánh trực tiếp lịch sử trả nợ; các biến này đã được kiểm tra tính đúng đắn theo
thời điểm nên $\IV$ cao là bản chất của dữ liệu chứ không phải dấu hiệu rò rỉ.

\paragraph{Bước 2 -- lọc theo tương quan.} Trong mỗi cặp biến có hệ số tương quan
tuyệt đối vượt 0,95, giữ lại biến có $\IV$ cao hơn. Bước này chủ yếu loại bỏ các
biến trùng lặp giữa các cửa sổ thời gian liền kề.

\paragraph{Bước 3 -- lọc theo độ ổn định.} Loại các đặc trưng có $\PSI > 0{,}25$
giữa giai đoạn huấn luyện và giai đoạn kiểm thử. Mô hình phụ thuộc vào biến không
ổn định sẽ suy giảm nhanh khi vận hành.

\section{Xây dựng nhãn và xác định quy tắc phê duyệt}
\label{sec:xay-dung-nhan}

\subsection{Định nghĩa trạng thái rủi ro và tiêu chuẩn nợ xấu}
\label{subsec:dinh-nghia-nhan}

\paragraph{Tiêu chuẩn nợ xấu.} Đề án sử dụng định nghĩa theo quy định phân loại
nợ hiện hành~\cite{tt11}: khoản vay bị coi là nợ xấu khi được phân vào nhóm 3, 4
hoặc 5, tương ứng với số ngày quá hạn từ 90 ngày trở lên. Định nghĩa này trùng
với ngưỡng mà chuẩn mực Basel~\cite{bcbs2006} sử dụng cho khái niệm vỡ nợ, nhờ đó
kết quả nghiên cứu có thể đối chiếu với các công trình quốc tế.

\paragraph{Cửa sổ quan sát.} Một hồ sơ chỉ được gán nhãn sau khi đã trải qua cửa
sổ quan sát hiệu quả trả nợ đủ dài. Đề án chọn cửa sổ 12 tháng dựa trên hai căn
cứ đối lập nhau cần cân bằng: cửa sổ quá ngắn bỏ sót phần lớn trường hợp vỡ nợ
muộn, làm tỷ lệ nợ xấu quan sát được thấp hơn thực tế và mô hình lệch về phía lạc
quan; cửa sổ quá dài buộc phải loại nhiều hồ sơ gần đây khỏi tập huấn luyện, làm
dữ liệu cũ đi và mô hình kém phản ánh điều kiện hiện tại.

Căn cứ định lượng cho lựa chọn này là \textbf{phân tích vintage}: đo tỷ lệ nợ xấu
tích luỹ theo từng tháng kể từ thời điểm cấp tín dụng. Điểm mà đường cong tích
luỹ bắt đầu đi ngang chính là độ dài cửa sổ hợp lý. Kết quả phân tích được trình
bày ở mục~\ref{subsec:kq-gan-nhan}.

\subsection{Phương pháp gán nhãn dữ liệu hồ sơ vay}
\label{subsec:phuong-phap-gan-nhan}

Quy trình gán nhãn gồm ba bước:

\begin{enumerate}[label=(\arabic*)]
  \item \textbf{Xác định thời điểm kết thúc cửa sổ quan sát} cho từng hồ sơ, bằng
  thời điểm cấp tín dụng cộng với độ dài cửa sổ.

  \item \textbf{Phân loại trạng thái.} Hồ sơ có thời điểm kết thúc cửa sổ vượt
  quá mốc trích xuất dữ liệu được gán \code{INDETERMINATE}. Trong số các hồ sơ
  còn lại, hồ sơ phát sinh nợ xấu \textit{trong} cửa sổ được gán \code{BAD}, phần
  còn lại được gán \code{GOOD}. Trường hợp vỡ nợ muộn hơn cửa sổ không được tính
  là \code{BAD}, vì nằm ngoài phạm vi dự báo của mô hình.

  \item \textbf{Loại nhóm chưa xác định} khỏi tập dữ liệu mô hình hóa.
\end{enumerate}

\subsection{Xác định điểm cắt phê duyệt}
\label{subsec:diem-cat}

\subsubsection{Vì sao ngưỡng 0,5 không phù hợp}

Ngưỡng mặc định 0,5 của bài toán phân loại nhị phân không phù hợp với bài toán
tín dụng vì hai lý do. Thứ nhất, dữ liệu mất cân bằng: với tỷ lệ nợ xấu khoảng
15\%, rất ít hồ sơ có $\PD$ vượt 0,5, nên dùng ngưỡng này gần như duyệt toàn bộ.
Thứ hai, chi phí hai loại sai lầm rất khác nhau: bỏ sót một hồ sơ xấu gây tổn
thất bằng một phần đáng kể dư nợ, trong khi từ chối sai một hồ sơ tốt chỉ làm mất
phần lãi biên.

\subsubsection{Ba tiêu chí xác định điểm cắt}

Đề án hiện thực ba phương pháp xác định điểm cắt, tương ứng ba cách đặt bài toán
khác nhau trong thực tế nghiệp vụ.

\paragraph{Tiêu chí 1 -- tối thiểu chi phí kỳ vọng.} Chọn ngưỡng $\tau$ làm cực
tiểu chi phí kỳ vọng:

\begin{equation}
\tau^* = \argmin_{\tau}
\left[ c_{FN} \cdot \text{FN}(\tau) + c_{FP} \cdot \text{FP}(\tau) \right] / n,
\label{eq:cost-threshold}
\end{equation}

với $c_{FN} : c_{FP} = 12 : 1$. Áp dụng khi tổ chức tín dụng ước lượng được chi
phí của hai loại sai lầm.

\paragraph{Tiêu chí 2 -- đạt tỷ lệ duyệt mục tiêu.} Chọn ngưỡng bằng phân vị
tương ứng với tỷ lệ duyệt mong muốn. Áp dụng khi ngưỡng bị ràng buộc bởi kế hoạch
tăng trưởng dư nợ.

\paragraph{Tiêu chí 3 -- bảo đảm trần nợ xấu.} Chọn ngưỡng cao nhất mà tỷ lệ nợ
xấu trong nhóm được duyệt vẫn dưới mức trần cho phép. Chọn ngưỡng \textit{cao
nhất} thỏa mãn ràng buộc nhằm tối đa hóa quy mô cấp tín dụng trong giới hạn khẩu
vị rủi ro.

\subsubsection{Hai ngưỡng cho ba vùng quyết định}

Hệ thống hỗ trợ kiểm duyệt cần \textbf{hai} ngưỡng chứ không phải một, chia thang
$\PD$ thành ba vùng:

\begin{equation}
\text{Khuyến nghị} =
\begin{cases}
\texttt{APPROVE}, & \PD \le \tau_{\text{duyệt}}, \\
\texttt{REVIEW},  & \tau_{\text{duyệt}} < \PD \le \tau_{\text{từ chối}}, \\
\texttt{REJECT},  & \PD > \tau_{\text{từ chối}}.
\end{cases}
\label{eq:three-zones}
\end{equation}

Hai ngưỡng được đặt đối xứng quanh điểm cắt tối ưu chi phí, sao cho tỷ lệ hồ sơ
rơi vào vùng chuyển thẩm định xấp xỉ một mức mục tiêu định trước. Tỷ lệ này là
\textit{tham số chính sách} chứ không phải tham số kỹ thuật: tăng lên thì mức độ
an toàn cao hơn nhưng khối lượng thẩm định thủ công lớn hơn; giảm xuống thì tự
động hóa cao hơn nhưng rủi ro khuyến nghị sai tăng. Đề án đặt mục tiêu vùng thẩm
định ở mức 20\%.

\section{Thiết kế mô hình học máy kiểm duyệt tín dụng}
\label{sec:thiet-ke-mo-hinh}

\subsection{Lựa chọn các thuật toán học máy dự kiến}
\label{subsec:lua-chon-thuat-toan}

Đề án lựa chọn mười hai mô hình để so sánh, trong đó năm thuật toán là trọng tâm
nghiên cứu -- hồi quy logistic, cây quyết định, rừng ngẫu nhiên, máy vectơ hỗ trợ
và gradient boosting -- và bảy mô hình còn lại là các biến thể hiện đại của cùng
những họ thuật toán đó, đưa vào nhằm đánh giá xem mức cải thiện thực tế có tương
xứng với chi phí triển khai hay không
(Bảng~\ref{tab:model-selection}).

\begin{table}[htbp]
\centering
\caption{Danh mục mô hình được thiết kế và so sánh}
\label{tab:model-selection}
\small
\begin{tabular}{@{}clL{3.0cm}L{5.2cm}@{}}
\toprule
\textbf{TT} & \textbf{Mô hình} & \textbf{Nhóm} & \textbf{Vai trò trong so sánh} \\
\midrule
1  & Hồi quy logistic      & Tuyến tính & Mô hình cơ sở tham chiếu \\
2  & Thẻ điểm WOE          & Tuyến tính & Chuẩn mực thực tiễn ngân hàng \\
3  & Cây quyết định (CART) & Cây đơn    & Mô hình phi tuyến đơn giản nhất \\
4  & Rừng ngẫu nhiên       & Bagging    & Học tập hợp song song \\
5  & Extra Trees           & Bagging    & Biến thể tăng mức ngẫu nhiên hóa \\
6  & Máy vectơ hỗ trợ      & SVM        & Tiếp cận theo lề cực đại \\
7  & Hist Gradient Boosting& Boosting   & Học tập hợp tuần tự \\
8  & XGBoost               & Boosting   & Biến thể có điều chuẩn tường minh \\
9  & LightGBM              & Boosting   & Biến thể sinh cây theo chiều lá \\
10 & CatBoost              & Boosting   & Biến thể xử lý biến định tính \\
11 & Mạng nơ-ron (MLP)     & Mạng nơ-ron& Mô hình phi tuyến toàn cục \\
12 & Kết hợp xếp tầng      & Ensemble   & Tổ hợp các mô hình cơ sở khác biệt \\
\bottomrule
\end{tabular}
\end{table}

\paragraph{Ghi chú về hiện thực SVM.} Như đã phân tích ở mục~\ref{subsec:svm},
SVM hạt nhân có chi phí huấn luyện từ $O(n^2)$ đến $O(n^3)$, không khả thi với
tập huấn luyện hàng chục nghìn hồ sơ. Đề án sử dụng phép xấp xỉ Nystroem với 600
điểm mốc, sau đó huấn luyện SVM tuyến tính trên không gian đặc trưng xấp xỉ. Xác
suất đầu ra được quy đổi từ giá trị hàm quyết định bằng phương pháp Platt qua
kiểm định chéo ba tầng.

\subsection{Thiết kế kiến trúc tổng thể của mô hình học máy hỗ trợ ra quyết định}
\label{subsec:kien-truc-tong-the}

\subsubsection{Kiến trúc phân tầng}

Hệ thống được tổ chức thành bốn tầng (Hình~\ref{fig:architecture}), tách bạch
giữa phần \textit{xây dựng mô hình} (thực hiện định kỳ, ngoại tuyến) và phần
\textit{sử dụng mô hình} (thực hiện mỗi khi có hồ sơ mới).

\begin{figure}[htbp]
\centering
\begin{tikzpicture}[
  node distance=0.5cm and 0.5cm,
  layer/.style={draw, rounded corners=3pt, align=center, font=\small\bfseries,
                minimum height=0.8cm, text width=2.0cm},
  comp/.style={draw, rounded corners=2pt, align=center, font=\footnotesize,
               minimum height=0.85cm, text width=2.4cm, fill=blue!4},
  arrow/.style={-{Stealth[length=2mm]}, thick, draw=black!65},
]
% Tầng 1: dữ liệu
\node[layer, fill=gray!12] (l1) {Tầng\\dữ liệu};
\node[comp, right=of l1] (d1) {Hồ sơ khách hàng\\và tài khoản};
\node[comp, right=of d1] (d2) {Đơn đề nghị vay};
\node[comp, right=of d2] (d3) {Bản tin CIC};
\node[comp, right=of d3] (d4) {Lịch sử giao dịch};

% Tầng 2: xử lý
\node[layer, fill=gray!12, below=1.0cm of l1] (l2) {Tầng\\xử lý dữ liệu};
\node[comp, right=of l2, fill=orange!8] (p1) {Kiểm tra\\chất lượng};
\node[comp, right=of p1, fill=orange!8] (p2) {Tiền xử lý\\và làm sạch};
\node[comp, right=of p2, fill=orange!8] (p3) {Trích xuất\\đặc trưng};
\node[comp, right=of p3, fill=orange!8] (p4) {Lựa chọn\\đặc trưng};

% Tầng 3: mô hình
\node[layer, fill=gray!12, below=1.0cm of l2] (l3) {Tầng\\mô hình};
\node[comp, right=of l3, fill=green!8] (m1) {Huấn luyện\\và tối ưu};
\node[comp, right=of m1, fill=green!8] (m2) {Hiệu chuẩn\\xác suất};
\node[comp, right=of m2, fill=green!8] (m3) {Giải thích\\(SHAP)};
\node[comp, right=of m3, fill=green!8] (m4) {Giám sát\\trôi dạt};

% Tầng 4: quyết định
\node[layer, fill=gray!12, below=1.0cm of l3] (l4) {Tầng hỗ trợ\\quyết định};
\node[comp, right=of l4, fill=red!6] (e1) {Luật\\nghiệp vụ};
\node[comp, right=of e1, fill=red!6] (e2) {Chấm điểm\\$\PD$};
\node[comp, right=of e2, fill=red!6] (e3) {Chính sách\\hạn mức};
\node[comp, right=of e3, fill=red!6] (e4) {Khuyến nghị\\và lý do};

\draw[arrow] (d2.south) -- ++(0,-0.32) -| (p1.north);
\draw[arrow] (p1) -- (p2);
\draw[arrow] (p2) -- (p3);
\draw[arrow] (p3) -- (p4);
\draw[arrow] (p3.south) -- ++(0,-0.32) -| (m1.north);
\draw[arrow] (m1) -- (m2);
\draw[arrow] (m2.south) -- ++(0,-0.32) -| (e2.north);
\draw[arrow] (e1) -- (e2);
\draw[arrow] (e2) -- (e3);
\draw[arrow] (e3) -- (e4);
\end{tikzpicture}
\caption{Kiến trúc phân tầng của mô hình hỗ trợ kiểm duyệt tín dụng}
\label{fig:architecture}
\end{figure}

\subsubsection{Logic tổng hợp khuyến nghị}

Khuyến nghị cuối cùng là kết quả của một chuỗi ưu tiên rõ ràng, trong đó luật
nghiệp vụ có quyền phủ quyết cao hơn mô hình thống kê:

\begin{enumerate}[label=\arabic*.]
  \item Có luật loại trừ cứng bị vi phạm $\Rightarrow$ \code{REJECT}
  \item Mô hình gặp lỗi khi chấm điểm $\Rightarrow$ \code{REVIEW}
  \item Không truy vấn được dữ liệu CIC $\Rightarrow$ \code{REVIEW}
  \item $\PD > \tau_{\text{từ chối}}$ $\Rightarrow$ \code{REJECT}
  \item Có luật yêu cầu thẩm định bị vi phạm $\Rightarrow$ \code{REVIEW}
  \item $\PD > \tau_{\text{duyệt}}$ $\Rightarrow$ \code{REVIEW}
  \item Hạn mức tính được bằng 0 $\Rightarrow$ \code{REJECT}
  \item Các trường hợp còn lại $\Rightarrow$ \code{APPROVE}
\end{enumerate}

Thứ tự này phản ánh một nguyên tắc thiết kế: \textit{ràng buộc tuân thủ đứng
trước đánh giá thống kê}. Một hồ sơ vi phạm điều kiện pháp lý hoặc chính sách
phải bị từ chối bất kể mô hình cho điểm rủi ro thấp đến đâu. Ngược lại, khi hệ
thống không đủ thông tin hoặc không đủ chắc chắn, kết quả mặc định là chuyển cho
con người xem xét chứ không phải tự quyết.

\subsubsection{Bộ luật nghiệp vụ}

Đề án xây dựng bộ luật gồm hai mươi mốt luật, chia theo sáu nhóm và ba mức hành
động (Bảng~\ref{tab:rule-groups}).

\begin{table}[htbp]
\centering
\caption{Cơ cấu bộ luật nghiệp vụ theo nhóm và mức hành động}
\label{tab:rule-groups}
\small
\begin{tabular}{@{}L{3.2cm}L{4.2cm}L{5.0cm}@{}}
\toprule
\textbf{Nhóm luật} & \textbf{Nội dung kiểm tra} & \textbf{Ví dụ} \\
\midrule
Điều kiện vay & Độ tuổi, thu nhập tối thiểu, thâm niên công việc &
Tuổi ngoài khoảng 20--62 $\Rightarrow$ từ chối \\
Lịch sử tín dụng & Nhóm nợ, số ngày quá hạn, nợ đã xử lý rủi ro &
Nhóm nợ xấu nhất 12 tháng $> 2$ $\Rightarrow$ từ chối \\
Khả năng trả nợ & $\DTI$, thu nhập còn lại sau trả nợ &
$\DTI > 0{,}6$ $\Rightarrow$ từ chối \\
Tài sản bảo đảm & Tỷ lệ vay trên giá trị tài sản &
$\LTV > 0{,}85$ $\Rightarrow$ chuyển thẩm định \\
Chất lượng dữ liệu & Thiếu dữ liệu CIC, khách hàng thin-file &
Không có bản tin CIC $\Rightarrow$ chuyển thẩm định \\
Dấu hiệu bất thường & Số lần truy vấn tín dụng, hành vi giao dịch &
Truy vấn tín dụng $> 5$ lần trong 3 tháng $\Rightarrow$ chuyển thẩm định \\
\bottomrule
\end{tabular}
\end{table}

Ba mức hành động gồm: \code{HARD\_REJECT} (từ chối, có quyền phủ quyết cao nhất),
\code{REVIEW} (buộc chuyển thẩm định thủ công) và \code{WARNING} (ghi nhận vào
danh sách lý do nhưng không thay đổi khuyến nghị).

\subsection{Xác định các hàm mất mát}
\label{subsec:ham-mat-mat}

\paragraph{Hàm mất mát huấn luyện.} Toàn bộ mô hình sử dụng hàm entropy chéo nhị
phân~\eqref{eq:logloss}. Lựa chọn này không phải quy ước tùy tiện: log-loss là
quy tắc chấm điểm chính đáng, nghĩa là kỳ vọng của nó đạt cực tiểu khi và chỉ khi
giá trị dự báo bằng đúng xác suất thật. Đặc điểm này có ý nghĩa quyết định vì hệ
thống sử dụng giá trị $\PD$ trực tiếp chứ không chỉ dùng để xếp hạng.

\paragraph{Bù mất cân bằng lớp.} Với tỷ lệ nợ xấu khoảng 15\%, hàm mất mát được
bổ sung trọng số lớp:

\begin{equation}
\ell_w(y, p) = -\big[\, w_1 \, y \log p + w_0 \,(1 - y) \log (1 - p) \,\big],
\qquad
w_1 = \frac{n_0}{n_1},\ w_0 = 1,
\label{eq:weighted-logloss}
\end{equation}

trong đó $n_0, n_1$ là số quan sát của lớp âm và lớp dương. Việc dùng trọng số
lớp làm xác suất đầu ra bị lệch hệ thống về phía cao hơn thực tế; đây là lý do
bước hiệu chuẩn xác suất trở thành bắt buộc chứ không phải tùy chọn.

\paragraph{Hàm chi phí nghiệp vụ.} Tách biệt với hàm mất mát huấn luyện, hàm chi
phí nghiệp vụ được dùng ở bước chọn ngưỡng quyết định:

\begin{equation}
\mathcal{C}(\tau) = \frac{1}{n}\Big[ c_{FN} \cdot \text{FN}(\tau)
                                   + c_{FP} \cdot \text{FP}(\tau) \Big].
\label{eq:business-cost}
\end{equation}

Việc tách hai hàm này là một lựa chọn thiết kế có chủ đích: mô hình được huấn
luyện để ước lượng xác suất đúng, còn chính sách chi phí được áp dụng ở bước
quyết định. Nhờ vậy, khi tổ chức tín dụng thay đổi khẩu vị rủi ro, chỉ cần điều
chỉnh ngưỡng mà không phải huấn luyện lại mô hình.

\subsection{Lựa chọn chỉ số đánh giá hiệu năng}
\label{subsec:chi-so-danh-gia}

Trên cơ sở bộ chỉ tiêu đã trình bày ở mục~\ref{subsec:so-sanh-mo-hinh}, đề
án bổ sung nhóm chỉ tiêu thứ năm mang tính nghiệp vụ nhằm quy đổi chất lượng
thống kê thành hệ quả kinh doanh. Bảng~\ref{tab:metrics-design} nêu vai trò của
từng nhóm trong quy trình đánh giá.

\begin{table}[htbp]
\centering
\caption{Hệ thống chỉ tiêu đánh giá và vai trò trong quy trình}
\label{tab:metrics-design}
\small
\begin{tabular}{@{}L{2.4cm}L{3.4cm}L{6.6cm}@{}}
\toprule
\textbf{Nhóm} & \textbf{Chỉ tiêu} & \textbf{Vai trò} \\
\midrule
Phân loại phổ quát & Accuracy, Precision, Recall, F1, PR-AUC &
Báo cáo đầy đủ; Recall được ưu tiên hơn Precision do chi phí bỏ sót hồ sơ xấu cao
\\[2pt]
Đặc thù tín dụng & $\AUC$, Gini, $\KS$ &
Tiêu chí chính để xếp hạng và lựa chọn mô hình; không phụ thuộc ngưỡng \\[2pt]
Hiệu chuẩn & Brier, ECE, đường cong hiệu chuẩn &
Điều kiện cần để dùng $\PD$ trong phân vùng quyết định và tính tổn thất kỳ vọng
\\[2pt]
Vận hành & Thời gian huấn luyện, thời gian chấm điểm, mức độ giải thích được &
Quyết định khả năng triển khai thực tế \\[2pt]
Nghiệp vụ & Tỷ lệ duyệt, tỷ lệ nợ xấu nhóm duyệt, chi phí kỳ vọng, bảng gains &
Quy đổi chất lượng thống kê thành hệ quả kinh doanh \\
\bottomrule
\end{tabular}
\end{table}

\subsection{Tích hợp cơ chế giải thích mô hình}
\label{subsec:tich-hop-giai-thich}

Cơ chế giải thích được tích hợp ở hai cấp độ.

\paragraph{Cấp độ toàn cục.} Phân tích SHAP trên một mẫu đại diện của tập kiểm
thử cho biết đặc trưng nào chi phối kết quả dự báo của mô hình nói chung, và mỗi
nhóm nguồn dữ liệu đóng góp bao nhiêu phần vào quyết định. Kết quả này phục vụ
hai mục đích: kiểm chứng rằng mô hình học được quan hệ hợp lý về mặt nghiệp vụ,
và định lượng giá trị của từng nguồn dữ liệu để làm cơ sở cho quyết định đầu tư
dữ liệu.

\paragraph{Cấp độ từng hồ sơ.} Với mỗi hồ sơ được chấm điểm, mô hình sinh danh
sách lý do gồm ba thành phần theo thứ tự ưu tiên: (i) các luật nghiệp vụ bị vi
phạm -- đây là loại lý do tường minh và dễ hiểu nhất; (ii) các chỉ tiêu định
lượng then chốt kèm giá trị thực tế của hồ sơ; và (iii) các đặc trưng có đóng góp
SHAP lớn nhất. Danh sách này được trả về cùng kết quả chấm điểm, cho phép cán bộ
tín dụng đánh giá xem khuyến nghị của mô hình có hợp lý hay không.

\section{Quy trình huấn luyện và tối ưu mô hình}
\label{sec:quy-trinh-huan-luyen}

\subsection{Thiết lập môi trường và công cụ huấn luyện}
\label{subsec:moi-truong-huan-luyen}

Toàn bộ quy trình được hiện thực bằng Python, tổ chức thành một gói phần mềm mà
mỗi gói con tương ứng một mục của chương này: \code{data} cho khâu thu thập và
kiểm tra chất lượng dữ liệu, \code{preprocessing} cho tiền xử lý, \code{features}
cho trích xuất và lựa chọn đặc trưng, \code{labeling} cho gán nhãn và điểm cắt,
\code{models} cho huấn luyện và đánh giá mô hình, \code{decision} cho bộ máy hỗ
trợ quyết định, \code{serving} cho dịch vụ REST và \code{monitoring} cho giám
sát. Cấu trúc thư mục chi tiết được trình bày tại Phụ lục~\ref{app:source}.

Các thư viện chính gồm: \comp{scikit-learn} cho các mô hình cơ sở và khung
pipeline; \comp{XGBoost}, \comp{LightGBM}, \comp{CatBoost} cho các biến thể
gradient boosting; \comp{SHAP} cho giải thích mô hình; \comp{pandas} và
\comp{NumPy} cho xử lý dữ liệu; \comp{FastAPI} cho dịch vụ REST. Toàn bộ tham số
được đặt tập trung trong một file cấu hình duy nhất, và hạt giống ngẫu nhiên được
cố định để bảo đảm khả năng tái lập.

\subsection{Quy trình thử nghiệm và kiểm tra chéo}
\label{subsec:kiem-tra-cheo}

\paragraph{Kiểm định ngoài thời gian cho đánh giá cuối cùng.} Việc đánh giá chất
lượng mô hình được thực hiện trên tập kiểm thử gồm các hồ sơ thuộc giai đoạn sau
giai đoạn huấn luyện. Đây là phương pháp phản ánh đúng điều kiện vận hành.

\paragraph{Kiểm định chéo theo chuỗi thời gian cho tinh chỉnh tham số.} Với việc
tinh chỉnh siêu tham số, đề án sử dụng \code{TimeSeriesSplit} -- phương pháp chia
tập huấn luyện thành nhiều lần khớp liên tiếp, trong đó mỗi lần dùng dữ liệu quá
khứ để huấn luyện và dữ liệu kế tiếp để đánh giá:

\begin{center}
\begin{tabular}{@{}ll@{}}
\toprule
\textbf{Lần khớp} & \textbf{Phân chia} \\
\midrule
1 & Huấn luyện: đoạn 1 \quad$\rightarrow$\quad Đánh giá: đoạn 2 \\
2 & Huấn luyện: đoạn 1--2 \quad$\rightarrow$\quad Đánh giá: đoạn 3 \\
3 & Huấn luyện: đoạn 1--3 \quad$\rightarrow$\quad Đánh giá: đoạn 4 \\
\bottomrule
\end{tabular}
\end{center}

Phương pháp này khác với kiểm định chéo $k$-fold thông thường ở chỗ nó không bao
giờ dùng dữ liệu tương lai để huấn luyện -- điều kiện bắt buộc với dữ liệu có cấu
trúc thời gian.

\subsection{Phương pháp tối ưu hóa siêu tham số}
\label{subsec:toi-uu-sieu-tham-so}

Đề án sử dụng \textbf{tìm kiếm ngẫu nhiên} thay vì tìm kiếm lưới. Lựa chọn này
dựa trên kết quả của Bergstra và Bengio~\cite{bergstra2012}: với cùng một ngân
sách tính toán, tìm kiếm ngẫu nhiên thường tìm được cấu hình tốt hơn tìm kiếm
lưới. Nguyên nhân là trong không gian siêu tham số nhiều chiều, chỉ một số ít
chiều thực sự ảnh hưởng tới kết quả; tìm kiếm lưới lãng phí phần lớn số lần thử
để dò các giá trị khác nhau của những chiều không quan trọng, trong khi tìm kiếm
ngẫu nhiên bao phủ được nhiều giá trị khác biệt hơn ở những chiều quan trọng.

Không gian tìm kiếm được xác định cho ba mô hình có nhiều siêu tham số nhất --
LightGBM, XGBoost và rừng ngẫu nhiên -- bao gồm các tham số về độ sâu/số lá, tốc
độ học, số quan sát tối thiểu tại lá, tỷ lệ lấy mẫu quan sát và đặc trưng, cùng
các hệ số điều chuẩn. Tiêu chí lựa chọn là $\AUC$ trung bình trên các lần khớp
của kiểm định chéo theo chuỗi thời gian.

Để bảo đảm tính nhất quán, việc tinh chỉnh siêu tham số phải được thực hiện trên
\textit{đúng tập đặc trưng} mà mô hình cuối cùng sẽ sử dụng. Nếu
tinh chỉnh trên tập đầy đủ rồi áp dụng tham số cho tập đã rút gọn, siêu tham số
tìm được không còn tối ưu.

\subsection{Mức độ kiểm soát hiện tượng quá khớp}
\label{subsec:kiem-soat-qua-khop}

Đề án áp dụng bốn cơ chế kiểm soát quá khớp, tác động ở các giai đoạn khác nhau
của quy trình.

\paragraph{Điều chuẩn trong hàm mục tiêu.} Hồi quy logistic và SVM sử dụng điều
chuẩn $L_2$; các mô hình boosting sử dụng đồng thời điều chuẩn $L_1$ và $L_2$
trên giá trị lá.

\paragraph{Ràng buộc cấu trúc mô hình.} Giới hạn độ sâu cây, số lá tối đa và số
quan sát tối thiểu tại mỗi lá. Với cây quyết định đơn, bổ sung tỉa nhánh theo độ
phức tạp chi phí~\eqref{eq:ccp}.

\paragraph{Dừng sớm.} Các mô hình boosting theo dõi $\AUC$ trên tập
\code{train\_es} và dừng khi chỉ tiêu này ngừng cải thiện trong một số vòng lặp
liên tiếp. Cần chỉ định $\AUC$ là chỉ tiêu dừng \textit{duy nhất}: khi sử dụng
trọng số lớp, hàm log-loss trên tập đánh giá xấu đi rất sớm do xác suất bị lệch
hệ thống, nên nếu để cơ chế mặc định theo dõi nhiều chỉ tiêu, quá trình huấn
luyện sẽ dừng quá sớm và mô hình bị dưới khớp.

\paragraph{Lấy mẫu ngẫu nhiên.} Rừng ngẫu nhiên và các mô hình boosting sử dụng
lấy mẫu con quan sát và đặc trưng ở mỗi bước, vừa giảm tương quan giữa các cây
vừa đóng vai trò điều chuẩn ngẫu nhiên.

Mức độ quá khớp được đo bằng chênh lệch $\AUC$ giữa tập kiểm định và tập kiểm
thử; chênh lệch lớn là dấu hiệu mô hình phụ thuộc quá mức vào đặc thù của giai
đoạn huấn luyện. Kết quả đo được trình bày tại mục~\ref{subsec:kq-on-dinh}.

\subsection{Phương pháp đánh giá độ ổn định mô hình theo thời gian}
\label{subsec:do-on-dinh}

Độ ổn định được đánh giá trên hai phương diện.

\paragraph{Ổn định của dữ liệu đầu vào.} Đo bằng chỉ số $\PSI$~\eqref{eq:psi} cho
từng đặc trưng, so sánh phân phối giữa giai đoạn huấn luyện và giai đoạn kiểm
thử. Ngưỡng cảnh báo 0,10 và ngưỡng báo động 0,25.

\paragraph{Ổn định của chất lượng dự báo.} Chia tập kiểm thử theo tháng và tính
$\AUC$, $\KS$ cùng tỷ lệ nợ xấu thực tế cho từng tháng. Xu hướng giảm liên tục
của $\AUC$ là dấu hiệu mô hình đang mất dần hiệu lực và cần được huấn luyện lại.

Hai phương diện này bổ trợ nhau: $\PSI$ phát hiện sớm sự thay đổi của dữ liệu đầu
vào ngay cả khi chưa có nhãn, trong khi chất lượng dự báo theo tháng chỉ đo được
sau khi nhãn đã sẵn sàng. Trong vận hành thực tế, $\PSI$ đóng vai trò cảnh báo
sớm còn chất lượng dự báo đóng vai trò xác nhận.

\section*{Kết luận chương 2}
\addcontentsline{toc}{section}{Kết luận chương 2}

Chương 2 đã trình bày quá trình xây dựng mô hình học máy hỗ trợ kiểm duyệt tín
dụng, từ
phân tích bài toán tới thiết kế quy trình huấn luyện mô hình, với bốn nhóm kết
quả.

Về \textit{phân tích bài toán}: xác định biến mục tiêu là xác suất phát sinh nợ
xấu trong 12 tháng, phân biệt ba trạng thái khi gán nhãn, thiết kế thang bảy hạng
rủi ro và năm thành phần của kết quả đầu ra; hệ thống hóa mười yêu cầu chức năng,
bốn yêu cầu phi chức năng và năm nhóm yêu cầu pháp lý.

Về \textit{xử lý dữ liệu}: thiết lập quy tắc ghép nối bảo đảm tính đúng đắn theo
thời điểm, xây dựng bộ mười lăm ràng buộc kiểm tra tính hợp lệ kèm cơ chế kiểm
chứng bằng phép tiêm lỗi, phân biệt hai cách xử lý khác nhau cho dữ liệu không
hợp lệ và giá trị ngoại lệ, thiết kế chiến lược phân chia bốn tập với vai trò
tách bạch.

Về \textit{đặc trưng}: trình bày cách trích xuất năm nhóm đặc trưng cùng thuật
toán luỹ kế bảo đảm tính đúng đắn theo thời điểm với chi phí tuyến tính, và quy
trình lựa chọn ba bước theo năng lực dự báo, mức tương quan và độ ổn định.

Về \textit{mô hình}: lựa chọn mười hai mô hình để so sánh, thiết kế kiến trúc
phân tầng với logic tổng hợp khuyến nghị theo tám mức ưu tiên, tách hàm mất mát
huấn luyện khỏi hàm chi phí nghiệp vụ, và thiết lập bốn cơ chế kiểm soát quá
khớp.

Toàn bộ thiết kế này được kiểm chứng bằng thực nghiệm ở
Chương~\ref{chap:thuc-nghiem}.

% ===========================================================================
% CHƯƠNG 3: TRIỂN KHAI THỰC NGHIỆM VÀ ĐÁNH GIÁ KẾT QUẢ
% ===========================================================================
\chapter{TRIỂN KHAI THỰC NGHIỆM VÀ ĐÁNH GIÁ KẾT QUẢ}
% Tiêu đề chạy rút gọn: tiêu đề chương đầy đủ dài hơn một dòng đầu trang,
% khiến fancyhdr phải nới \headheight và làm lệch khung chữ của cả chương.
\markboth{Chương 3. Triển khai thực nghiệm và đánh giá kết quả}{}
\label{chap:thuc-nghiem}

\noindent\textit{Chương 3 trình bày kịch bản thực nghiệm, kết quả huấn luyện mô
hình học máy và đánh giá hiệu năng mô hình kiểm duyệt tín dụng. Nội dung chương
tập trung mô phỏng vận hành và bàn luận, đối chiếu kết quả thực nghiệm
thu được.}

\section{Thiết lập kịch bản và môi trường thực nghiệm}
\label{sec:moi-truong-thuc-nghiem}

\subsection{Hạ tầng phần cứng, phần mềm và thư viện sử dụng}
\label{subsec:ha-tang}

Toàn bộ thực nghiệm được thực hiện trên một máy tính cá nhân với cấu hình và môi
trường phần mềm trình bày tại Bảng~\ref{tab:environment}.

\begin{table}[htbp]
\centering
\caption{Cấu hình phần cứng và môi trường phần mềm}
\label{tab:environment}
\small
\begin{tabular}{@{}L{4.2cm}L{9.6cm}@{}}
\toprule
\textbf{Thành phần} & \textbf{Thông số} \\
\midrule
Hệ điều hành & Windows 11 (64-bit) \\
Ngôn ngữ lập trình & Python 3.12 \\
Xử lý dữ liệu & pandas, NumPy, PyArrow (định dạng lưu trữ Parquet) \\
Mô hình cơ sở & scikit-learn 1.9 \\
Mô hình boosting & XGBoost, LightGBM, CatBoost \\
Giải thích mô hình & SHAP \\
Dịch vụ REST & FastAPI, Uvicorn \\
Trực quan hóa & Matplotlib, Streamlit \\
\bottomrule
\end{tabular}
\end{table}

Các số đo thời gian trong chương này được thực hiện trên một máy đơn, một tiến
trình, nên phản ánh chi phí tính toán của thuật toán chứ không bao gồm thời gian
truyền qua mạng hay ảnh hưởng của tải đồng thời; chúng có giá trị để \textit{so
sánh tương đối} giữa các mô hình (xem thêm mục~\ref{subsec:han-che-dinh-huong}).

Toàn bộ tham số của thực nghiệm được đặt tập trung trong một file cấu hình duy
nhất và hạt giống ngẫu nhiên được cố định, bảo đảm kết quả tái lập được. Quy
trình thực nghiệm gồm sáu bước, mỗi bước là một chương trình độc lập: chuẩn bị dữ
liệu, phân tích và lựa chọn đặc trưng, tinh chỉnh siêu tham số, huấn luyện mô
hình, mô phỏng quy trình kiểm duyệt, và xuất kết quả.

\subsection{Kịch bản thực nghiệm các mô hình học máy}
\label{subsec:kich-ban}

Thực nghiệm được thiết kế theo nguyên tắc \textbf{so sánh công bằng}: toàn bộ
\NModels{} mô hình được huấn luyện trên cùng một tập dữ liệu, cùng một tập đặc
trưng, cùng cách chia dữ liệu và được đánh giá bằng cùng một bộ chỉ tiêu. Mỗi mô
hình đi qua đúng bốn bước:

\begin{enumerate}[label=(\arabic*)]
  \item Huấn luyện trên tập \code{train\_fit}; các mô hình boosting bổ sung cơ
  chế dừng sớm trên tập \code{train\_es};
  \item Hiệu chuẩn xác suất bằng hồi quy đẳng hướng trên tập \code{valid};
  \item Lựa chọn ngưỡng quyết định theo hàm chi phí trên tập \code{valid};
  \item Đánh giá cuối cùng trên tập \code{test} -- chỉ thực hiện một lần.
\end{enumerate}

Về mặt phương pháp, mọi quyết định liên quan tới mô hình đều dựa trên tập
\code{valid}. Tập \code{test} hoàn toàn không tham gia vào bất kỳ lựa
chọn nào, nhờ đó ước lượng hiệu năng cuối cùng không bị lạc quan.

\subsection{Thiết kế và kiểm chứng tập dữ liệu thực nghiệm}
\label{sec:data-design}

\subsubsection{Cơ sở lựa chọn dữ liệu mô phỏng}

Dữ liệu tín dụng cá nhân thuộc nhóm dữ liệu nhạy cảm, chịu ràng buộc chặt chẽ về
bảo mật. Việc tiếp cận dữ liệu thực của một tổ chức tín dụng cho mục đích nghiên
cứu học thuật đòi hỏi quy trình phê duyệt phức tạp và thường không khả thi trong
khuôn khổ một đề án tốt nghiệp. Mặt khác, các bộ dữ liệu tín dụng công khai
thường có ba hạn chế khiến chúng không phù hợp với đề án này: quy mô nhỏ (vài
nghìn quan sát), không có thông tin thời gian nên không cho phép kiểm định ngoài
thời gian, và đặc biệt là \textit{không có dữ liệu giao dịch} -- nhóm dữ liệu mà
đề án muốn đánh giá đóng góp.

Vì vậy đề án sử dụng dữ liệu mô phỏng được sinh bằng chương trình tự xây dựng.
Cách tiếp cận này có một ưu điểm phương pháp luận quan trọng ngoài việc giải
quyết vấn đề tiếp cận dữ liệu: vì quá trình sinh dữ liệu xác định tường minh xác
suất vỡ nợ thật của từng hồ sơ, ta biết được \textbf{trần Bayes} -- giới hạn lý
thuyết mà không mô hình nào có thể vượt qua. Nhờ đó có thể đánh giá một mô hình
còn cách giới hạn lý thuyết bao xa, điều không thể làm được với dữ liệu thực.

Giới hạn tương ứng của cách tiếp cận này là các chỉ tiêu chất lượng thu được
\textbf{không nên được hiểu là dự báo về hiệu quả trên dữ liệu thực}. Giá trị
của thực nghiệm nằm ở việc so sánh \textit{tương đối} giữa các mô hình trong cùng điều kiện và ở việc
kiểm chứng tính đúng đắn của toàn bộ quy trình từ dữ liệu tới quyết định.

\subsubsection{Mô hình sinh dữ liệu biến tiềm ẩn}

Dữ liệu được sinh theo một mô hình xác suất có cấu trúc, trong đó mỗi khách hàng
$i$ được gán một \textbf{biến tiềm ẩn} $z_i \sim \mathcal{N}(0, 1)$ biểu diễn
mức độ tín nhiệm nội tại -- đại lượng không quan sát được trong thực tế. Từ
$z_i$, chương trình sinh ra ba nhóm dữ liệu quan sát được:

\begin{itemize}
  \item \textbf{Thông tin nhân khẩu và tài chính}: thu nhập, nghề nghiệp, thâm
  niên công việc phụ thuộc $z_i$ với nhiễu;
  \item \textbf{Lịch sử tín dụng tại CIC}: điểm CIC, nhóm nợ, số ngày quá hạn là
  các \textit{phép đo có nhiễu} của $z_i$;
  \item \textbf{Hành vi giao dịch}: tính ổn định dòng tiền, tỷ trọng chi tiêu
  nhóm rủi ro cũng phụ thuộc $z_i$ với nhiễu độc lập.
\end{itemize}

Nhãn vỡ nợ được sinh từ một hàm log-odds phi tuyến của $z_i$ và các đặc trưng đã
tính đúng theo thời điểm:

\begin{equation}
\PD_i^{\text{thật}} = \sigma\big( \eta(z_i, \bm{x}_i) + \beta_0 \big),
\qquad
y_i \sim \text{Bernoulli}\big( \PD_i^{\text{thật}} \big),
\label{eq:datagen}
\end{equation}

trong đó hệ số chặn $\beta_0$ được hiệu chỉnh bằng phương pháp số để đạt tỷ lệ vỡ
nợ mục tiêu.

Cấu trúc này tạo ra một đặc điểm quan trọng: vì CIC và hành vi giao dịch chỉ là
các phép đo \textit{có nhiễu} của $z_i$, mô hình học máy chỉ có thể phục hồi một
phần thông tin về $z_i$. Điều này tạo ra một giới hạn trên tự nhiên cho chất
lượng dự báo -- chính là trần Bayes -- giống như tình huống thực tế, nơi không có
mô hình nào đạt $\AUC = 1$.

Quá trình sinh dữ liệu còn đưa vào ba hiện tượng đặc thù của dữ liệu tín dụng
thực: nhóm khách hàng thiếu lịch sử tín dụng (chiếm \ThinFileRate), giá trị khuyết
ở một số trường khai báo, và hiện tượng trôi dạt theo thời gian với mức độ xác
định trước.

\subsubsection{Quy mô và đặc điểm tập dữ liệu}

Bảng~\ref{tab:dataset} trình bày quy mô các bảng dữ liệu thu được. Tập dữ liệu
gồm \NCustomers{} khách hàng, \NAccounts{} tài khoản, \NTransactions{} giao dịch
và \NTotal{} hồ sơ vay trải trên khoảng thời gian hai năm, với thời gian sinh
\GenSeconds{} giây.

\begin{table}[!htbp]
\centering
\caption{Quy mô các bảng dữ liệu sử dụng trong đề án}
\label{tab:dataset}
\small
\setlength{\tabcolsep}{5.0pt}
\begin{tabular}{@{}L{3.0cm}R{2.4cm}L{7.6cm}@{}}
\toprule
\textbf{Bảng dữ liệu} & \textbf{Số bản ghi} & \textbf{Nội dung} \\
\midrule
Khách hàng & 40.000 & Thông tin nhân khẩu học, nghề nghiệp, thu nhập \\
Tài khoản & 73.023 & Tài khoản thanh toán/tiết kiệm và số dư bình quân \\
Giao dịch & 4.538.127 & Giao dịch tài khoản dùng để tính đặc trưng hành vi \\
Hồ sơ vay & 60.000 & Hồ sơ đề nghị cấp tín dụng \\
Bản tin CIC & 60.000 & Thông tin lịch sử tín dụng từ CIC \\
Nhãn rủi ro & 60.000 & Trạng thái vỡ nợ sau cửa sổ quan sát \\
\bottomrule
\end{tabular}
\end{table}


Hình~\ref{fig:data-overview} mô tả phân bố của một số biến chính trong tập dữ
liệu.

\begin{figure}[!htbp]
\centering
\cffig[0.94\textwidth]{data_overview}
\caption{Phân bố các biến chính trong tập dữ liệu thực nghiệm}
\label{fig:data-overview}
\end{figure}

\subsubsection{Kiểm tra chất lượng dữ liệu}
\label{subsec:kq-chat-luong-du-lieu}

Bộ \NConstraints{} ràng buộc nghiệp vụ được áp dụng cho ba bảng dữ liệu chính,
phát hiện \NViolations{} bản ghi vi phạm. Kết quả này không đáng ngạc nhiên với
dữ liệu mô phỏng, nhưng bản thân nó \textit{không chứng minh được} rằng cơ chế
kiểm tra hoạt động đúng -- một bộ kiểm tra hỏng cũng cho kết quả tương tự.

Để phân biệt hai khả năng, đề án thực hiện phép tiêm lỗi: tạo bản sao dữ liệu, cố
ý đưa vào 1\% lỗi ứng với từng loại ràng buộc, rồi chạy lại cơ chế kiểm tra. Kết
quả tại Bảng~\ref{tab:validity} cho thấy toàn bộ lỗi đã tiêm đều được phát hiện
(tỷ lệ phát hiện \ValidityDetected).

\begin{table}[!htbp]
\centering
\caption{Kiểm chứng cơ chế phát hiện dữ liệu không hợp lệ bằng phép tiêm lỗi}
\label{tab:validity}
\small
\setlength{\tabcolsep}{5.0pt}
\begin{tabular}{@{}L{3.3cm}L{5.2cm}R{1.7cm}R{1.7cm}L{1.9cm}@{}}
\toprule
\textbf{Mã ràng buộc} & \textbf{Nội dung ràng buộc} & \textbf{Số lỗi tiêm vào} & \textbf{Số lỗi phát hiện} & \textbf{Kết quả} \\
\midrule
CUS\_\allowbreak AGE\_\allowbreak RANGE & Tuổi khách hàng phải nằm trong khoảng 18-75 & 400 & 400 & Phát hiện đủ \\
CUS\_\allowbreak INCOME\_\allowbreak POSITIVE & Thu nhập khai báo phải lớn hơn 0 & 400 & 400 & Phát hiện đủ \\
APP\_\allowbreak AMOUNT\_\allowbreak POSITIVE & Số tiền đề nghị vay phải lớn hơn 0 & 600 & 600 & Phát hiện đủ \\
APP\_\allowbreak RATE\_\allowbreak RANGE & Lãi suất phải nằm trong khoảng 0-40\%/năm & 600 & 600 & Phát hiện đủ \\
CIC\_\allowbreak SCORE\_\allowbreak RANGE & Điểm CIC phải nằm trong khoảng 300-850 & 600 & 600 & Phát hiện đủ \\
CIC\_\allowbreak GROUP\_\allowbreak ORDER & Nhóm nợ xấu nhất 12 tháng không thể tốt hơn nhóm nợ hiện tại & 600 & 600 & Phát hiện đủ \\
\bottomrule
\end{tabular}

\vspace{2pt}
{\footnotesize\itshape Phép tiêm lỗi được thực hiện trên một bản sao dữ liệu với tỷ lệ 1\% mỗi loại lỗi; dữ liệu gốc không bị thay đổi.}
\end{table}


Về giá trị khuyết, \NMissingVars{} trên tổng số đặc trưng có dữ liệu trống, với
tỷ lệ cao nhất là \MaxMissingPct{} -- tương ứng nhóm khách hàng thiếu lịch sử tín
dụng (Bảng~\ref{tab:missing}). Về giá trị ngoại lệ, \NOutlierVars{} đặc trưng có
quan sát nằm ngoài ngưỡng $3 \times$IQR (Bảng~\ref{tab:outliers}); các giá trị
này được xử lý bằng phép cắt ngọn thay vì loại bỏ quan sát.

\begin{table}[!htbp]
\centering
\caption{Các đặc trưng có giá trị khuyết và chiến lược xử lý}
\label{tab:missing}
\small
\setlength{\tabcolsep}{5.0pt}
\begin{tabular}{@{}L{4.4cm}L{2.0cm}R{2.4cm}L{5.4cm}@{}}
\toprule
\textbf{Đặc trưng} & \textbf{Nhóm} & \textbf{Tỷ lệ khuyết (\%)} & \textbf{Cách xử lý} \\
\midrule
oldest\_\allowbreak account\_\allowbreak months & Lịch sử tín dụng (CIC) & 11,65 & Quy ước nghiệp vụ + cờ chỉ báo khuyết \\
credit\_\allowbreak utilization & Lịch sử tín dụng (CIC) & 10,93 & Quy ước nghiệp vụ + cờ chỉ báo khuyết \\
cic\_score & Lịch sử tín dụng (CIC) & 9,66 & Quy ước nghiệp vụ + cờ chỉ báo khuyết \\
current\_\allowbreak group & Lịch sử tín dụng (CIC) & 9,66 & Quy ước nghiệp vụ + cờ chỉ báo khuyết \\
total\_\allowbreak outstanding & Lịch sử tín dụng (CIC) & 9,66 & Quy ước nghiệp vụ + cờ chỉ báo khuyết \\
dpd\_\allowbreak max\_\allowbreak 12m & Lịch sử tín dụng (CIC) & 9,66 & Quy ước nghiệp vụ + cờ chỉ báo khuyết \\
worst\_\allowbreak group\_\allowbreak 12m & Lịch sử tín dụng (CIC) & 9,66 & Quy ước nghiệp vụ + cờ chỉ báo khuyết \\
num\_\allowbreak bad\_\allowbreak debt\_\allowbreak institutions & Lịch sử tín dụng (CIC) & 9,66 & Quy ước nghiệp vụ + cờ chỉ báo khuyết \\
monthly\_\allowbreak debt\_\allowbreak payment & Lịch sử tín dụng (CIC) & 9,66 & Quy ước nghiệp vụ + cờ chỉ báo khuyết \\
\bottomrule
\end{tabular}
\end{table}


\begin{table}[!htbp]
\centering
\caption{Các đặc trưng có tỷ lệ giá trị ngoại lệ cao nhất (quy tắc $3\times$IQR)}
\label{tab:outliers}
\small
\setlength{\tabcolsep}{5.0pt}
\begin{tabular}{@{}L{5.2cm}R{3.0cm}R{2.6cm}R{2.6cm}@{}}
\toprule
\textbf{Đặc trưng} & \textbf{Tỷ lệ ngoại lệ (\%)} & \textbf{Phân vị 1\%} & \textbf{Phân vị 99\%} \\
\midrule
utility\_\allowbreak amt\_\allowbreak 7d & 24,75 & 0 & 1,49 \\
ecom\_\allowbreak amt\_\allowbreak 7d & 21,90 & 0 & 3,22 \\
current\_\allowbreak group & 20,84 & 1 & 5 \\
outflow\_\allowbreak ratio\_\allowbreak 30d & 18,67 & 0 & 14447060,00 \\
dpd\_\allowbreak x\_\allowbreak institutions & 17,97 & 0 & 87,20 \\
cash\_\allowbreak wd\_\allowbreak amt\_\allowbreak 7d & 14,86 & 0 & 4,91 \\
cash\_\allowbreak ratio\_\allowbreak 7d & 14,86 & 0 & 1,00 \\
cash\_\allowbreak wd\_\allowbreak cnt\_\allowbreak 7d & 14,85 & 0 & 2 \\
repay\_\allowbreak amt\_\allowbreak 7d & 13,29 & 0 & 8,64 \\
repay\_\allowbreak burden\_\allowbreak 7d & 13,29 & 0 & 6209040,00 \\
\bottomrule
\end{tabular}

\vspace{2pt}
{\footnotesize\itshape Giá trị ngoại lệ được xử lý bằng phép cắt ngọn tại phân vị 1\% và 99\% học từ tập huấn luyện, không loại bỏ quan sát.}
\end{table}


\subsubsection{Kết quả gán nhãn rủi ro}
\label{subsec:kq-gan-nhan}

Bảng~\ref{tab:label} trình bày phân bố ba trạng thái rủi ro sau khi gán nhãn theo
tiêu chuẩn nợ nhóm 3--5 với cửa sổ quan sát 12 tháng.

\begin{table}[!htbp]
\centering
\caption{Phân bố trạng thái rủi ro sau khi gán nhãn}
\label{tab:label}
\small
\setlength{\tabcolsep}{5.0pt}
\begin{tabular}{@{}L{8.0cm}R{2.6cm}R{2.6cm}@{}}
\toprule
\textbf{Trạng thái} & \textbf{Số hồ sơ} & \textbf{Tỷ trọng (\%)} \\
\midrule
Tốt (không phát sinh nợ xấu) & 50.867 & 84,78 \\
Xấu (nợ nhóm 3-5, quá hạn ≥ 90 ngày) & 9.133 & 15,22 \\
\bottomrule
\end{tabular}
\end{table}


Cơ sở định lượng cho lựa chọn cửa sổ 12 tháng được trình bày tại
Bảng~\ref{tab:vintage}: đường cong tỷ lệ nợ xấu tích luỹ cho thấy cửa sổ 12 tháng
ghi nhận được \VintageCoverage{} tổng số trường hợp vỡ nợ, trong khi cửa sổ 6
tháng chỉ ghi nhận được khoảng 56\%. Rút ngắn cửa sổ xuống 6 tháng sẽ khiến gần
một nửa số hồ sơ thực sự xấu bị gán nhãn tốt -- sai lệch đủ lớn để làm hỏng mô
hình.

\begin{table}[!htbp]
\centering
\caption{Tỷ lệ nợ xấu tích luỹ theo số tháng kể từ khi cấp tín dụng}
\label{tab:vintage}
\small
\setlength{\tabcolsep}{5.0pt}
\begin{tabular}{@{}R{3.4cm}R{3.5cm}R{3.5cm}R{3.5cm}@{}}
\toprule
\textbf{Tháng kể từ khi cấp tín dụng} & \textbf{Số hồ sơ phát sinh nợ xấu} & \textbf{Tỷ lệ nợ xấu tích luỹ (\%)} & \textbf{Tỷ lệ ghi nhận được (\%)} \\
\midrule
3 & 1.460 & 2,43 & 16,0 \\
6 & 5.078 & 8,46 & 55,6 \\
9 & 7.632 & 12,72 & 83,6 \\
12 & 9.133 & 15,22 & 100 \\
15 & 9.133 & 15,22 & 100 \\
18 & 9.133 & 15,22 & 100 \\
\bottomrule
\end{tabular}

\vspace{2pt}
{\footnotesize\itshape Cửa sổ quan sát 12 tháng được chọn tại điểm đường cong tích luỹ bắt đầu đi ngang.}
\end{table}


\subsection{Phương pháp phân chia tập dữ liệu thử nghiệm}
\label{subsec:kq-chia-tap}

Kết quả phân chia dữ liệu theo thời gian được trình bày tại
Bảng~\ref{tab:split}.

\begin{table}[!htbp]
\centering
\caption{Phân chia dữ liệu theo thời gian}
\label{tab:split}
\small
\setlength{\tabcolsep}{5.0pt}
\begin{tabular}{@{}L{3.6cm}R{2.0cm}R{1.6cm}L{5.6cm}@{}}
\toprule
\textbf{Tập} & \textbf{Số hồ sơ} & \textbf{Tỷ lệ} & \textbf{Vai trò} \\
\midrule
Huấn luyện (train\_\allowbreak fit) & 35.700 & 59,5\% & Ước lượng tham số mô hình \\
Dừng sớm (train\_es) & 6.300 & 10,5\% & Xác định số vòng lặp tối ưu \\
Kiểm định (valid) & 9.000 & 15,0\% & Hiệu chuẩn xác suất, chọn ngưỡng \\
Kiểm thử (test) & 9.000 & 15,0\% & Đánh giá cuối cùng, dùng một lần \\
\bottomrule
\end{tabular}

\vspace{2pt}
{\footnotesize\itshape Giai đoạn huấn luyện: 2024-07-01 đến 2025-11-24; giai đoạn kiểm thử: 2026-03-12 đến 2026-06-30.}
\end{table}


Tỷ lệ nợ xấu tăng dần qua các giai đoạn, từ \DRTrain{}
ở tập huấn luyện lên \DRTest{} ở tập kiểm thử. Đây chính là biểu hiện của hiện
tượng trôi dạt đã được đưa vào quá trình sinh dữ liệu. Sự chênh lệch này khiến
bài toán trở nên khó hơn so với chia ngẫu nhiên và làm cho kết quả đánh giá sát
với điều kiện vận hành thực tế hơn.

\section{Kết quả huấn luyện và tính toán mô hình}
\label{sec:ket-qua-huan-luyen}

\subsection{Kết quả xây dựng và lựa chọn đặc trưng}
\label{subsec:kq-dac-trung}

Tập đặc trưng ban đầu gồm \NFeaturesInitial{} biến, phân bố theo năm nhóm nguồn
dữ liệu như trình bày tại Bảng~\ref{tab:feature-groups}.

\begin{table}[!htbp]
\centering
\caption{Cơ cấu tập đặc trưng theo nhóm và nguồn dữ liệu}
\label{tab:feature-groups}
\small
\setlength{\tabcolsep}{5.0pt}
\begin{tabular}{@{}L{4.2cm}L{4.2cm}R{3.0cm}@{}}
\toprule
\textbf{Nhóm} & \textbf{Nguồn dữ liệu} & \textbf{Số đặc trưng} \\
\midrule
Hành vi giao dịch & GIAO DỊCH & 99 \\
Lịch sử tín dụng (CIC) & BẢN TIN CIC & 16 \\
Thông tin cá nhân & HỒ SƠ KH & 15 \\
Biến tổ hợp & TÍNH TOÁN & 13 \\
Tham số khoản vay & ĐƠN VAY & 11 \\
\bottomrule
\end{tabular}
\end{table}


\subsubsection{Phân tích Information Value}

Bảng~\ref{tab:iv-top} liệt kê mười lăm đặc trưng có năng lực dự báo cao nhất.

\begin{table}[!htbp]
\centering
\caption{Mười lăm đặc trưng có năng lực dự báo cao nhất theo Information Value}
\label{tab:iv-top}
\small
\setlength{\tabcolsep}{5.0pt}
\begin{tabular}{@{}L{5.0cm}L{2.4cm}R{1.8cm}L{5.0cm}@{}}
\toprule
\textbf{Đặc trưng} & \textbf{Nhóm} & \textbf{IV} & \textbf{Mức độ} \\
\midrule
cic\_score & Lịch sử tín dụng (CIC) & 1,244 & Rất mạnh (cần kiểm tra rò rỉ nhãn) \\
worst\_\allowbreak group\_\allowbreak 12m & Lịch sử tín dụng (CIC) & 1,179 & Rất mạnh (cần kiểm tra rò rỉ nhãn) \\
dpd\_\allowbreak max\_\allowbreak 12m & Lịch sử tín dụng (CIC) & 1,106 & Rất mạnh (cần kiểm tra rò rỉ nhãn) \\
current\_\allowbreak group & Lịch sử tín dụng (CIC) & 0,959 & Rất mạnh (cần kiểm tra rò rỉ nhãn) \\
dpd\_\allowbreak x\_\allowbreak institutions & Biến tổ hợp & 0,948 & Rất mạnh (cần kiểm tra rò rỉ nhãn) \\
num\_\allowbreak bad\_\allowbreak debt\_\allowbreak institutions & Lịch sử tín dụng (CIC) & 0,821 & Rất mạnh (cần kiểm tra rò rỉ nhãn) \\
disposable\_\allowbreak income & Biến tổ hợp & 0,739 & Rất mạnh (cần kiểm tra rò rỉ nhãn) \\
disposable\_\allowbreak ratio & Biến tổ hợp & 0,660 & Rất mạnh (cần kiểm tra rò rỉ nhãn) \\
total\_\allowbreak avg\_\allowbreak balance & Thông tin cá nhân & 0,564 & Rất mạnh (cần kiểm tra rò rỉ nhãn) \\
inflow\_90d & Hành vi giao dịch & 0,546 & Rất mạnh (cần kiểm tra rò rỉ nhãn) \\
inflow\_\allowbreak 180d & Hành vi giao dịch & 0,521 & Rất mạnh (cần kiểm tra rò rỉ nhãn) \\
monthly\_\allowbreak income & Thông tin cá nhân & 0,503 & Rất mạnh (cần kiểm tra rò rỉ nhãn) \\
txn\_\allowbreak amt\_\allowbreak std\_\allowbreak 90d & Hành vi giao dịch & 0,448 & Mạnh \\
txn\_\allowbreak amt\_\allowbreak std\_\allowbreak 180d & Hành vi giao dịch & 0,443 & Mạnh \\
net\_\allowbreak flow\_\allowbreak 180d & Hành vi giao dịch & 0,432 & Mạnh \\
\bottomrule
\end{tabular}
\end{table}


Kết quả cho thấy nhóm đặc trưng CIC chiếm ưu thế ở phần đầu bảng xếp hạng, với
giá trị $\IV$ cao nhất đạt \MaxIV{}. Điều này phù hợp với trực giác nghiệp vụ:
lịch sử trả nợ trong quá khứ là chỉ báo mạnh nhất về khả năng trả nợ trong tương
lai. Giá trị $\IV$ vượt ngưỡng 0,5 thông thường là dấu hiệu cần kiểm tra rò rỉ
nhãn; trong trường hợp này, việc kiểm tra ràng buộc thời gian
(CIC\_QUERY\_BEFORE\_APPLY) xác nhận mọi bản tin CIC đều được truy vấn trước thời
điểm nộp hồ sơ, nên $\IV$ cao phản ánh bản chất của dữ liệu chứ không phải lỗi
phương pháp.

Bảng~\ref{tab:iv-group} tổng hợp $\IV$ theo nhóm nguồn dữ liệu.

\begin{table}[!htbp]
\centering
\caption{Information Value tổng hợp theo nhóm nguồn dữ liệu}
\label{tab:iv-group}
\small
\setlength{\tabcolsep}{5.0pt}
\begin{tabular}{@{}L{3.426cm}R{2.474cm}R{2.474cm}R{2.855cm}R{2.665cm}@{}}
\toprule
\textbf{Nhóm} & \textbf{Số đặc trưng} & \textbf{Tổng IV} & \textbf{IV trung bình} & \textbf{IV lớn nhất} \\
\midrule
Hành vi giao dịch & 99 & 13,945 & 0,141 & 0,546 \\
Lịch sử tín dụng (CIC) & 16 & 6,780 & 0,424 & 1,244 \\
Biến tổ hợp & 13 & 4,253 & 0,327 & 0,948 \\
Thông tin cá nhân & 15 & 1,604 & 0,107 & 0,564 \\
Tham số khoản vay & 11 & 0,589 & 0,054 & 0,190 \\
\bottomrule
\end{tabular}
\end{table}


Xét theo $\IV$ trung bình, nhóm CIC dẫn đầu rõ rệt: mỗi đặc trưng CIC mang nhiều
thông tin hơn hẳn. Nhưng xét theo \textit{tổng} $\IV$, nhóm hành vi giao dịch lại
chiếm vị trí cao nhất nhờ số lượng đặc trưng lớn. Hai cách nhìn này không mâu
thuẫn mà bổ sung cho nhau: dữ liệu CIC là nguồn thông tin cô đọng nhất, còn dữ
liệu giao dịch là nguồn thông tin dồi dào nhất. Với nhóm khách hàng thiếu lịch sử
tín dụng, nơi dữ liệu CIC gần như trống, nhóm thứ hai trở thành nguồn thông tin
chủ yếu.

\begin{figure}[!htbp]
\centering
\cffig[0.95\textwidth]{iv_ranking}
\caption{Xếp hạng đặc trưng theo Information Value}
\label{fig:iv-ranking}
\end{figure}

\subsubsection{Kết quả lựa chọn tập đặc trưng}

Quy trình lọc ba bước rút gọn tập đặc trưng từ \NFeaturesInitial{} xuống
\NFeaturesSelected{} biến (Bảng~\ref{tab:selection}).

\begin{table}[!htbp]
\centering
\caption{Kết quả lựa chọn tập đặc trưng qua ba bộ lọc}
\label{tab:selection}
\small
\setlength{\tabcolsep}{5.0pt}
\begin{tabular}{@{}L{7.4cm}R{3.2cm}R{3.2cm}@{}}
\toprule
\textbf{Bước} & \textbf{Số đặc trưng} & \textbf{Đã loại} \\
\midrule
Tập đặc trưng ban đầu & 154 & 0 \\
Sau lọc theo IV (≥ 0.02) & 126 & 28 \\
Sau lọc theo tương quan (|r| ≤ 0.95) & 116 & 10 \\
Sau lọc theo độ ổn định (PSI ≤ 0.25) & 116 & 0 \\
\bottomrule
\end{tabular}
\end{table}


\begin{table}[!htbp]
\centering
\caption{Cơ cấu tập đặc trưng trước và sau lựa chọn}
\label{tab:selection-group}
\small
\setlength{\tabcolsep}{5.0pt}
\begin{tabular}{@{}L{4.0cm}R{2.8cm}R{3.0cm}R{3.2cm}@{}}
\toprule
\textbf{Nhóm} & \textbf{Ban đầu} & \textbf{Sau lựa chọn} & \textbf{Tỷ lệ giữ lại (\%)} \\
\midrule
Hành vi giao dịch & 99 & 76 & 76,8 \\
Lịch sử tín dụng (CIC) & 16 & 14 & 87,5 \\
Thông tin cá nhân & 15 & 8 & 53,3 \\
Biến tổ hợp & 13 & 12 & 92,3 \\
Tham số khoản vay & 11 & 6 & 54,5 \\
\bottomrule
\end{tabular}
\end{table}


Phân tích cơ cấu sau lựa chọn (Bảng~\ref{tab:selection-group}) cho thấy các bộ
lọc hoạt động đúng như thiết kế: nhóm biến tổ hợp giữ lại tỷ lệ cao nhất -- hợp
lý vì mỗi biến trong nhóm này được thiết kế để biểu diễn một chỉ tiêu nghiệp vụ
riêng biệt nên ít trùng lặp. Ngược lại, nhóm tham số khoản vay bị loại nhiều
nhất, chủ yếu do năng lực dự báo thấp.

Bước lọc theo độ ổn định không loại bỏ đặc trưng nào: chỉ số $\PSI$
lớn nhất trên toàn bộ tập đặc trưng chỉ đạt \MaxPSI{}, thấp hơn nhiều so với
ngưỡng cảnh báo 0,10. Kết quả này cho thấy mức trôi dạt trong dữ liệu thực nghiệm
tập trung ở \textit{quan hệ giữa đặc trưng và nhãn} chứ không ở phân phối biên
của đặc trưng -- một dạng trôi dạt khó phát hiện hơn bằng giám sát $\PSI$ đơn
thuần.

\subsection{Kết quả trên các mô hình cơ sở}
\label{subsec:kq-mo-hinh-co-so}

Nhóm mô hình cơ sở gồm hồi quy logistic, thẻ điểm WOE và cây quyết định -- các mô
hình có cấu trúc đơn giản và khả năng giải thích cao.

Hồi quy logistic đạt $\AUC = \LRAUC{}$ trên tập kiểm thử, thẻ điểm WOE đạt
\WOEAUC{}, còn cây quyết định đơn đạt \DTAUC{}. Khoảng cách giữa hai mô hình
tuyến tính và cây quyết định là đáng kể -- gần 0,04 điểm $\AUC$ -- cho thấy hạn
chế về phương sai cao của cây đơn đã được phân tích ở
mục~\ref{subsec:decision-tree} thể hiện rõ trong thực nghiệm.

Việc hồi quy logistic -- mô hình đơn giản nhất trong nhóm -- lại đạt kết quả gần
sát với các mô hình phức tạp là một quan sát quan trọng, sẽ được bàn kỹ ở
mục~\ref{subsec:so-sanh-thuat-toan}.

\subsection{Kết quả trên các mô hình học máy nâng cao}
\label{subsec:kq-mo-hinh-nang-cao}

Nhóm mô hình nâng cao gồm các mô hình học tập hợp (rừng ngẫu nhiên, Extra Trees,
bốn biến thể gradient boosting), máy vectơ hỗ trợ, mạng nơ-ron và mô hình kết hợp
xếp tầng.

Trong nhóm này, mô hình kết hợp xếp tầng đạt $\AUC$ danh nghĩa cao nhất
(\TopAUC{}), ngay sau là ba biến thể gradient boosting: CatBoost \CatAUC{},
XGBoost \XGBAUC{} và LightGBM \LGBAUC{}. Máy vectơ hỗ trợ đạt \SVMAUC{}, rừng
ngẫu nhiên đạt \RFAUC{} và Hist Gradient Boosting đạt \GBAUC{}.

Cần nhấn mạnh ngay rằng khoảng cách giữa các mô hình dẫn đầu là \textit{không
đáng kể}: chênh lệch giữa mô hình đứng đầu và mô hình thứ hai chỉ \GapTopSecond{}
điểm $\AUC$. Mục~\ref{subsec:kq-chon-mo-hinh} phân tích vì sao mức chênh lệch
này không đủ cơ sở để kết luận một thuật toán vượt trội, và vì sao mô hình được
chọn triển khai lại không phải mô hình có $\AUC$ cao nhất.

Kết quả của SVM cần được bình luận riêng. Mặc dù sử dụng phép xấp xỉ Nystroem
thay vì giải bài toán đối ngẫu đầy đủ, mô hình vẫn đạt mức chất lượng tương đương
rừng ngẫu nhiên và chỉ kém mô hình tốt nhất \GapBestSVM{} điểm $\AUC$, với thời
gian huấn luyện khoảng \SVMTrainSec{} giây. Kết quả này cho thấy phép xấp xỉ
hạt nhân là giải pháp khả thi để áp dụng SVM ở quy mô dữ liệu mà SVM hạt nhân đầy đủ không huấn
luyện được trong thời gian chấp nhận được.

\subsection{Tối ưu hóa siêu tham số cho mô hình}
\label{subsec:kq-tuning}

Kết quả trình bày ở hai mục trên đã sử dụng siêu tham số được tinh chỉnh ở bước
này; mục này trình bày chi tiết quá trình tinh chỉnh đó.

Ba mô hình có nhiều siêu tham số nhất -- LightGBM, XGBoost và rừng ngẫu nhiên --
được tinh chỉnh bằng tìm kiếm ngẫu nhiên kết hợp kiểm định chéo theo chuỗi thời
gian trên tập huấn luyện. Kết quả trình bày tại Bảng~\ref{tab:tuning}.

\begin{table}[!htbp]
\centering
\caption{Kết quả tinh chỉnh siêu tham số bằng tìm kiếm ngẫu nhiên}
\label{tab:tuning}
\small
\setlength{\tabcolsep}{5.0pt}
\begin{tabular}{@{}L{4.6cm}R{2.6cm}R{2.6cm}R{3.4cm}@{}}
\toprule
\textbf{Mô hình} & \textbf{Số lần thử} & \textbf{AUC (CV)} & \textbf{Thời gian (phút)} \\
\midrule
lightgbm & 14 & 0,8499 & 1,5 \\
xgboost & 14 & 0,8623 & 2,9 \\
random\_\allowbreak forest & 14 & 0,8576 & 9,5 \\
\bottomrule
\end{tabular}

\vspace{2pt}
{\footnotesize\itshape Kiểm định chéo theo chuỗi thời gian (TimeSeriesSplit) ba lần khớp trên tập huấn luyện.}
\end{table}


Hai điểm cần lưu ý về phương pháp. Thứ nhất, việc tinh chỉnh được thực hiện trên
\textit{đúng tập đặc trưng} mà mô hình cuối cùng sử dụng -- tập \NFeaturesSelected{}
biến sau lựa chọn, không phải tập đầy đủ ban đầu. Nếu tinh chỉnh trên tập đầy đủ
rồi áp dụng tham số cho tập đã rút gọn, siêu tham số tìm được không còn tối ưu.
Thứ hai, toàn bộ quá trình tinh chỉnh chỉ sử dụng tập huấn luyện; tập kiểm định
và tập kiểm thử không tham gia, nhờ đó không tạo ra rò rỉ thông tin.

\section{Đánh giá hiệu năng và so sánh kết quả thực nghiệm}
\label{sec:danh-gia-hieu-nang}

\subsection{Đánh giá qua các chỉ số phân loại}
\label{subsec:kq-chi-so-phan-loai}

Bảng~\ref{tab:model-quality} trình bày đầy đủ kết quả của \NModels{} mô hình theo
bảy chỉ tiêu chất lượng dự báo.

\begin{table}[!htbp]
\centering
\caption{So sánh chất lượng dự báo của các mô hình trên tập kiểm thử}
\label{tab:model-quality}
\small
\setlength{\tabcolsep}{3.0pt}
\begin{tabular}{@{}L{4.374cm}R{1.350cm}R{1.350cm}R{1.350cm}R{1.350cm}R{1.350cm}R{1.350cm}R{1.350cm}@{}}
\toprule
\textbf{Mô hình} & \textbf{AUC} & \textbf{Gini} & \textbf{KS} & \textbf{PR-AUC} & \textbf{F1} & \textbf{Brier} & \textbf{ECE} \\
\midrule
\textbf{Kết hợp xếp tầng (Stacking)} & \textbf{0,8620} & \textbf{0,7240} & \textbf{0,5702} & \textbf{0,6141} & \textbf{0,4445} & \textbf{0,0981} & \textbf{0,0063} \\
CatBoost & 0,8614 & 0,7229 & 0,5689 & 0,6257 & 0,4731 & 0,0976 & 0,0075 \\
XGBoost & 0,8612 & 0,7223 & 0,5729 & 0,6213 & 0,4689 & 0,0980 & 0,0101 \\
LightGBM & 0,8590 & 0,7181 & 0,5647 & 0,6107 & 0,4482 & 0,0988 & 0,0053 \\
Hồi quy logistic & 0,8579 & 0,7158 & 0,5629 & 0,5914 & 0,4418 & 0,1007 & 0,0070 \\
Gradient Boosting (histogram) & 0,8575 & 0,7150 & 0,5594 & 0,5999 & 0,4453 & 0,0993 & 0,0107 \\
Thẻ điểm WOE + Hồi quy logistic & 0,8570 & 0,7140 & 0,5621 & 0,6025 & 0,4627 & 0,1000 & 0,0112 \\
Mạng nơ-ron truyền thẳng (MLP) & 0,8560 & 0,7119 & 0,5602 & 0,6012 & 0,4649 & 0,0999 & 0,0083 \\
Máy vectơ hỗ trợ (SVM-RBF) & 0,8540 & 0,7080 & 0,5619 & 0,5953 & 0,4517 & 0,1005 & 0,0070 \\
Rừng ngẫu nhiên & 0,8527 & 0,7054 & 0,5531 & 0,5930 & 0,4654 & 0,1008 & 0,0064 \\
Cây cực ngẫu nhiên (Extra Trees) & 0,8471 & 0,6942 & 0,5448 & 0,5808 & 0,4360 & 0,1030 & 0,0078 \\
Cây quyết định (CART) & 0,8206 & 0,6412 & 0,4970 & 0,5299 & 0,4052 & 0,1092 & 0,0066 \\
\bottomrule
\end{tabular}

\vspace{2pt}
{\footnotesize\itshape Sắp xếp theo ROC-AUC giảm dần. Trần Bayes (tính từ xác suất vỡ nợ thực của quá trình sinh dữ liệu) là 0,9109.}
\end{table}


Mô hình được chọn triển khai -- \BestModel{} -- đạt $\AUC = \BestAUC{}$,
$\text{Gini} = \BestGini{}$, $\KS = \BestKS{}$ và $F_1 = \BestFOne{}$. Về chất
lượng hiệu chuẩn, điểm Brier đạt \BestBrier{} và sai số hiệu chuẩn kỳ vọng chỉ
\BestECE{}; mức sai số thấp này cho phép sử dụng trực tiếp xác suất đầu ra làm
căn cứ phân vùng quyết định.

Hình~\ref{fig:roc-curves} và Hình~\ref{fig:calibration} minh họa khả năng phân
biệt và mức độ hiệu chuẩn của các mô hình.

\begin{figure}[!htbp]
\centering
\cffig[0.95\textwidth]{roc_curves}
\caption{Đường cong ROC của các mô hình trên tập kiểm thử}
\label{fig:roc-curves}
\end{figure}

\begin{figure}[!htbp]
\centering
\cffig[0.95\textwidth]{calibration}
\caption{Đường cong hiệu chuẩn xác suất sau khi hiệu chuẩn đẳng hướng}
\label{fig:calibration}
\end{figure}

\subsubsection{So sánh với trần Bayes}

Vì quá trình sinh dữ liệu xác định tường minh xác suất vỡ nợ thật của từng hồ sơ,
có thể tính được trần Bayes -- giới hạn lý thuyết của mọi mô hình trên tập dữ
liệu này. Trần Bayes trên tập kiểm thử là \BayesAUC{}.

Mô hình tốt nhất đạt \BestAUC{}, tức còn cách trần Bayes \GapBayes{} điểm
$\AUC$. Khoảng cách này có ý nghĩa quan trọng về mặt diễn giải: nó cho thấy mô
hình đã khai thác được phần lớn thông tin có trong dữ liệu, và phần còn thiếu
không phải do thuật toán yếu mà do bản thân dữ liệu quan sát được chỉ là phép đo
có nhiễu của mức tín nhiệm thật. Nói cách khác, việc đầu tư thêm vào thuật toán
phức tạp hơn sẽ cho lợi ích giảm dần; cải thiện đáng kể hơn chỉ đến từ việc bổ
sung nguồn dữ liệu mới mang thông tin bổ sung.

\subsection{Đánh giá khả năng phân tách rủi ro}
\label{subsec:kq-phan-tach}

Khả năng phân tách rủi ro được đánh giá qua ba công cụ bổ trợ nhau.

\paragraph{Thống kê KS.} Mô hình được chọn đạt $\KS = \BestKS{}$, nghĩa là tại
điểm cắt tối ưu, khoảng cách giữa phân phối tích lũy của nhóm tốt và nhóm xấu đạt
\KSPoints{} điểm phần trăm. Theo thực hành ngành, giá trị $\KS$ trên 0,40 được
coi là mô hình có khả năng phân tách tốt.

\begin{figure}[!htbp]
\centering
\cffig[0.9\textwidth]{ks_curve}
\caption{Đường cong Kolmogorov--Smirnov của mô hình được chọn triển khai}
\label{fig:ks-curve}
\end{figure}

\paragraph{Phân tích theo thập phân vị.} Bảng~\ref{tab:gains} trình bày phân bố
hồ sơ xấu theo mười nhóm điểm rủi ro.

\begin{table}[!htbp]
\centering
\caption{Phân tích theo thập phân vị điểm rủi ro}
\label{tab:gains}
\small
\setlength{\tabcolsep}{3.0pt}
\begin{tabular}{@{}R{1.113cm}R{1.370cm}R{1.455cm}R{1.712cm}R{1.712cm}R{1.798cm}R{1.798cm}R{1.113cm}R{1.541cm}@{}}
\toprule
\textbf{Nhóm} & \textbf{Số hồ sơ} & \textbf{Số hồ sơ xấu} & \textbf{PD trung bình} & \textbf{Tỷ lệ nợ xấu (\%)} & \textbf{Tỷ trọng nợ xấu (\%)} & \textbf{Nợ xấu tích luỹ (\%)} & \textbf{Lift} & \textbf{Lift tích luỹ} \\
\midrule
1 & 900 & 647 & 0,7130 & 71,89 & 40,67 & 40,67 & 4,07 & 4,07 \\
2 & 900 & 362 & 0,3976 & 40,22 & 22,75 & 63,42 & 2,28 & 3,17 \\
3 & 900 & 213 & 0,2196 & 23,67 & 13,39 & 76,81 & 1,34 & 2,56 \\
4 & 900 & 128 & 0,1556 & 14,22 & 8,05 & 84,85 & 0,80 & 2,12 \\
5 & 900 & 85 & 0,0975 & 9,44 & 5,34 & 90,19 & 0,53 & 1,80 \\
6 & 900 & 59 & 0,0663 & 6,56 & 3,71 & 93,90 & 0,37 & 1,57 \\
7 & 900 & 43 & 0,0482 & 4,78 & 2,70 & 96,61 & 0,27 & 1,38 \\
8 & 900 & 21 & 0,0258 & 2,33 & 1,32 & 97,93 & 0,13 & 1,22 \\
9 & 900 & 18 & 0,0119 & 2 & 1,13 & 99,06 & 0,11 & 1,10 \\
10 & 900 & 15 & 0,0091 & 1,67 & 0,94 & 100 & 0,09 & 1 \\
\bottomrule
\end{tabular}

\vspace{2pt}
{\footnotesize\itshape Nhóm 1 gồm 10\% hồ sơ có điểm rủi ro cao nhất. Lift là tỷ số giữa mật độ nợ xấu của nhóm và mật độ nợ xấu trung bình.}
\end{table}


Kết quả cho thấy thập phân vị rủi ro cao nhất tập trung \TopDecileBad{} tổng số
hồ sơ xấu -- gấp hơn bốn lần mức phân bố ngẫu nhiên. Ba thập phân vị rủi ro cao
nhất chứa \TopThreeDecileBad{} tổng số hồ sơ xấu. Đây là thước đo trực tiếp giá trị vận hành của mô hình:
nếu tổ chức tín dụng chỉ có nguồn lực để thẩm định kỹ 30\% số hồ sơ, việc ưu tiên
theo điểm rủi ro của mô hình cho phép bao phủ hơn ba phần tư số hồ sơ có vấn đề.

\begin{figure}[!htbp]
\centering
\cffig[0.95\textwidth]{gains_lift}
\caption{Đường cong gains và lift theo thập phân vị điểm rủi ro}
\label{fig:gains}
\end{figure}

\paragraph{Phân bố điểm rủi ro.} Hình~\ref{fig:score-distribution} cho thấy phân
bố điểm rủi ro của hai nhóm hồ sơ và vị trí của các ngưỡng quyết định.

\begin{figure}[!htbp]
\centering
\cffig[0.95\textwidth]{score_distribution}
\caption{Phân bố xác suất vỡ nợ dự báo theo nhóm hồ sơ thực tế}
\label{fig:score-distribution}
\end{figure}

\subsection{So sánh hiệu năng giữa các thuật toán}
\label{subsec:so-sanh-thuat-toan}

\subsubsection{Chênh lệch giữa các thuật toán là nhỏ}

Quan sát nổi bật nhất từ Bảng~\ref{tab:model-quality} là \textbf{khoảng cách giữa
các mô hình rất hẹp}. Chênh lệch $\AUC$ giữa mô hình được chọn triển khai
(\BestModel{}) và hồi quy logistic chỉ là \GapBestLR{} -- chưa tới nửa điểm phần
trăm. Trong khi đó, hồi quy logistic có ba ưu thế rõ rệt: thời gian huấn luyện
ngắn hơn hàng chục lần, mỗi hệ số có thể diễn giải trực tiếp, và xác suất đầu ra
vốn đã được hiệu chuẩn tốt mà không cần bước hiệu chuẩn bổ sung.

Kết quả này phù hợp với phát hiện của Lessmann và cộng sự~\cite{lessmann2015}
trên các bộ dữ liệu tín dụng thực: các mô hình học tập hợp cho kết quả tốt hơn
nhưng mức chênh lệch so với hồi quy logistic không lớn.

\subsubsection{Giới hạn độ phân giải của phép đo}
\label{subsec:do-phan-giai}

Khoảng cách hẹp giữa các mô hình đặt ra một câu hỏi phương pháp mà phần lớn
nghiên cứu so sánh thuật toán không nêu: \textit{chênh lệch $\AUC$ nhỏ tới đâu
thì còn mang thông tin?}

Để trả lời, đề án chạy lại toàn bộ bước huấn luyện trên cùng tập dữ liệu, cùng
siêu tham số và cùng hạt giống ngẫu nhiên. Kết quả tái lập chính xác tới bốn chữ
số thập phân ở cả \NModels{} mô hình. Tuy nhiên khi chạy lại \textit{bước tinh
chỉnh siêu tham số}, giá trị $\AUC$ kiểm định chéo đo được dao động khoảng
$3 \times 10^{-4}$ giữa các lần chạy. Nguyên nhân mang tính kỹ thuật thuần túy:
các thư viện boosting dựng biểu đồ tần suất song song nhiều luồng, nên thứ tự
cộng dồn số thực dấu phẩy động thay đổi theo cách hệ điều hành cấp phát luồng.

Mức dao động này nhỏ, nhưng không nhỏ hơn khoảng cách giữa các mô hình dẫn đầu.
Hệ quả trực tiếp đã được quan sát trong quá trình thực nghiệm: hai bộ siêu tham
số của XGBoost có $\AUC$ kiểm định chéo cách nhau $1{,}6 \times 10^{-4}$ đã đổi
chỗ cho nhau giữa hai lần chạy, kéo theo thứ hạng của chính mô hình đó trên tập
kiểm thử cũng thay đổi.

Kết luận rút ra có ý nghĩa với mọi nghiên cứu so sánh thuật toán trên bài toán
này: \textbf{thứ hạng ở chữ số thập phân thứ tư của $\AUC$ không phải là bằng
chứng về sự vượt trội}. Các mô hình chênh nhau dưới ngưỡng đó cần được coi là
tương đương, và việc lựa chọn phải dựa trên các tiêu chí khác -- chi phí tính
toán, khả năng giải thích, độ ổn định -- chứ không dựa trên bảng xếp hạng.

\subsubsection{Đánh đổi giữa chất lượng và chi phí tính toán}

Bảng~\ref{tab:model-cost} trình bày chi phí huấn luyện và thời gian chấm điểm một
hồ sơ của từng mô hình.

\begin{table}[!htbp]
\centering
\caption{Chi phí huấn luyện và thời gian chấm điểm một hồ sơ}
\label{tab:model-cost}
\small
\setlength{\tabcolsep}{3.0pt}
\begin{tabular}{@{}L{3.799cm}L{2.849cm}R{1.899cm}R{1.994cm}R{1.994cm}R{1.709cm}@{}}
\toprule
\textbf{Mô hình} & \textbf{Nhóm} & \textbf{Thời gian HL (s)} & \textbf{Chấm điểm p50 (ms)} & \textbf{Chấm điểm p95 (ms)} & \textbf{Số vòng lặp} \\
\midrule
Kết hợp xếp tầng (Stacking) & Kết hợp xếp tầng & 214,8 & 141,38 & 202,39 & -- \\
CatBoost & Học tập hợp tuần tự & 91,5 & 4,61 & 10,46 & 353 \\
XGBoost & Học tập hợp tuần tự & 20,7 & 16,78 & 18,53 & 706 \\
LightGBM & Học tập hợp tuần tự & 10,8 & 7,85 & 8,81 & 156 \\
Hồi quy logistic & Tuyến tính & 3,2 & 10,63 & 16,55 & 94 \\
Gradient Boosting (histogram) & Học tập hợp tuần tự & 50,5 & 13,14 & 23,23 & 184 \\
Thẻ điểm WOE + Hồi quy logistic & Tuyến tính & 7,6 & 11,12 & 22,77 & 74 \\
Mạng nơ-ron truyền thẳng (MLP) & Mạng nơ-ron & 24,1 & 6,42 & 13,44 & 35 \\
Máy vectơ hỗ trợ (SVM-RBF) & Máy vectơ hỗ trợ & 16,1 & 15,10 & 20,64 & -- \\
Rừng ngẫu nhiên & Học tập hợp song song & 24 & 188,81 & 236,25 & -- \\
Cây cực ngẫu nhiên (Extra Trees) & Học tập hợp song song & 11,8 & 383,00 & 516,63 & -- \\
Cây quyết định (CART) & Cây quyết định & 5,5 & 12,06 & 14,93 & -- \\
\bottomrule
\end{tabular}
\end{table}


Hình~\ref{fig:auc-vs-latency} biểu diễn quan hệ giữa chất lượng dự báo và thời
gian chấm điểm.

\begin{figure}[!htbp]
\centering
\cffig[0.92\textwidth]{auc_vs_latency}
\caption{Đánh đổi giữa chất lượng dự báo và thời gian chấm điểm một hồ sơ}
\label{fig:auc-vs-latency}
\end{figure}

Đồ thị cho thấy một số mô hình nằm rõ rệt phía dưới đường biên hiệu quả. Rõ nhất
là mô hình kết hợp xếp tầng: nó đổi \GapTopDeploy{} điểm $\AUC$ lấy thời gian
chấm điểm cao gấp \LatRatioTopDeploy{} lần (\TopAUCLatPFive{} ms so với
\BestLatPFive{} ms ở phân vị 95). Extra Trees và rừng ngẫu nhiên vừa chậm hơn vừa
kém chính xác hơn nhóm boosting. Với bài toán kiểm duyệt trực tuyến -- nơi kết
quả phải trả về trong lúc cán bộ tín dụng đang làm việc với hồ sơ -- những mô
hình này không phải lựa chọn hợp lý dù chỉ xét riêng góc độ kỹ thuật.

\subsubsection{Tiêu chí chọn mô hình triển khai}
\label{subsec:kq-chon-mo-hinh}

Hai phân tích trên dẫn tới một quy tắc lựa chọn cụ thể, thay cho cách lấy thẳng
mô hình đứng đầu bảng xếp hạng $\AUC$. Quy tắc gồm hai bước:

\begin{enumerate}[label=(\arabic*)]
  \item \textbf{Xác định nhóm tương đương về chất lượng:} mọi mô hình có $\AUC$
  nằm trong khoảng \TieTolerance{} so với mô hình dẫn đầu. Ngưỡng này được chọn
  lớn hơn một bậc so với dao động đo đạc đã xác định ở
  mục~\ref{subsec:do-phan-giai}, đồng thời vẫn nhỏ hơn mức chênh lệch có ý nghĩa
  nghiệp vụ.
  \item \textbf{Trong nhóm đó, chọn mô hình có thời gian chấm điểm thấp nhất} --
  tức mô hình nằm trên biên Pareto của Hình~\ref{fig:auc-vs-latency}.
\end{enumerate}

Áp dụng trên kết quả thực nghiệm, nhóm tương đương gồm \NTiedModels{} mô hình:
kết hợp xếp tầng (\TopAUC{}), CatBoost (\CatAUC{}) và XGBoost (\XGBAUC{}). Thời
gian chấm điểm ở phân vị 95 của ba mô hình lần lượt là \StackLat{} ms,
\CatLat{} ms và \XGBLat{} ms. Mô hình được chọn triển khai là \BestModel{}.

Lựa chọn này có hai căn cứ. Thứ nhất, nếu lấy mô hình kết hợp xếp tầng, thời gian
xử lý một hồ sơ sẽ vượt mục tiêu 100 ms đặt ra tại
mục~\ref{subsec:yeu-cau} -- nghĩa là đánh đổi một yêu cầu đã công bố lấy
\GapTopDeploy{} điểm $\AUC$ nằm trong vùng nhiễu đo. Thứ hai, quy tắc này cho kết
quả \textit{ổn định}: khi thứ hạng trong nhóm dẫn đầu đảo chiều giữa các lần chạy,
mô hình được chọn không đổi, vì tiêu chí quyết định là thời gian chấm điểm --
đại lượng chênh nhau hàng chục lần chứ không phải vài phần vạn.

\subsubsection{Tác động nghiệp vụ}

Bảng~\ref{tab:model-business} quy đổi chất lượng thống kê thành các chỉ tiêu
nghiệp vụ tại ngưỡng tối ưu chi phí.

\begin{table}[!htbp]
\centering
\caption{Tác động nghiệp vụ của các mô hình tại ngưỡng tối ưu chi phí}
\label{tab:model-business}
\small
\setlength{\tabcolsep}{3.0pt}
\begin{tabular}{@{}L{2.718cm}R{1.500cm}R{1.521cm}R{1.565cm}R{1.500cm}R{1.750cm}R{1.750cm}R{1.521cm}@{}}
\toprule
\textbf{Mô hình} & \textbf{Ngưỡng} & \textbf{Tỷ lệ duyệt} & \textbf{Nợ xấu sau duyệt} & \textbf{Chi phí} & \textbf{Accuracy} & \textbf{Precision} & \textbf{Recall} \\
\midrule
\textbf{Kết hợp xếp tầng (Stacking)} & \textbf{0,0501} & \textbf{44,0\%} & \textbf{2,98\%} & \textbf{0,5532} & \textbf{0,5910} & \textbf{0,2925} & \textbf{0,9258} \\
CatBoost & 0,0746 & 50,5\% & 3,52\% & 0,5497 & 0,6459 & 0,3210 & 0,8994 \\
XGBoost & 0,0789 & 49,2\% & 3,32\% & 0,5431 & 0,6366 & 0,3161 & 0,9076 \\
LightGBM & 0,0758 & 45,2\% & 3,19\% & 0,5586 & 0,6003 & 0,2965 & 0,9183 \\
Hồi quy logistic & 0,0729 & 43,4\% & 2,94\% & 0,5550 & 0,5856 & 0,2899 & 0,9277 \\
Gradient Boosting (histogram) & 0,0749 & 44,7\% & 3,21\% & 0,5624 & 0,5952 & 0,2938 & 0,9189 \\
Thẻ điểm WOE + Hồi quy logistic & 0,0848 & 48,0\% & 3,26\% & 0,5464 & 0,6259 & 0,3101 & 0,9114 \\
Mạng nơ-ron truyền thẳng (MLP) & 0,0778 & 48,7\% & 3,40\% & 0,5510 & 0,6311 & 0,3126 & 0,9063 \\
Máy vectơ hỗ trợ (SVM-RBF) & 0,0702 & 46,1\% & 3,28\% & 0,5588 & 0,6074 & 0,2999 & 0,9145 \\
Rừng ngẫu nhiên & 0,0682 & 49,4\% & 3,62\% & 0,5618 & 0,6350 & 0,3140 & 0,8988 \\
Cây cực ngẫu nhiên (Extra Trees) & 0,0866 & 42,8\% & 3,17\% & 0,5714 & 0,5777 & 0,2854 & 0,9233 \\
Cây quyết định (CART) & 0,0707 & 38,1\% & 4,08\% & 0,6444 & 0,5267 & 0,2605 & 0,9120 \\
\bottomrule
\end{tabular}

\vspace{2pt}
{\footnotesize\itshape Ngưỡng được chọn theo hàm chi phí với tỷ lệ $c_{FN}:c_{FP} = 12:1$ trên tập kiểm định.}
\end{table}


\subsubsection{Mức độ quá khớp và độ ổn định theo thời gian}
\label{subsec:kq-on-dinh}

\paragraph{Mức độ quá khớp.} Bảng~\ref{tab:overfit} đối chiếu $\AUC$ trên tập
kiểm định với $\AUC$ trên tập kiểm thử. Chênh lệch trung bình giữa hai tập là
\MeanOverfitGap{} và mức lớn nhất là \MaxOverfitGap{} -- đều rất nhỏ so với bản
thân giá trị $\AUC$. Kết quả này cho thấy bốn cơ chế kiểm soát quá khớp trình bày
ở mục~\ref{subsec:kiem-soat-qua-khop} đã phát huy tác dụng: không mô hình nào phụ
thuộc quá mức vào đặc thù của giai đoạn huấn luyện.

\begin{table}[!htbp]
\centering
\caption{Mức độ quá khớp đo bằng chênh lệch AUC giữa hai tập dữ liệu}
\label{tab:overfit}
\small
\setlength{\tabcolsep}{5.0pt}
\begin{tabular}{@{}L{5.4cm}R{2.9cm}R{2.9cm}R{2.6cm}@{}}
\toprule
\textbf{Mô hình} & \textbf{AUC (kiểm định)} & \textbf{AUC (kiểm thử)} & \textbf{Chênh lệch} \\
\midrule
Kết hợp xếp tầng (Stacking) & 0,8639 & 0,8620 & 0,0019 \\
CatBoost & 0,8633 & 0,8614 & 0,0019 \\
XGBoost & 0,8646 & 0,8612 & 0,0034 \\
LightGBM & 0,8595 & 0,8590 & 0,0005 \\
Hồi quy logistic & 0,8590 & 0,8579 & 0,0011 \\
Gradient Boosting (histogram) & 0,8616 & 0,8575 & 0,0041 \\
Thẻ điểm WOE + Hồi quy logistic & 0,8610 & 0,8570 & 0,0040 \\
Mạng nơ-ron truyền thẳng (MLP) & 0,8554 & 0,8560 & -0,0006 \\
Máy vectơ hỗ trợ (SVM-RBF) & 0,8578 & 0,8540 & 0,0038 \\
Rừng ngẫu nhiên & 0,8565 & 0,8527 & 0,0038 \\
Cây cực ngẫu nhiên (Extra Trees) & 0,8521 & 0,8471 & 0,0050 \\
Cây quyết định (CART) & 0,8231 & 0,8206 & 0,0025 \\
\bottomrule
\end{tabular}

\vspace{2pt}
{\footnotesize\itshape Chênh lệch dương lớn cho thấy mô hình phụ thuộc quá mức vào đặc thù của giai đoạn huấn luyện.}
\end{table}


\paragraph{Trôi dạt phân phối dữ liệu.} Bảng~\ref{tab:drift} trình bày mười đặc
trưng có mức trôi dạt lớn nhất. Toàn bộ tập đặc trưng có $\PSI$ dưới ngưỡng cảnh
báo, với giá trị lớn nhất là \MaxPSI{}; số đặc trưng ở mức cảnh báo là
\NPSIWarn{} và ở mức báo động là \NPSIAlert{}.

\begin{table}[!htbp]
\centering
\caption{Mười đặc trưng có mức trôi dạt phân phối lớn nhất}
\label{tab:drift}
\small
\setlength{\tabcolsep}{5.0pt}
\begin{tabular}{@{}L{6.4cm}R{2.6cm}L{4.4cm}@{}}
\toprule
\textbf{Đặc trưng} & \textbf{PSI} & \textbf{Trạng thái} \\
\midrule
salary\_\allowbreak cnt\_\allowbreak 90d & 0,0598 & ỔN ĐỊNH \\
salary\_\allowbreak cnt\_\allowbreak 30d & 0,0571 & ỔN ĐỊNH \\
income\_\allowbreak stability & 0,0533 & ỔN ĐỊNH \\
salary\_\allowbreak cnt\_\allowbreak 180d & 0,0504 & ỔN ĐỊNH \\
income\_\allowbreak verification\_\allowbreak ratio & 0,0395 & ỔN ĐỊNH \\
outflow\_\allowbreak ratio\_\allowbreak 30d & 0,0388 & ỔN ĐỊNH \\
salary\_\allowbreak amt\_\allowbreak 30d & 0,0374 & ỔN ĐỊNH \\
net\_\allowbreak flow\_\allowbreak 30d & 0,0363 & ỔN ĐỊNH \\
txn\_\allowbreak amt\_\allowbreak std\_\allowbreak 30d & 0,0336 & ỔN ĐỊNH \\
inflow\_\allowbreak to\_\allowbreak installment & 0,0291 & ỔN ĐỊNH \\
\bottomrule
\end{tabular}

\vspace{2pt}
{\footnotesize\itshape Ngưỡng cảnh báo PSI $= 0{,}10$; ngưỡng báo động PSI $= 0{,}25$.}
\end{table}


\paragraph{Chất lượng dự báo theo thời gian.} Bảng~\ref{tab:perf-time} theo dõi
chất lượng dự báo theo từng tháng của giai đoạn kiểm thử. Giá trị $\AUC$ dao động
trong khoảng hẹp quanh mức 0,86 mà không có xu hướng giảm rõ rệt, và chênh lệch
giữa tỷ lệ nợ xấu thực tế với mức $\PD$ trung bình dự báo luôn dưới một điểm phần
trăm. Trong phạm vi bốn tháng quan sát, mô hình giữ được cả khả năng phân biệt
lẫn mức hiệu chuẩn. Tuy nhiên, bốn tháng là khoảng thời gian ngắn so với chu kỳ
suy giảm điển hình của mô hình tín dụng trong thực tế, nên kết luận về độ ổn định
chỉ có giá trị trong phạm vi dữ liệu thực nghiệm.

\begin{table}[!htbp]
\centering
\caption{Chất lượng dự báo của mô hình theo từng tháng của giai đoạn kiểm thử}
\label{tab:perf-time}
\small
\setlength{\tabcolsep}{3.0pt}
\begin{tabular}{@{}L{1.671cm}R{1.495cm}R{2.023cm}R{2.287cm}R{1.847cm}R{1.500cm}R{1.500cm}R{1.500cm}@{}}
\toprule
\textbf{Tháng} & \textbf{Số hồ sơ} & \textbf{Nợ xấu thực tế (\%)} & \textbf{PD dự báo trung bình (\%)} & \textbf{Chênh lệch (đpt)} & \textbf{AUC} & \textbf{KS} & \textbf{Gini} \\
\midrule
2026-03 & 1.652 & 17,86 & 17,76 & -0,09 & 0,8665 & 0,5953 & 0,7330 \\
2026-04 & 2.485 & 17,46 & 17,13 & -0,33 & 0,8584 & 0,5730 & 0,7168 \\
2026-05 & 2.429 & 18,53 & 17,74 & -0,79 & 0,8551 & 0,5673 & 0,7102 \\
2026-06 & 2.434 & 16,93 & 17,26 & 0,33 & 0,8679 & 0,5782 & 0,7358 \\
\bottomrule
\end{tabular}

\vspace{2pt}
{\footnotesize\itshape Chênh lệch tính bằng điểm phần trăm giữa mức PD dự báo trung bình và tỷ lệ nợ xấu thực tế của tháng.}
\end{table}


\begin{figure}[!htbp]
\centering
\cffig[0.95\textwidth]{performance_over_time}
\caption{Diễn biến chất lượng dự báo theo tháng trên giai đoạn kiểm thử}
\label{fig:perf-time}
\end{figure}

\section{Mô phỏng và giải thích mô hình}
\label{sec:mo-phong-he-thong}

\subsection{Thử nghiệm giải thích quyết định mô hình}
\label{subsec:kq-giai-thich}

\subsubsection{Giải thích ở cấp độ toàn cục}

Phân tích SHAP được thực hiện trên mẫu đại diện của tập kiểm thử.
Bảng~\ref{tab:shap-group} tổng hợp đóng góp theo nhóm nguồn dữ liệu.

\begin{table}[!htbp]
\centering
\caption{Tỷ trọng đóng góp vào quyết định của mô hình theo nhóm đặc trưng}
\label{tab:shap-group}
\small
\setlength{\tabcolsep}{3.0pt}
\begin{tabular}{@{}L{3.648cm}R{1.911cm}R{1.911cm}R{2.258cm}R{2.085cm}R{2.432cm}@{}}
\toprule
\textbf{Nhóm đặc trưng} & \textbf{Số đặc trưng} & \textbf{Tổng |SHAP|} & \textbf{|SHAP| trung bình} & \textbf{|SHAP| lớn nhất} & \textbf{Tỷ trọng đóng góp (\%)} \\
\midrule
Hành vi giao dịch & 76 & 1,2895 & 0,0170 & 0,1623 & 34,00 \\
Lịch sử tín dụng (CIC) & 14 & 1,1692 & 0,0835 & 0,4217 & 30,83 \\
Biến tổ hợp & 12 & 0,9091 & 0,0758 & 0,1899 & 23,97 \\
Thông tin cá nhân & 8 & 0,2662 & 0,0333 & 0,0909 & 7,02 \\
Tham số khoản vay & 6 & 0,1587 & 0,0264 & 0,0432 & 4,18 \\
\bottomrule
\end{tabular}

\vspace{2pt}
{\footnotesize\itshape Tính bằng tổng giá trị SHAP tuyệt đối trung bình của các đặc trưng trong từng nhóm, chuẩn hóa về tổng bằng 100\%.}
\end{table}


Kết quả cho thấy nhóm đặc trưng hành vi giao dịch đóng góp
\ShapBehaviourPct{}\% vào quyết định của mô hình -- cao hơn nhóm CIC
(\ShapCICPct{}\%) vốn được coi là nguồn thông tin truyền thống quan trọng nhất.
Nhóm biến tổ hợp đóng góp \ShapDerivedPct{}\%, còn hai nhóm thông tin cá nhân và
tham số khoản vay chỉ đóng góp lần lượt \ShapStaticPct{}\% và \ShapLoanPct{}\%.

Phát hiện này có ý nghĩa thực tiễn: dữ liệu giao dịch là nguồn mà tổ chức tín
dụng tự sở hữu, không phát sinh chi phí truy vấn bên ngoài và luôn có sẵn kể cả
với nhóm khách hàng thiếu lịch sử tín dụng. Việc nhóm này đóng góp phần lớn nhất
vào quyết định của mô hình là căn cứ định lượng cho khuyến nghị khai thác dữ liệu
giao dịch trong công tác kiểm duyệt tín dụng.

Kết quả trên cần được đặt cạnh phân tích $\IV$ ở mục~\ref{subsec:kq-dac-trung} để
tránh diễn giải sai. Xét từng đặc trưng riêng lẻ, biến CIC mạnh hơn hẳn biến giao
dịch. Nhưng trong mô hình hoàn chỉnh, tổng đóng góp của nhóm giao dịch lại lớn
hơn, vì nhóm này gồm nhiều đặc trưng bổ sung thông tin cho nhau ở những khía cạnh
khác nhau của hành vi tài chính.

\begin{figure}[!htbp]
\centering
\cffig[0.95\textwidth]{shap_summary}
\caption{Biểu đồ tóm tắt SHAP của hai mươi đặc trưng quan trọng nhất}
\label{fig:shap-summary}
\end{figure}

\begin{figure}[!htbp]
\centering
\cffig[0.9\textwidth]{shap_group_share}
\caption{Tỷ trọng đóng góp vào quyết định theo nhóm nguồn dữ liệu}
\label{fig:shap-group}
\end{figure}

\subsubsection{Kiểm chứng chéo bằng phép hoán vị}

Kết quả SHAP được kiểm chứng độc lập bằng phép hoán vị: với mỗi đặc trưng, giá
trị của nó được xáo trộn ngẫu nhiên và mức suy giảm $\AUC$ được ghi nhận. Thứ
hạng các đặc trưng theo hai phương pháp về cơ bản trùng khớp ở nhóm đầu bảng, xác
nhận rằng kết quả SHAP không phải là sản phẩm của một đặc thù kỹ thuật nào đó
trong cách tính.

\begin{figure}[!htbp]
\centering
\cffig[0.95\textwidth]{permutation_importance_bar}
\caption{Tầm quan trọng đặc trưng theo phép hoán vị}
\label{fig:permutation}
\end{figure}

\subsubsection{Giải thích ở cấp độ từng hồ sơ}

Với mỗi hồ sơ, mô hình sinh danh sách lý do gồm ba thành phần: luật nghiệp vụ bị
vi phạm, chỉ tiêu định lượng then chốt kèm giá trị thực tế, và các đặc trưng có
đóng góp SHAP lớn nhất. Danh sách này được trả về cùng kết quả chấm điểm, cho
phép cán bộ tín dụng đánh giá tính hợp lý của khuyến nghị thay vì phải chấp nhận
một con số không có căn cứ.

\subsection{Mô phỏng xác định ngưỡng phê duyệt tín dụng}
\label{subsec:kq-nguong}

\subsubsection{So sánh ba tiêu chí xác định điểm cắt}

Bảng~\ref{tab:cutoff} trình bày kết quả của ba tiêu chí xác định điểm cắt, tính
trên tập kiểm định.

\begin{table}[!htbp]
\centering
\caption{So sánh ba tiêu chí xác định điểm cắt phê duyệt}
\label{tab:cutoff}
\small
\setlength{\tabcolsep}{3.0pt}
\begin{tabular}{@{}L{4.255cm}R{1.850cm}R{1.850cm}R{2.220cm}R{2.220cm}R{1.850cm}@{}}
\toprule
\textbf{Tiêu chí} & \textbf{Điểm cắt PD} & \textbf{Tỷ lệ duyệt (\%)} & \textbf{Nợ xấu nhóm duyệt (\%)} & \textbf{Nợ xấu nhóm từ chối (\%)} & \textbf{Chi phí kỳ vọng} \\
\midrule
Tối thiểu chi phí (c\_\allowbreak FN:c\_\allowbreak FP = 12:1) & 0,07 & 51,64 & 3,33 & 31,23 & 0,54 \\
Tỷ lệ duyệt mục tiêu 70\% & 0,17 & 72,30 & 6,01 & 45,05 & 0,67 \\
Trần nợ xấu nhóm duyệt 5\% & 0,15 & 65,30 & 4,92 & 39,23 & 0,60 \\
\bottomrule
\end{tabular}
\end{table}


Ba tiêu chí cho ba điểm cắt khác nhau rõ rệt, phản ánh ba cách đặt bài toán khác
nhau. Tiêu chí tối thiểu chi phí cho điểm cắt thấp nhất (\CostOptimalCut{}) và do
đó thận trọng nhất, với tỷ lệ nợ xấu trong nhóm được duyệt chỉ
\CutoffCostBadRate{}. Tiêu chí tỷ lệ duyệt mục tiêu 70\% cho điểm cắt cao nhất,
đổi lại tỷ lệ nợ xấu nhóm duyệt tăng lên \CutoffRateBadRate{}. Tiêu chí trần nợ
xấu 5\% nằm ở giữa, với tỷ lệ nợ xấu nhóm duyệt \CutoffCapBadRate{}.

Việc trình bày song song ba tiêu chí -- thay vì áp đặt một ngưỡng duy nhất -- là
có chủ đích: lựa chọn giữa chúng là \textit{quyết định chính sách} thuộc thẩm
quyền của tổ chức tín dụng, không phải quyết định kỹ thuật của người xây dựng mô
hình.

\subsubsection{Độ nhạy của kết quả theo điểm cắt}

Bảng~\ref{tab:cutoff-sens} và Hình~\ref{fig:cutoff-sens} định lượng đánh đổi giữa
quy mô cấp tín dụng và chất lượng danh mục.

\begin{table}[!htbp]
\centering
\caption{Độ nhạy của kết quả kiểm duyệt theo điểm cắt}
\label{tab:cutoff-sens}
\small
\setlength{\tabcolsep}{3.0pt}
\begin{tabular}{@{}R{2.0cm}R{2.2cm}R{2.2cm}R{2.6cm}R{2.4cm}R{2.2cm}@{}}
\toprule
\textbf{Điểm cắt PD} & \textbf{Số hồ sơ duyệt} & \textbf{Tỷ lệ duyệt (\%)} & \textbf{Nợ xấu nhóm duyệt (\%)} & \textbf{Nợ xấu bị chặn (\%)} & \textbf{Chi phí kỳ vọng} \\
\midrule
0.02 & 2.085 & 23,17 & 1,10 & 21,56 & 0,63 \\
0.04 & 2.775 & 30,83 & 1,66 & 23,58 & 0,59 \\
0.06 & 3.505 & 38,94 & 2,25 & 26,11 & 0,56 \\
0.08 & 4.893 & 54,37 & 3,56 & 32,63 & 0,54 \\
0.1 & 4.893 & 54,37 & 3,56 & 32,63 & 0,54 \\
0.15 & 5.877 & 65,30 & 4,92 & 39,23 & 0,60 \\
0.2 & 6.843 & 76,03 & 6,62 & 49,19 & 0,73 \\
0.25 & 6.999 & 77,77 & 6,94 & 51,37 & 0,76 \\
0.32 & 7.533 & 83,70 & 8,56 & 59,24 & 0,93 \\
0.4 & 7.623 & 84,70 & 8,87 & 60,86 & 0,96 \\
\bottomrule
\end{tabular}
\end{table}


\begin{figure}[!htbp]
\centering
\cffig[0.9\textwidth]{cutoff_sensitivity}
\caption{Đánh đổi giữa tỷ lệ duyệt và tỷ lệ nợ xấu theo điểm cắt}
\label{fig:cutoff-sens}
\end{figure}

Đường cong này là công cụ đối thoại giữa bộ phận mô hình và bộ phận kinh doanh:
thay vì tranh luận về một giá trị ngưỡng trừu tượng, mỗi lựa chọn được quy đổi
thành hệ quả cụ thể về tỷ lệ duyệt và tỷ lệ nợ xấu.

\subsubsection{Hai ngưỡng cho ba vùng quyết định}

Từ điểm cắt tối ưu chi phí, hai ngưỡng phân vùng được xác định:
$\tau_{\text{duyệt}} = \PDApproveMax{}$ và
$\tau_{\text{từ chối}} = \PDRejectMin{}$. Kết quả áp dụng trên tập kiểm thử được
trình bày tại Bảng~\ref{tab:decision-mix}.

\begin{table}[!htbp]
\centering
\caption{Cơ cấu khuyến nghị của mô hình và tỷ lệ nợ xấu thực tế theo từng nhóm}
\label{tab:decision-mix}
\small
\setlength{\tabcolsep}{5.0pt}
\begin{tabular}{@{}L{5.854cm}R{2.537cm}R{2.537cm}R{3.317cm}@{}}
\toprule
\textbf{Kết quả} & \textbf{Số hồ sơ} & \textbf{Tỷ trọng (\%)} & \textbf{Tỷ lệ nợ xấu thực tế (\%)} \\
\midrule
Khuyến nghị duyệt & 1.459 & 29,18 & 3,43 \\
Chuyển thẩm định thủ công & 1.637 & 32,74 & 10,93 \\
Khuyến nghị từ chối & 1.904 & 38,08 & 34,45 \\
\bottomrule
\end{tabular}
\end{table}


Kết quả xác nhận mô hình hoạt động đúng thiết kế: tỷ lệ nợ xấu thực tế tăng đơn
điệu theo ba nhóm khuyến nghị, lần lượt là \BadApprove{} ở nhóm khuyến nghị
duyệt, \BadReview{} ở nhóm chuyển thẩm định và \BadReject{} ở nhóm khuyến nghị từ
chối. Mức chênh lệch hơn mười lần giữa nhóm đầu và nhóm cuối cho thấy khả năng
phân tách thực tế của mô hình.

Tỷ lệ hồ sơ được xử lý tự động đạt \AutomationRate{}, nghĩa là khoảng một phần ba
số hồ sơ (\ReviewRate{}) vẫn cần cán bộ tín dụng xem xét. Đây là kết quả phù hợp
với định hướng \textit{hỗ trợ} của đề án: những hồ sơ mà mô hình không đủ chắc
chắn được chuyển cho con người thay vì bị quyết định máy móc.

\begin{figure}[!htbp]
\centering
\cffig[0.98\textwidth]{decision_mix}
\caption{Cơ cấu khuyến nghị và tỷ lệ nợ xấu thực tế theo từng nhóm}
\label{fig:decision-mix}
\end{figure}

\subsubsection{Tần suất vi phạm luật nghiệp vụ}

Bảng~\ref{tab:rule-hits} trình bày tần suất vi phạm của \NRules{} luật nghiệp vụ
trong quá trình mô phỏng.

\begin{table}[!htbp]
\centering
\caption{Tần suất vi phạm các luật nghiệp vụ}
\label{tab:rule-hits}
\small
\setlength{\tabcolsep}{5.0pt}
\begin{tabular}{@{}L{2.8cm}L{6.0cm}L{2.6cm}R{2.6cm}@{}}
\toprule
\textbf{Mã luật} & \textbf{Tên luật} & \textbf{Hành động} & \textbf{Tỷ lệ vi phạm} \\
\midrule
R501 & Dòng tiền vào thực tế thấp hơn nhiều so với thu nhập khai báo & REVIEW & 23,08\% \\
R202 & Thu nhập còn lại sau trả nợ không đủ chi phí sinh hoạt & HARD\_\allowbreak REJECT & 21,10\% \\
R104 & Số ngày quá hạn tối đa 12 tháng vượt 30 ngày & REVIEW & 12,98\% \\
R101 & Nhóm nợ xấu nhất 12 tháng vượt nhóm 2 & HARD\_\allowbreak REJECT & 11,74\% \\
R107 & Tỷ lệ sử dụng hạn mức trên 95\% & WARNING & 11,10\% \\
R103 & Có nợ đã xử lý rủi ro / xóa nợ & HARD\_\allowbreak REJECT & 10,18\% \\
R401 & Không lấy được dữ liệu CIC & REVIEW & 9,74\% \\
R402 & Khách hàng thiếu lịch sử tín dụng (thin-file) & WARNING & 9,74\% \\
R002 & Thu nhập dưới mức tối thiểu 6 triệu VND & HARD\_\allowbreak REJECT & 9,28\% \\
R102 & Đang có nợ nhóm 3 trở lên & HARD\_\allowbreak REJECT & 9,28\% \\
R201 & Tỷ lệ nợ trên thu nhập (DTI) vượt 60\% & HARD\_\allowbreak REJECT & 7,46\% \\
R502 & Tỷ trọng chi tiêu nhóm rủi ro cao bất thường & REVIEW & 6,04\% \\
\bottomrule
\end{tabular}
\end{table}


\begin{figure}[!htbp]
\centering
\cffig[0.95\textwidth]{rule_hits}
\caption{Tần suất vi phạm các luật nghiệp vụ}
\label{fig:rule-hits}
\end{figure}

\subsection{Thử nghiệm tích hợp API và luồng xử lý hồ sơ}
\label{subsec:kq-api}

\subsubsection{Kết quả tích hợp dịch vụ REST}

Hệ thống được đóng gói thành một dịch vụ REST với năm điểm cuối. Kết quả thử
nghiệm tích hợp trình bày tại Bảng~\ref{tab:api}.

\begin{table}[!htbp]
\centering
\caption{Kết quả thử nghiệm tích hợp dịch vụ REST}
\label{tab:api}
\small
\setlength{\tabcolsep}{5.0pt}
\begin{tabular}{@{}L{4.600cm}L{4.161cm}R{1.702cm}L{3.783cm}@{}}
\toprule
\textbf{Điểm cuối} & \textbf{Mô tả} & \textbf{Mã HTTP} & \textbf{Kết quả} \\
\midrule
GET /health & Kiểm tra trạng thái dịch vụ & 200 & Đạt \\
POST /api/v1/screen & Kiểm duyệt từng hồ sơ (50 lần gọi) & 200 & 50/50 thành công, 62.4 ms/hồ sơ \\
POST /api/v1/screen/batch & Kiểm duyệt theo lô 50 hồ sơ & 200 & 50 hồ sơ, 2245 ms (44.79 ms/hồ sơ) \\
GET /api/v1/features/\{id\} & Truy xuất vector đặc trưng đầu vào & 200 & 157 đặc trưng \\
GET /admin/rules & Danh mục luật nghiệp vụ đang áp dụng & 200 & 21 luật \\
\bottomrule
\end{tabular}
\end{table}


Toàn bộ điểm cuối hoạt động đúng. Điểm cuối kiểm duyệt theo lô cho thời gian xử
lý bình quân mỗi hồ sơ thấp hơn so với gọi từng hồ sơ riêng lẻ, phản ánh phần chi
phí cố định của mỗi lần gọi dịch vụ được phân bổ trên nhiều hồ sơ.

\subsubsection{Thời gian xử lý một hồ sơ}

Bảng~\ref{tab:processing-time} phân rã thời gian xử lý một hồ sơ theo từng chặng.

\begin{table}[!htbp]
\centering
\caption{Thời gian xử lý một hồ sơ theo từng chặng}
\label{tab:processing-time}
\small
\setlength{\tabcolsep}{5.0pt}
\begin{tabular}{@{}L{4.6cm}R{2.6cm}R{2.2cm}R{2.2cm}R{2.2cm}@{}}
\toprule
\textbf{Chặng xử lý} & \textbf{Trung bình (ms)} & \textbf{p50 (ms)} & \textbf{p95 (ms)} & \textbf{p99 (ms)} \\
\midrule
Chấm điểm mô hình & 40,174 & 34,773 & 74,507 & 98,067 \\
Áp dụng luật nghiệp vụ & 0,048 & 0,046 & 0,086 & 0,126 \\
Áp dụng chính sách & 0,053 & 0,051 & 0,088 & 0,154 \\
Tổng cộng & 40,327 & 34,928 & 74,657 & 98,372 \\
\bottomrule
\end{tabular}

\vspace{2pt}
{\footnotesize\itshape Đo trên một tiến trình đơn, không tính thời gian truyền qua mạng và thời gian truy vấn dữ liệu từ hệ thống nguồn.}
\end{table}


Kết quả cho thấy \textbf{chấm điểm mô hình chiếm gần như toàn bộ} thời gian xử
lý: trung bình \SysMean{} ms cho một hồ sơ, trong đó phần áp dụng luật nghiệp vụ
và chính sách chỉ chiếm chưa tới 0,1 ms. Phân vị 95\% của tổng thời gian là
\SysPNinetyFive{} ms và phân vị 99\% là \SysPNinetyNine{} ms, đều nằm dưới mục
tiêu 100 ms đã đặt ra ở mục~\ref{subsec:yeu-cau}.

Thời gian chấm điểm trong mô phỏng quy trình (\SysMean{} ms) cao hơn số đo ở bước
huấn luyện (\BestLatPFive{} ms tại phân vị 95\%) do hai phép đo thực hiện trong
điều kiện khác nhau: bước huấn luyện đo trên ma trận đặc trưng đã chuẩn bị sẵn,
còn mô phỏng quy trình đo toàn bộ đường xử lý từ dữ liệu đầu vào thô, bao gồm cả
việc dựng khung dữ liệu, ép kiểu và ánh xạ hạng mục. Phần chênh lệch chính là chi
phí của lớp tích hợp. Đây là kết quả có giá trị thực tiễn: khi triển khai, chi
phí chuẩn bị dữ liệu đầu vào có thể lớn hơn chi phí suy luận của bản thân mô
hình, nên việc tối ưu hiệu năng cần bắt đầu từ khâu này.

\section{Đánh giá chung}
\label{sec:danh-gia-chung}

\subsection{So sánh với quy trình thẩm định truyền thống}
\label{subsec:so-sanh-truyen-thong}

Để đánh giá giá trị thực tế của mô hình, đề án mô phỏng quy trình thẩm định
truyền thống bằng chính bộ luật nghiệp vụ đang áp dụng, không sử dụng mô hình học
máy: hồ sơ vi phạm luật loại trừ cứng bị từ chối, hồ sơ vi phạm luật cảnh báo
được chuyển thẩm định, phần còn lại được duyệt. Hai phương pháp được chạy trên
cùng một tập hồ sơ; kết quả trình bày tại Bảng~\ref{tab:traditional}.

\begin{table}[!htbp]
\centering
\caption{So sánh mô hình đề xuất với quy trình thẩm định truyền thống}
\label{tab:traditional}
\small
\setlength{\tabcolsep}{5.0pt}
\begin{tabular}{@{}L{6.312cm}R{4.142cm}R{4.142cm}@{}}
\toprule
\textbf{Chỉ tiêu} & \textbf{Truyền thống} & \textbf{Đề xuất} \\
\midrule
Tỷ lệ duyệt (\%) & 35.82 & 29.18 \\
Tỷ lệ thẩm định thủ công (\%) & 27.22 & 32.74 \\
Tỷ lệ từ chối (\%) & 36.96 & 38.08 \\
Nợ xấu nhóm duyệt (\%) & 6.14 & 3.43 \\
Nợ xấu chặn được (\%) & 33.93 & 34.45 \\
Tỷ lệ tự động hóa (\%) & 72.78 & 67.26 \\
Thời gian xử lý (ms/hồ sơ) & 0.262 & 40.327 \\
\bottomrule
\end{tabular}

\vspace{2pt}
{\footnotesize\itshape Quy trình truyền thống được mô phỏng bằng bộ luật nghiệp vụ đang áp dụng, không sử dụng mô hình học máy.}
\end{table}


\begin{figure}[!htbp]
\centering
\cffig[0.95\textwidth]{traditional_comparison}
\caption{So sánh mô hình đề xuất với quy trình thẩm định truyền thống}
\label{fig:traditional}
\end{figure}

Kết quả phản ánh một đánh đổi gồm bốn khía cạnh.

\paragraph{Về chất lượng danh mục -- cải thiện rõ rệt.} Tỷ lệ nợ xấu trong nhóm
hồ sơ được duyệt giảm từ \TradBadRate{} xuống \PropBadRate{}, tức giảm
\BadRateReduction{} điểm phần trăm, tương đương mức giảm tương đối
\BadRateReductionPct{}. Đây là chỉ tiêu có ý nghĩa kinh tế trực tiếp: với cùng
quy mô dư nợ được cấp, tổn thất tín dụng kỳ vọng giảm gần một nửa.

\paragraph{Về quy mô cấp tín dụng -- giảm.} Đổi lại, tỷ lệ hồ sơ được duyệt giảm
từ \TradApproveRate{} xuống \ApproveRate{}. Hệ thống thận trọng hơn quy trình
truyền thống: nó phát hiện và loại ra những hồ sơ mà bộ quy tắc cố định bỏ sót.

\paragraph{Về khối lượng công việc thủ công -- tăng nhẹ.} Tỷ lệ hồ sơ chuyển thẩm
định thủ công tăng từ \TradReviewRate{} lên \PropReviewRate{}. Đây là hệ quả trực
tiếp của thiết kế ba vùng: mô hình chủ động thừa nhận vùng không chắc chắn của
mình thay vì đưa ra khuyến nghị dứt khoát cho mọi hồ sơ.

\paragraph{Về thời gian xử lý.} Quy trình chỉ dùng luật nhanh hơn nhiều về mặt
tính toán (dưới 0,2 ms so với \SysMean{} ms). Tuy nhiên so sánh này không phản
ánh thực tế vận hành: trong quy trình truyền thống, phần lớn thời gian nằm ở công
đoạn cán bộ tín dụng thu thập, đối chiếu và đánh giá hồ sơ -- tính bằng phút hoặc
giờ -- chứ không phải ở việc kiểm tra bộ quy tắc. Mức chênh lệch vài chục mili
giây giữa hai phương pháp là không đáng kể trên thang thời gian đó.

\paragraph{Diễn giải tổng hợp.} Đánh đổi mà mô hình mang lại là: chấp nhận giảm
\ApproveRateDrop{} điểm phần trăm tỷ lệ duyệt và tăng \ReviewRateRise{} điểm phần
trăm khối lượng thẩm định thủ công, để đổi lấy việc giảm gần một nửa tỷ lệ nợ xấu
trong nhóm được duyệt. Mức đánh đổi này có phù hợp hay không phụ thuộc khẩu vị
rủi ro và chiến lược kinh doanh của từng tổ chức. Khác biệt cơ bản so với bộ quy
tắc cố định là mô hình \textit{cho phép điều chỉnh} đánh đổi thông qua hai ngưỡng
phân vùng, trong khi bộ quy tắc không có cơ chế tương đương.

\subsection{Đối chiếu kết quả với các công trình nghiên cứu liên quan}
\label{subsec:doi-chieu-nghien-cuu}

\paragraph{Về thứ hạng giữa các thuật toán.} Kết quả của đề án phù hợp với kết
luận của Lessmann và cộng sự~\cite{lessmann2015}: nhóm học tập hợp đứng đầu nhưng
khoảng cách với hồi quy logistic không lớn. Chênh lệch $\AUC$ đo được trong đề án
là \GapBestLR{}, thậm chí còn hẹp hơn khoảng 0,01--0,03 mà công trình trên ghi
nhận trên các bộ dữ liệu tín dụng thực. Chênh lệch hẹp hơn có thể giải thích bằng
đặc điểm của dữ liệu mô phỏng: quan hệ giữa biến tiềm ẩn và nhãn được sinh theo
một hàm log-odds có cấu trúc tương đối trơn, nên phần quan hệ phi tuyến mà các mô
hình học tập hợp khai thác được so với mô hình tuyến tính là hạn chế. Trên dữ
liệu thực, nơi quan hệ giữa các biến phức tạp hơn, khoảng cách này nhiều khả năng
sẽ rộng hơn. Các tổng quan hệ thống gần đây~\cite{shi2022} cũng đi tới kết luận
tương tự về thứ hạng giữa các nhóm thuật toán.

\paragraph{Về vai trò của đặc trưng so với thuật toán.} Kết quả phân tích SHAP
của đề án củng cố nhận định phổ biến trong các tổng quan: yếu tố quyết định chất
lượng mô hình nằm ở chất lượng và độ phong phú của đặc trưng nhiều hơn ở việc lựa
chọn thuật toán. Việc nhóm đặc trưng hành vi giao dịch đóng góp \ShapBehaviourPct{}\%
-- vượt cả nhóm CIC truyền thống -- là minh chứng cụ thể cho nhận định này, đồng
thời phù hợp với quan sát của Onay và Öztürk~\cite{onay2018} về vai trò ngày càng
lớn của dữ liệu hành vi trong bối cảnh dữ liệu lớn.

\paragraph{Về khoảng cách so với các nghiên cứu trong nước.} Các nghiên cứu ứng
dụng học máy trong đánh giá rủi ro tín dụng tại Việt Nam~\cite{thu2025,thao2025}
chủ yếu tập trung vào việc so sánh chất lượng dự báo giữa các thuật toán trên dữ
liệu dạng bảng có sẵn. Đề án này bổ sung ba nội dung mà các công trình đó chưa đề
cập đầy đủ: (i) tích hợp mô hình vào một quy trình hỗ trợ kiểm duyệt hoàn chỉnh có
luật nghiệp vụ và cơ chế giải thích; (ii) đánh giá chất lượng hiệu chuẩn xác suất
-- điều kiện cần để dùng $\PD$ trong phân vùng quyết định; và (iii) đối chiếu
định lượng với quy trình thẩm định truyền thống.

\paragraph{Về điểm khác biệt trong thiết kế thực nghiệm.} Phần lớn nghiên cứu về
chấm điểm tín dụng đánh giá mô hình bằng cách chia dữ liệu ngẫu nhiên. Đề án sử
dụng kiểm định ngoài thời gian, cho kết quả thấp hơn nhưng sát thực tế hơn. Ngoài
ra, việc bổ sung hai chỉ tiêu thường bị bỏ qua -- độ hiệu chuẩn xác suất và thời
gian chấm điểm -- cho phép đánh giá khả năng triển khai chứ không chỉ chất lượng
thống kê.

\subsection{Nhận diện hạn chế và định hướng phát triển}
\label{subsec:han-che-dinh-huong}

\subsubsection{Hạn chế của kết quả thực nghiệm}

\paragraph{Hạn chế về dữ liệu.} Đây là hạn chế cơ bản nhất. Dữ liệu mô phỏng được
sinh theo một mô hình xác suất do chính đề án thiết kế, nên cấu trúc phụ thuộc
giữa các biến phản ánh giả định của người thiết kế chứ không phải quy luật thực
tế. Ba hệ quả cụ thể: (i) giá trị tuyệt đối của các chỉ tiêu chất lượng không có
ý nghĩa dự báo về hiệu quả trên dữ liệu thực; (ii) quan hệ giữa các đặc trưng
thực tế có thể phức tạp hơn nhiều so với mô hình sinh; (iii) dữ liệu thực còn
chứa các dạng nhiễu mà mô hình sinh không tái tạo được, như sai sót nhập liệu có
hệ thống hay hành vi khai báo thiếu trung thực.

\paragraph{Hạn chế về phạm vi đo lường.} Các số đo thời gian được thực hiện trên
một máy đơn, một tiến trình, không bao gồm thời gian truyền qua mạng, truy vấn cơ
sở dữ liệu nguồn hay ảnh hưởng của tải đồng thời. Chúng có giá trị so sánh tương
đối giữa các mô hình nhưng không phải dự báo về thời gian phản hồi trong môi
trường sản xuất.

\paragraph{Hạn chế về phạm vi nghiệp vụ.} Đề án chỉ mô hình hóa thành phần $\PD$
trong phân rã tổn thất kỳ vọng; $\LGD$ và $\EAD$ không thuộc phạm vi. Ngoài ra,
đề án chưa xử lý bài toán \eng{reject inference}: mô hình được huấn luyện trên
các hồ sơ đã được cấp tín dụng, trong khi trong thực tế còn một nhóm hồ sơ bị từ
chối không có kết quả trả nợ. Bỏ qua nhóm này có thể dẫn tới sai lệch chọn mẫu.

\paragraph{Hạn chế về thời gian quan sát.} Đánh giá độ ổn định chỉ dựa trên bốn
tháng dữ liệu kiểm thử -- quá ngắn so với chu kỳ suy giảm điển hình của mô hình
tín dụng. Kết luận về độ ổn định vì vậy mang tính tham khảo.

\paragraph{Hạn chế về kiểm định công bằng.} Đề án chưa thực hiện kiểm định công
bằng thuật toán (\eng{fairness testing}) -- đánh giá xem mô hình có tạo ra chênh
lệch có hệ thống giữa các nhóm khách hàng hay không. Đây là yêu cầu ngày càng
được chú trọng đối với mô hình sử dụng trong cấp tín dụng.

\subsubsection{Định hướng phát triển}

Từ các hạn chế trên, đề án xác định năm hướng phát triển tiếp theo:

\begin{enumerate}[label=(\arabic*)]
  \item \textbf{Kiểm chứng trên dữ liệu thực.} Áp dụng toàn bộ quy trình đã xây
  dựng trên dữ liệu ẩn danh của một tổ chức tín dụng để đối chiếu kết quả. Đây là
  hướng có giá trị cao nhất vì nó trực tiếp giải quyết hạn chế cơ bản nhất.

  \item \textbf{Bổ sung xử lý reject inference.} Áp dụng các phương pháp suy luận
  nhãn cho nhóm hồ sơ bị từ chối để khắc phục sai lệch chọn mẫu.

  \item \textbf{Mở rộng sang mô hình $\LGD$ và $\EAD$.} Hoàn thiện khung tính tổn
  thất kỳ vọng, cho phép định giá theo rủi ro một cách đầy đủ.

  \item \textbf{Bổ sung kiểm định công bằng và giám sát liên tục.} Xây dựng bộ chỉ
  tiêu đánh giá công bằng và cơ chế cảnh báo tự động khi phát hiện suy giảm chất
  lượng hoặc chênh lệch giữa các nhóm khách hàng.

  \item \textbf{Nghiên cứu cơ chế cập nhật mô hình.} Thiết kế quy trình huấn luyện
  lại định kỳ có kiểm soát, bao gồm tiêu chí kích hoạt, quy trình kiểm định mô
  hình mới trước khi thay thế và cơ chế quay lui khi phát hiện vấn đề.
\end{enumerate}

\section*{Kết luận chương 3}
\addcontentsline{toc}{section}{Kết luận chương 3}

Chương 3 đã trình bày toàn bộ kết quả thực nghiệm của đề án, với năm nhóm kết quả
chính.

\textit{Về dữ liệu và đặc trưng}: tập dữ liệu gồm \NTotal{} hồ sơ vay với
\NFeaturesInitial{} đặc trưng ban đầu, rút gọn còn \NFeaturesSelected{} biến sau
ba bộ lọc. Cơ chế kiểm tra chất lượng dữ liệu được kiểm chứng bằng phép tiêm lỗi
với tỷ lệ phát hiện \ValidityDetected{}; phân tích vintage xác nhận cửa sổ quan
sát 12 tháng ghi nhận được \VintageCoverage{} tổng số trường hợp vỡ nợ.

\textit{Về chất lượng mô hình}: \NModels{} mô hình được huấn luyện và so sánh
trên cùng điều kiện. Mô hình \BestModel{} đạt $\AUC = \BestAUC{}$,
$\text{Gini} = \BestGini{}$ và $\KS = \BestKS{}$, còn cách trần Bayes \BayesAUC{}
\GapBayes{} điểm. Sai số hiệu chuẩn kỳ vọng \BestECE{} cho phép sử dụng trực
tiếp giá trị $\PD$ trong phân vùng quyết định.

\textit{Về so sánh thuật toán}: chênh lệch giữa mô hình tốt nhất và hồi quy
logistic chỉ \GapBestLR{} điểm $\AUC$, phù hợp với các công trình đối sánh quốc
tế và đặt ra câu hỏi về việc đánh đổi giữa mức cải thiện nhỏ với chi phí mất khả
năng giải thích.

\textit{Về giải thích mô hình}: nhóm đặc trưng hành vi giao dịch đóng góp
\ShapBehaviourPct{}\% vào quyết định của mô hình, vượt nhóm CIC truyền thống
(\ShapCICPct{}\%).

\textit{Về mô phỏng quy trình kiểm duyệt}: mô phỏng trên tập kiểm thử cho tỷ lệ
xử lý tự động \AutomationRate{} với tỷ lệ nợ xấu thực tế tăng đơn điệu qua ba nhóm khuyến nghị
(\BadApprove{} -- \BadReview{} -- \BadReject{}). So với quy trình thẩm định
truyền thống, mô hình giảm tỷ lệ nợ xấu trong nhóm được duyệt từ \TradBadRate{}
xuống \PropBadRate{} (giảm tương đối \BadRateReductionPct{}), đổi lại giảm
\ApproveRateDrop{} điểm phần trăm tỷ lệ duyệt và tăng \ReviewRateRise{} điểm phần
trăm khối lượng thẩm định thủ công. Thời gian xử lý một hồ sơ ở phân vị 95\% là
\SysPNinetyFive{} ms, đạt mục tiêu đề ra.

% ===========================================================================
% KẾT LUẬN
% ===========================================================================
\chapter*{KẾT LUẬN}
\addcontentsline{toc}{chapter}{KẾT LUẬN}
\markboth{KẾT LUẬN}{KẾT LUẬN}

\section*{1. Kết quả nghiên cứu}
\addcontentsline{toc}{section}{1. Kết quả nghiên cứu}

Đề án đã hoàn thành mục tiêu nghiên cứu, thiết kế và xây dựng thử nghiệm mô hình
học máy hỗ trợ kiểm duyệt tín dụng. Các kết quả chính đạt được
gồm năm nhóm.

\subsection*{1.1. Về cơ sở lý luận}

Đề án đã hệ thống hóa cơ sở lý thuyết về tín dụng, kiểm duyệt tín dụng và đánh
giá rủi ro tín dụng, trong đó làm rõ hai nội dung có ý nghĩa định hướng cho toàn
bộ thiết kế. Thứ nhất là sự phân biệt giữa \textit{kiểm duyệt} -- quá trình đánh
giá, sàng lọc hồ sơ -- với \textit{phê duyệt} -- hành vi ra quyết định của cấp có
thẩm quyền; sự phân biệt này xác định vai trò hỗ trợ chứ không thay thế của hệ
thống. Thứ hai là việc xác định vị trí của bài toán trong khung lý thuyết rủi ro
tín dụng: trong phân rã $\EL = \PD \times \LGD \times \EAD$, đề án tập trung ước
lượng thành phần $\PD$.

Đề án cũng đã phân tích bốn hạn chế mang tính cấu trúc của phương pháp kiểm duyệt
truyền thống, ba đặc thù của bài toán tín dụng so với bài toán phân loại thông
thường, và trình bày nguyên lý cùng ưu nhược điểm của năm nhóm thuật toán học máy
được nghiên cứu.

\subsection*{1.2. Về quy trình xử lý dữ liệu}

Đề án đã xây dựng quy trình hoàn chỉnh từ thu thập dữ liệu tới tập đặc trưng sẵn
sàng cho mô hình, với ba đóng góp cụ thể:

\begin{itemize}
  \item \textbf{Bộ ràng buộc kiểm tra chất lượng dữ liệu} gồm \NConstraints{}
  ràng buộc nghiệp vụ, kèm cơ chế kiểm chứng bằng phép tiêm lỗi nhằm xác nhận
  bản thân cơ chế kiểm tra hoạt động đúng.

  \item \textbf{Thuật toán tính đặc trưng bảo đảm tính đúng đắn theo thời điểm}
  bằng kỹ thuật luỹ kế kết hợp ghép gần nhất theo thời gian, đưa chi phí tính
  toán từ $O(N \times |W| \times \bar{m})$ xuống tuyến tính. Quan trọng hơn chi
  phí, phương pháp này bảo đảm tính đúng đắn bằng chính cấu trúc thuật toán chứ
  không phụ thuộc vào kiểm tra bổ sung.

  \item \textbf{Quy trình lựa chọn đặc trưng ba bước} theo năng lực dự báo, mức
  tương quan và độ ổn định, rút gọn tập đặc trưng từ \NFeaturesInitial{} xuống
  \NFeaturesSelected{} biến.
\end{itemize}

\subsection*{1.3. Về mô hình học máy}

Đề án đã huấn luyện, hiệu chuẩn và so sánh \NModels{} mô hình trên cùng một điều
kiện thực nghiệm. Mô hình được chọn triển khai là \BestModel{}, đạt
$\AUC = \BestAUC{}$, $\text{Gini} = \BestGini{}$, $\KS = \BestKS{}$ và sai số
hiệu chuẩn kỳ vọng \BestECE{}.

Ngoài các con số cụ thể, thực nghiệm cho ba kết quả có giá trị nhận thức. Thứ
nhất, \textbf{khoảng cách giữa các thuật toán là nhỏ}: chênh lệch $\AUC$ giữa mô
hình được chọn và hồi quy logistic chỉ \GapBestLR{}, trong khi hồi quy logistic
có ưu thế rõ rệt về khả năng giải thích, chi phí huấn luyện và chất lượng hiệu
chuẩn tự nhiên. Thứ hai, \textbf{bảng xếp hạng $\AUC$ có giới hạn độ phân giải}:
phép đo dao động khoảng $3 \times 10^{-4}$ giữa các lần chạy, tức cùng bậc với
khoảng cách giữa các mô hình dẫn đầu, nên thứ hạng ở chữ số thập phân thứ tư
không phải bằng chứng về sự vượt trội. Đề án vì vậy chọn mô hình triển khai theo
tiêu chí biên Pareto -- nhanh nhất trong nhóm tương đương về chất lượng -- thay
vì lấy thẳng mô hình đứng đầu bảng. Thứ ba, \textbf{khoảng cách tới trần Bayes
cho biết giới hạn của việc cải tiến thuật toán}: mô hình được chọn còn cách trần
Bayes \BayesAUC{} khoảng
\GapBayes{} điểm $\AUC$, và phần chênh lệch này bắt nguồn từ việc dữ liệu quan sát được
chỉ là phép đo có nhiễu của mức tín nhiệm thật, chứ không phải do thuật toán yếu.
Hệ quả thực tiễn là cải thiện đáng kể chỉ đến từ việc bổ sung nguồn dữ liệu mới,
không phải từ thuật toán phức tạp hơn.

\subsection*{1.4. Về giá trị của các nguồn dữ liệu}

Phân tích SHAP cho thấy nhóm đặc trưng hành vi giao dịch đóng góp
\ShapBehaviourPct{}\% vào quyết định của mô hình, cao hơn nhóm lịch sử tín dụng
CIC (\ShapCICPct{}\%) vốn được coi là nguồn thông tin truyền thống quan trọng
nhất. Đây là căn cứ định lượng cho khuyến nghị khai thác dữ liệu giao dịch trong
công tác kiểm duyệt tín dụng: nguồn dữ liệu này do tổ chức tín dụng tự sở hữu,
không phát sinh chi phí truy vấn bên ngoài, phản ánh tình trạng tài chính hiện
tại và luôn có sẵn kể cả với nhóm khách hàng thiếu lịch sử tín dụng.

\subsection*{1.5. Về quy trình hỗ trợ kiểm duyệt}

Đề án đã hiện thực quy trình hoàn chỉnh kết hợp \NRules{} luật nghiệp vụ với mô
hình học máy theo logic ưu tiên tám mức, đóng gói thành dịch vụ REST với năm điểm
cuối hoạt động đúng. Mô phỏng trên tập kiểm thử cho tỷ lệ nợ xấu thực tế tăng đơn
điệu qua ba nhóm khuyến nghị (\BadApprove{} -- \BadReview{} -- \BadReject{}), tỷ
lệ xử lý tự động \AutomationRate{} và thời gian xử lý một hồ sơ ở phân vị 95\% là
\SysPNinetyFive{} ms.

So với quy trình thẩm định truyền thống chỉ dựa trên luật nghiệp vụ, mô hình
giảm tỷ lệ nợ xấu trong nhóm hồ sơ được duyệt từ \TradBadRate{} xuống
\PropBadRate{}, tức giảm tương đối \BadRateReductionPct{}. Đổi lại, tỷ lệ duyệt
giảm \ApproveRateDrop{} điểm phần trăm và khối lượng thẩm định thủ công tăng
\ReviewRateRise{} điểm phần trăm. Đánh đổi này điều chỉnh được thông qua hai
ngưỡng phân vùng, trong khi bộ quy tắc cố định không có cơ chế tương đương.

\section*{2. Hạn chế của nghiên cứu}
\addcontentsline{toc}{section}{2. Hạn chế của nghiên cứu}

Các hạn chế của kết quả thực nghiệm đã được phân tích chi tiết tại
mục~\ref{subsec:han-che-dinh-huong}. Ở tầm tổng thể, đề án có ba hạn chế cơ bản.

\textit{Thứ nhất là hạn chế về dữ liệu.} Do đặc thù bảo mật của dữ liệu tín dụng
cá nhân, đề án sử dụng dữ liệu mô phỏng sinh theo một mô hình xác suất do chính
đề án thiết kế, nên cấu trúc phụ thuộc giữa các biến phản ánh giả định của người
thiết kế chứ không phải quy luật thực tế. Giá trị tuyệt đối của các chỉ tiêu chất
lượng vì vậy không có ý nghĩa dự báo về hiệu quả trên dữ liệu thực; giá trị của
thực nghiệm nằm ở việc so sánh \textit{tương đối} giữa các mô hình trong cùng
điều kiện và ở việc kiểm chứng tính đúng đắn của toàn bộ quy trình.

\textit{Thứ hai là hạn chế về phương pháp.} Đề án chưa xử lý bài toán suy luận
nhãn cho hồ sơ bị từ chối (\eng{reject inference}) nên có thể chịu sai lệch chọn
mẫu; chưa thực hiện kiểm định công bằng thuật toán; và việc đánh giá độ ổn định
chỉ dựa trên bốn tháng dữ liệu kiểm thử, quá ngắn so với chu kỳ suy giảm điển
hình của mô hình tín dụng.

\textit{Thứ ba là hạn chế về phạm vi.} Đề án chỉ mô hình hóa thành phần $\PD$
trong phân rã tổn thất kỳ vọng; $\LGD$ và $\EAD$ không thuộc phạm vi nghiên cứu.
Mô hình được xây dựng ở mức nguyên mẫu trong môi trường nghiên cứu, chưa triển
khai trên hệ thống nghiệp vụ thực tế của một tổ chức tín dụng, và các số đo thời
gian chỉ có giá trị so sánh tương đối giữa các mô hình.

\section*{3. Hướng phát triển}
\addcontentsline{toc}{section}{3. Hướng phát triển}

Trên cơ sở các hạn chế đã nhận diện, đề án đề xuất năm hướng phát triển tiếp
theo, sắp xếp theo mức độ ưu tiên:

\begin{enumerate}[label=(\arabic*)]
  \item \textbf{Kiểm chứng trên dữ liệu thực} của một tổ chức tín dụng. Đây là
  hướng có giá trị cao nhất vì trực tiếp giải quyết hạn chế cơ bản nhất của đề
  án; mã nguồn đã được thiết kế để chuyển sang dữ liệu thực chỉ bằng cách thay
  thành phần tiếp nhận dữ liệu.

  \item \textbf{Bổ sung xử lý suy luận nhãn cho hồ sơ bị từ chối}, mở rộng phạm
  vi áp dụng của mô hình tới toàn bộ hồ sơ nộp vào.

  \item \textbf{Mở rộng sang mô hình $\LGD$ và $\EAD$} để hoàn thiện khung tính
  tổn thất kỳ vọng, cho phép định giá theo rủi ro đầy đủ.

  \item \textbf{Bổ sung kiểm định công bằng thuật toán và cơ chế cảnh báo tự
  động} khi phát hiện suy giảm chất lượng hoặc chênh lệch giữa các nhóm khách
  hàng.

  \item \textbf{Nghiên cứu cơ chế cập nhật mô hình có kiểm soát}, gồm tiêu chí
  kích hoạt huấn luyện lại, quy trình kiểm định mô hình mới và cơ chế quay lui
  khi phát hiện vấn đề trong vận hành.
\end{enumerate}

\vspace{0.8cm}

Tóm lại, đề án đã xây dựng được một mô hình học máy hỗ trợ kiểm duyệt tín dụng
hoàn chỉnh từ khâu tiếp nhận dữ liệu tới khâu đưa ra khuyến nghị kèm lý giải,
đồng thời đánh giá định lượng giá trị mà mô hình mang lại so với quy trình thẩm
định
truyền thống. Kết quả nghiên cứu cho thấy việc kết hợp mô hình học máy với bộ
luật nghiệp vụ, trong đó mô hình đóng vai trò cung cấp thông tin định lượng còn
quyết định cuối cùng vẫn thuộc về người có thẩm quyền, là hướng tiếp cận vừa khả
thi về kỹ thuật vừa phù hợp với yêu cầu quản trị rủi ro và tuân thủ của hoạt động
cấp tín dụng.


% ===========================================================================
% TÀI LIỆU THAM KHẢO
% ===========================================================================
\cleardoublepage
% ===========================================================================
% DANH MỤC TÀI LIỆU THAM KHẢO
% ===========================================================================
\chapter*{DANH MỤC TÀI LIỆU THAM KHẢO}
\addcontentsline{toc}{chapter}{DANH MỤC TÀI LIỆU THAM KHẢO}
\markboth{DANH MỤC TÀI LIỆU THAM KHẢO}{DANH MỤC TÀI LIỆU THAM KHẢO}

% Danh mục tài liệu chứa nhiều tên riêng và tiêu đề tiếng Anh mà bộ quy tắc ngắt
% âm tiết tiếng Việt không xử lý được. Nới lỏng tiêu chí xuống dòng trong phạm vi
% chương này để tránh tràn lề; chấp nhận khoảng trắng giữa từ thưa hơn đôi chút.
% Danh mục tài liệu được căn lề trái (không căn đều hai bên).
%
% Lý do: định dạng XeLaTeX đang dùng không nạp sẵn bộ quy tắc ngắt âm tiết
% (kiểm chứng bằng \showhyphens: không có điểm ngắt nào được đề xuất). Với phần
% chữ tiếng Việt điều này không gây vấn đề vì tiếng Việt vốn không ngắt âm tiết;
% nhưng các tiêu đề và tên riêng tiếng Anh dài trong danh mục tài liệu sẽ không
% ngắt được, buộc TeX phải giãn khoảng trắng rất thưa để căn đều hai bên.
%
% Căn lề trái là cách trình bày chuẩn và phổ biến cho danh mục tài liệu tham
% khảo, đồng thời loại bỏ hoàn toàn hiện tượng dòng thưa nói trên.
\begingroup
\raggedright
\emergencystretch=2.5em

\section*{Tài liệu tham khảo tiếng Việt}

\begin{enumerate}[label={[\arabic*]}, ref=\arabic*, leftmargin=*, labelsep=0.6em,
                  itemsep=5pt]

\item\label{ref:tt15} Ngân hàng Nhà nước Việt Nam (2023), \textit{Thông tư số
15/2023/TT-NHNN quy định về hoạt động thông tin tín dụng của Ngân hàng Nhà nước
Việt Nam}, Hà Nội.

\item\label{ref:tt11} Ngân hàng Nhà nước Việt Nam (2021), \textit{Thông tư số
11/2021/TT-NHNN quy định về phân loại tài sản có, mức trích, phương pháp trích
lập dự phòng rủi ro và việc sử dụng dự phòng để xử lý rủi ro trong hoạt động của
tổ chức tín dụng, chi nhánh ngân hàng nước ngoài}, Hà Nội.

\item\label{ref:nhnn-ml} Ngân hàng Nhà nước Việt Nam (2023), \textit{Machine
Learning trong quản trị rủi ro tín dụng của ngân hàng}, Tạp chí Ngân hàng.

\item\label{ref:thu2025} Huỳnh Thị Ái Thu (2025), \textit{Ứng dụng học máy trong
xây dựng mô hình cảnh báo sớm rủi ro tín dụng khách hàng cá nhân tại ngân hàng
thương mại Việt Nam}, Luận văn Thạc sĩ, Trường Đại học Kinh tế Thành phố Hồ Chí
Minh.

\item\label{ref:thao2025} Nguyễn Thị Phương Thảo (2025), \textit{Xây dựng mô hình
chấm điểm xếp hạng tín dụng khách hàng doanh nghiệp tại các ngân hàng thương mại
ứng dụng các mô hình học máy hiện đại}, Luận văn Thạc sĩ, Trường Đại học Kinh tế
Thành phố Hồ Chí Minh.

\item\label{ref:sang2021} Nguyễn Minh Sáng, Nguyễn Thị Hồng Vinh (2021),
``Lược khảo các nghiên cứu về ứng dụng công nghệ phân tích dữ liệu lớn và trí tuệ
nhân tạo trong quản trị rủi ro tín dụng của ngân hàng'', \textit{Tạp chí Kinh tế
và Ngân hàng châu Á}, số 186.

\item\label{ref:hoa2025} Nguyễn Thị Hoa, Nguyễn Thị Thu (2025), ``Ứng dụng phân
tích dữ liệu lớn và trí tuệ nhân tạo trong tối ưu hóa quản trị rủi ro tín dụng
tại các ngân hàng số ở Việt Nam'', \textit{Tạp chí Khoa học Trường Đại học Mở Hà
Nội}, số 2A.

\end{enumerate}

\section*{Tài liệu tham khảo tiếng Anh}

\begin{enumerate}[label={[\arabic*]}, ref=\arabic*, leftmargin=*, labelsep=0.6em,
                  itemsep=5pt, start=8]

\item\label{ref:lessmann2015} Lessmann, S., Baesens, B., Seow, H.-V., \& Thomas,
L. C. (2015), ``Benchmarking state-of-the-art classification algorithms for credit
scoring: An update of research'', \textit{European Journal of Operational
Research}, 247(1), pp. 124--136.

\item\label{ref:shi2022} Shi, S., Tse, R., Luo, W., D'Addona, S., \& Pau, G.
(2022), ``Machine learning-driven credit risk: A systemic review'', \textit{Neural
Computing and Applications}, 34, pp. 14327--14339.

\item\label{ref:onay2018} Onay, C., \& Öztürk, E. (2018), ``A review of credit
scoring research in the age of Big Data'', \textit{Journal of Financial Regulation
and Compliance}, 26(3), pp. 382--405.

\item\label{ref:ayari2026} Ayari, H., Guetari, R., \& Kraïem, N. (2026),
``Machine learning powered financial credit scoring: A systematic literature
review'', \textit{Artificial Intelligence Review}, 59, 13.

\item\label{ref:data2023} ``Machine Learning for Credit Risk Prediction: A
Systematic Literature Review'' (2023), \textit{Data}, 8(11), 169.

\item\label{ref:cortes1995} Cortes, C., \& Vapnik, V. (1995), ``Support-Vector
Networks'', \textit{Machine Learning}, 20(3), pp. 273--297.

\item\label{ref:vapnik1998} Vapnik, V. N. (1998), \textit{Statistical Learning
Theory}, John Wiley \& Sons.

\item\label{ref:williams2001} Williams, C. K. I., \& Seeger, M. (2001), ``Using
the Nyström Method to Speed Up Kernel Machines'', \textit{Advances in Neural
Information Processing Systems}, 13, pp. 682--688.

\item\label{ref:platt1999} Platt, J. C. (1999), ``Probabilistic Outputs for
Support Vector Machines and Comparisons to Regularized Likelihood Methods'',
\textit{Advances in Large Margin Classifiers}, MIT Press, pp. 61--74.

\item\label{ref:breiman1984} Breiman, L., Friedman, J. H., Olshen, R. A., \&
Stone, C. J. (1984), \textit{Classification and Regression Trees}, Wadsworth.

\item\label{ref:hastie2009} Hastie, T., Tibshirani, R., \& Friedman, J. (2009),
\textit{The Elements of Statistical Learning: Data Mining, Inference, and
Prediction}, 2nd edition, Springer.

\item\label{ref:mitchell1997} Mitchell, T. M. (1997), \textit{Machine Learning},
McGraw-Hill.

\item\label{ref:chen2016} Chen, T., \& Guestrin, C. (2016), ``XGBoost: A Scalable
Tree Boosting System'', \textit{Proceedings of the 22nd ACM SIGKDD International
Conference on Knowledge Discovery and Data Mining}, pp. 785--794.

\item\label{ref:ke2017} Ke, G., Meng, Q., Finley, T., Wang, T., Chen, W., Ma, W.,
Ye, Q., \& Liu, T.-Y. (2017), ``LightGBM: A Highly Efficient Gradient Boosting
Decision Tree'', \textit{Advances in Neural Information Processing Systems}, 30,
pp. 3146--3154.

\item\label{ref:prokhorenkova2018} Prokhorenkova, L., Gusev, G., Vorobev, A.,
Dorogush, A. V., \& Gulin, A. (2018), ``CatBoost: unbiased boosting with
categorical features'', \textit{Advances in Neural Information Processing
Systems}, 31, pp. 6638--6648.

\item\label{ref:breiman2001} Breiman, L. (2001), ``Random Forests'',
\textit{Machine Learning}, 45(1), pp. 5--32.

\item\label{ref:friedman2001} Friedman, J. H. (2001), ``Greedy function
approximation: A gradient boosting machine'', \textit{The Annals of Statistics},
29(5), pp. 1189--1232.

\item\label{ref:lundberg2017} Lundberg, S. M., \& Lee, S.-I. (2017), ``A Unified
Approach to Interpreting Model Predictions'', \textit{Advances in Neural
Information Processing Systems}, 30, pp. 4765--4774.

\item\label{ref:lundberg2020} Lundberg, S. M., Erion, G., Chen, H., DeGrave, A.,
Prutkin, J. M., Nair, B., Katz, R., Himmelfarb, J., Bansal, N., \& Lee, S.-I.
(2020), ``From local explanations to global understanding with explainable AI for
trees'', \textit{Nature Machine Intelligence}, 2(1), pp. 56--67.

\item\label{ref:zadrozny2002} Zadrozny, B., \& Elkan, C. (2002), ``Transforming
classifier scores into accurate multiclass probability estimates'',
\textit{Proceedings of the 8th ACM SIGKDD International Conference on Knowledge
Discovery and Data Mining}, pp. 694--699.

\item\label{ref:niculescu2005} Niculescu-Mizil, A., \& Caruana, R. (2005),
``Predicting good probabilities with supervised learning'', \textit{Proceedings of
the 22nd International Conference on Machine Learning}, pp. 625--632.

\item\label{ref:bergstra2012} Bergstra, J., \& Bengio, Y. (2012), ``Random Search
for Hyper-Parameter Optimization'', \textit{Journal of Machine Learning
Research}, 13, pp. 281--305.

\item\label{ref:siddiqi2017} Siddiqi, N. (2017), \textit{Intelligent Credit
Scoring: Building and Implementing Better Credit Risk Scorecards}, 2nd edition,
John Wiley \& Sons.

\item\label{ref:thomas2017} Thomas, L. C., Crook, J., \& Edelman, D. (2017),
\textit{Credit Scoring and Its Applications}, 2nd edition, Society for Industrial
and Applied Mathematics.

\item\label{ref:baesens2016} Baesens, B., Roesch, D., \& Scheule, H. (2016),
\textit{Credit Risk Analytics: Measurement Techniques, Applications, and Examples
in SAS}, John Wiley \& Sons.

\item\label{ref:bcbs2006} Basel Committee on Banking Supervision (2006),
\textit{International Convergence of Capital Measurement and Capital Standards: A
Revised Framework}, Bank for International Settlements.

\item\label{ref:sculley2015} Sculley, D., Holt, G., Golovin, D., Davydov, E.,
Phillips, T., Ebner, D., Chaudhary, V., Young, M., Crespo, J.-F., \& Dennison, D.
(2015), ``Hidden Technical Debt in Machine Learning Systems'', \textit{Advances in
Neural Information Processing Systems}, 28, pp. 2503--2511.

\item\label{ref:gama2014} Gama, J., Žliobaitė, I., Bifet, A., Pechenizkiy, M., \&
Bouchachia, A. (2014), ``A survey on concept drift adaptation'', \textit{ACM
Computing Surveys}, 46(4), pp. 1--37.

\item\label{ref:chawla2002} Chawla, N. V., Bowyer, K. W., Hall, L. O., \&
Kegelmeyer, W. P. (2002), ``SMOTE: Synthetic Minority Over-sampling Technique'',
\textit{Journal of Artificial Intelligence Research}, 16, pp. 321--357.

\item\label{ref:elkan2001} Elkan, C. (2001), ``The Foundations of Cost-Sensitive
Learning'', \textit{Proceedings of the 17th International Joint Conference on
Artificial Intelligence}, pp. 973--978.

\end{enumerate}
\endgroup


% ===========================================================================
% PHỤ LỤC
% ===========================================================================
\cleardoublepage
\appendix
\renewcommand{\appendixname}{PHỤ LỤC}
% ===========================================================================
% PHỤ LỤC
% ===========================================================================

\chapter{CẤU TRÚC MÃ NGUỒN VÀ HƯỚNG DẪN TÁI LẬP KẾT QUẢ}
% Tiêu đề chạy rút gọn: tiêu đề chương đầy đủ dài hơn một dòng đầu trang,
% khiến fancyhdr phải nới \headheight và làm lệch khung chữ của cả chương.
\markboth{Phụ lục A. Cấu trúc mã nguồn và hướng dẫn tái lập kết quả}{}
\label{app:source}

\section{Tổ chức mã nguồn}

Toàn bộ quy trình được hiện thực bằng Python và tổ chức thành một gói phần mềm có
cấu trúc phản ánh đúng mạch nội dung nghiên cứu của đề án. Mỗi gói con tương ứng
một mục của Chương~\ref{chap:xay-dung-he-thong}, nhờ đó người đọc có thể đối
chiếu trực tiếp giữa phần trình bày phương pháp và phần hiện thực.

\begin{verbatim}
config/config.yaml           Cấu hình tập trung toàn bộ quy trình
src/creditreview/
  config.py                  Nạp và truy cập cấu hình
  data/                      2.2 Thu thập và tổng hợp dữ liệu
    generator.py           Sinh dữ liệu mô phỏng (biến tiềm ẩn)
    reference.py           Bảng tham chiếu nghiệp vụ
    profiling.py           Thống kê mô tả, kiểm tra hợp lệ, tiêm lỗi
  preprocessing/             2.3 Tiền xử lý và làm sạch dữ liệu
    cleaning.py              Xử lý giá trị khuyết thiếu và ngoại lệ
    encoding.py              Chuẩn hóa và mã hóa biến
    splitting.py             Phân chia tập dữ liệu
  features/                  2.4 Trích xuất và lựa chọn đặc trưng
    definitions.py           Đăng ký đặc trưng theo nhóm nguồn dữ liệu
    extraction.py            Trích xuất đặc trưng đúng theo thời điểm
    selection.py             Lựa chọn tập đặc trưng qua ba bộ lọc
  labeling/                  2.5 Xây dựng nhãn và quy tắc phê duyệt
    target.py                Gán nhãn rủi ro, phân tích vintage
    cutoff.py                Xác định điểm cắt phê duyệt
  models/                    2.6, 2.7 Thiết kế và huấn luyện mô hình
    zoo.py                   Nhà máy mô hình (12 mô hình)
    train.py                 Huấn luyện, hiệu chuẩn, đánh giá
    tuning.py                Tinh chỉnh siêu tham số
    metrics.py               Bộ chỉ tiêu đánh giá
    woe.py                   Biến đổi WOE và thẻ điểm
    explain.py               Giải thích mô hình (SHAP, phép hoán vị)
    plots.py                 Sinh đồ thị
  decision/                  Bộ máy hỗ trợ ra quyết định
    rules.py                 Bộ luật nghiệp vụ
    policy.py                Chính sách hạn mức và định giá
    engine.py                Tổng hợp khuyến nghị theo thứ tự ưu tiên
  serving/                   Dịch vụ REST
    api.py                   Các điểm cuối nghiệp vụ và vận hành
    repository.py            Kho đặc trưng phục vụ tra cứu
  monitoring/
    drift.py                 Giám sát PSI và chất lượng theo thời gian
  dashboard/
    app.py                   Giao diện trực quan hóa kết quả
scripts/                     Bảy bước tái lập kết quả
results/                     Mô hình, bảng và đồ thị kết quả
thesis/                      Mã nguồn LaTeX của đề án
\end{verbatim}

\section{Hướng dẫn tái lập kết quả}

\subsection{Cài đặt môi trường}

\begin{verbatim}
pip install -r requirements.txt
\end{verbatim}

Yêu cầu Python phiên bản 3.10 trở lên.

\subsection{Chạy toàn bộ quy trình}

Cách nhanh nhất là chạy toàn bộ bằng một lệnh:

\begin{verbatim}
python scripts/run_all.py
\end{verbatim}

Hoặc chạy tuần tự từng bước để theo dõi kết quả trung gian:

\begin{verbatim}
# Bước 1 - Chuẩn bị dữ liệu: sinh, kiểm tra chất lượng,
#          gán nhãn, phân chia tập
python scripts/01_prepare_data.py

# Bước 2 - Phân tích Information Value và lựa chọn đặc trưng
python scripts/02_build_features.py

# Bước 3 - Tinh chỉnh siêu tham số (tùy chọn)
python scripts/03_tune_hyperparams.py --trials 14

# Bước 4 - Huấn luyện và so sánh 12 mô hình
python scripts/04_train_models.py

# Bước 5 - Mô phỏng quy trình kiểm duyệt và đánh giá chung
python scripts/05_simulate_evaluate.py

# Bước 6 - Xuất bảng LaTeX cho đề án
python scripts/06_export_latex_tables.py

# Bước 7 - Xuất từ điển dữ liệu các bảng đầu vào (Excel, tùy chọn)
python scripts/07_export_data_dictionary.py
\end{verbatim}

Hạt giống ngẫu nhiên được cố định trong \code{config/config.yaml}
(\code{project.seed}), bảo đảm kết quả lặp lại được giữa các lần chạy.

\subsection{Chạy nhanh để kiểm tra}

Để kiểm tra quy trình hoạt động mà không mất nhiều thời gian:

\begin{verbatim}
python scripts/01_prepare_data.py --n-customers 3000 \
    --n-applications 4500
python scripts/02_build_features.py
python scripts/04_train_models.py --no-stacking --no-shap \
    --models logistic_regression lightgbm
python scripts/05_simulate_evaluate.py --n-applications 1000 --no-api
\end{verbatim}

\subsection{Khởi động dịch vụ và giao diện}

\begin{verbatim}
# Dịch vụ REST hỗ trợ kiểm duyệt
uvicorn creditreview.serving.api:app --port 8000
#   Tài liệu API tự sinh: http://localhost:8000/docs

# Giao diện trực quan hóa kết quả
streamlit run src/creditreview/dashboard/app.py
\end{verbatim}

\subsection{Biên dịch đề án}

\begin{verbatim}
cd thesis
xelatex main.tex && xelatex main.tex
\end{verbatim}

Cần dùng XeLaTeX do tài liệu sử dụng tiếng Việt và phông Times New Roman. Chạy
hai lần để cập nhật mục lục và tham chiếu chéo.

Số liệu trong đề án được sinh tự động từ kết quả thực nghiệm qua
\code{scripts/06\_export\_latex\_tables.py}: chạy lại thực nghiệm là toàn bộ bảng
biểu và các con số trích dẫn trong phần diễn giải tự cập nhật, không phải sửa
tay.

\chapter{DANH MỤC LUẬT NGHIỆP VỤ}
% Tiêu đề chạy rút gọn: tiêu đề chương đầy đủ dài hơn một dòng đầu trang,
% khiến fancyhdr phải nới \headheight và làm lệch khung chữ của cả chương.
\markboth{Phụ lục B. Danh mục luật nghiệp vụ}{}
\label{app:rules}

Bảng dưới đây liệt kê đầy đủ bộ luật nghiệp vụ được cài đặt trong mô hình. Ba
mức hành động gồm: \code{HARD\_REJECT} -- từ chối, có quyền phủ quyết cao nhất;
\code{REVIEW} -- buộc chuyển thẩm định thủ công; \code{WARNING} -- ghi nhận vào
danh sách lý do nhưng không thay đổi khuyến nghị.

\begingroup
\small
\setlength{\tabcolsep}{4.0pt}
\begin{longtable}{@{}L{1.436cm}L{3.638cm}L{2.393cm}L{2.400cm}L{4.308cm}@{}}
\caption{Danh mục đầy đủ luật nghiệp vụ được cài đặt trong mô hình}\label{tab:rule-catalogue}\\
\toprule
\textbf{Mã luật} & \textbf{Tên luật} & \textbf{Nhóm} & \textbf{Hành động} & \textbf{Diễn giải} \\
\midrule
\endfirsthead
\multicolumn{5}{@{}l}{\small\textit{(tiếp theo trang trước)}}\\
\toprule
\textbf{Mã luật} & \textbf{Tên luật} & \textbf{Nhóm} & \textbf{Hành động} & \textbf{Diễn giải} \\
\midrule
\endhead
\bottomrule
\multicolumn{5}{r@{}}{\small\textit{(xem tiếp trang sau)}}\\
\endfoot
\bottomrule
\endlastfoot
R001 & Tuổi nằm ngoài khoảng cho phép [20; 62] & Điều kiện cơ bản & HARD\_\allowbreak REJECT & Khách hàng phải trong độ tuổi từ 20 đến 62 tại thời điểm đề nghị vay \\
R002 & Thu nhập dưới mức tối thiểu 6 triệu VND & Điều kiện cơ bản & HARD\_\allowbreak REJECT & Thu nhập hàng tháng không đạt mức tối thiểu của sản phẩm \\
R003 & Thâm niên công việc dưới 6 tháng & Điều kiện cơ bản & REVIEW & Thâm niên công việc chưa đủ để đánh giá tính ổn định của thu nhập \\
R101 & Nhóm nợ xấu nhất 12 tháng vượt nhóm 2 & Lịch sử tín dụng & HARD\_\allowbreak REJECT & Khách hàng phát sinh nợ xấu trong 12 tháng gần nhất \\
R102 & Đang có nợ nhóm 3 trở lên & Lịch sử tín dụng & HARD\_\allowbreak REJECT & Khách hàng đang có khoản nợ được phân loại từ nhóm 3 (nợ dưới tiêu chuẩn) \\
R103 & Có nợ đã xử lý rủi ro / xóa nợ & Lịch sử tín dụng & HARD\_\allowbreak REJECT & Lịch sử có khoản nợ đã được xử lý rủi ro hoặc xóa nợ \\
R104 & Số ngày quá hạn tối đa 12 tháng vượt 30 ngày & Lịch sử tín dụng & REVIEW & Từng chậm trả vượt ngưỡng cho phép trong 12 tháng gần nhất \\
R105 & Quan hệ tín dụng với hơn 6 tổ chức & Lịch sử tín dụng & REVIEW & Số lượng tổ chức tín dụng đang có quan hệ ở mức cao, tiềm ẩn rủi ro vay đảo nợ \\
R106 & Truy vấn tín dụng 3 tháng vượt 5 lần & Lịch sử tín dụng & WARNING & Khách hàng tìm kiếm tín dụng dồn dập tại nhiều tổ chức \\
R107 & Tỷ lệ sử dụng hạn mức trên 95\% & Lịch sử tín dụng & WARNING & Hạn mức tín dụng hiện có đã được sử dụng gần hết \\
R201 & Tỷ lệ nợ trên thu nhập (DTI) vượt 60\% & Khả năng trả nợ & HARD\_\allowbreak REJECT & Tổng nghĩa vụ trả nợ vượt quá khả năng chi trả từ thu nhập \\
R202 & Thu nhập còn lại sau trả nợ không đủ chi phí sinh hoạt & Khả năng trả nợ & HARD\_\allowbreak REJECT & Sau khi trừ nghĩa vụ trả nợ và chi phí sinh hoạt tối thiểu, thu nhập còn lại âm \\
R203 & Tỷ lệ trả góp trên thu nhập (PTI) vượt 45\% & Khả năng trả nợ & REVIEW & Riêng khoản vay mới đã chiếm tỷ trọng lớn trong thu nhập \\
R204 & Số tiền vay vượt 60 lần thu nhập hàng tháng & Khả năng trả nợ & REVIEW & Quy mô khoản vay rất lớn so với thu nhập, cần thẩm định bổ sung \\
R301 & Tỷ lệ vay trên giá trị TSBĐ (LTV) vượt 85\% & Tài sản bảo đảm & REVIEW & Giá trị tài sản bảo đảm không đủ tỷ lệ bao phủ yêu cầu \\
R401 & Không lấy được dữ liệu CIC & Chất lượng dữ liệu & REVIEW & Hệ thống CIC không phản hồi - chuyển thẩm định thủ công theo quy trình dự phòng \\
R402 & Khách hàng thiếu lịch sử tín dụng (thin-file) & Chất lượng dữ liệu & WARNING & Không có/ít lịch sử tín dụng; đánh giá dựa chủ yếu vào dữ liệu hành vi giao dịch \\
R403 & Thiếu dữ liệu hành vi giao dịch & Chất lượng dữ liệu & REVIEW & Số giao dịch trong 180 ngày quá ít để tính đặc trưng hành vi đáng tin cậy \\
R501 & Dòng tiền vào thực tế thấp hơn nhiều so với thu nhập khai báo & Dấu hiệu bất thường & REVIEW & Khách hàng nhận lương qua tài khoản nhưng dòng tiền vào 30 ngày dưới 35\% thu nhập khai báo - cần xác minh thu nhập \\
R502 & Tỷ trọng chi tiêu nhóm rủi ro cao bất thường & Dấu hiệu bất thường & REVIEW & Trên 15\% dòng tiền ra thuộc nhóm giao dịch rủi ro cao \\
R503 & Không phát sinh giao dịch trong 60 ngày gần nhất & Dấu hiệu bất thường & WARNING & Tài khoản không hoạt động trong thời gian dài trước khi đề nghị vay \\
\end{longtable}
\endgroup


\chapter{DANH MỤC ĐẶC TRƯNG}
% Tiêu đề chạy rút gọn: tiêu đề chương đầy đủ dài hơn một dòng đầu trang,
% khiến fancyhdr phải nới \headheight và làm lệch khung chữ của cả chương.
\markboth{Phụ lục C. Danh mục đặc trưng}{}
\label{app:features}

Danh mục đầy đủ các đặc trưng được sử dụng, phân theo nhóm nguồn dữ liệu. Cột
``Nguồn dữ liệu'' cho biết đặc trưng được tính từ bảng dữ liệu nào.

\begingroup
\small
\setlength{\tabcolsep}{4.0pt}
\begin{longtable}{@{}L{3.761cm}L{2.121cm}L{2.314cm}L{1.061cm}L{4.918cm}@{}}
\caption{Danh mục đặc trưng đầu vào của mô hình}\label{tab:feature-catalogue}\\
\toprule
\textbf{Tên đặc trưng} & \textbf{Nhóm} & \textbf{Nguồn dữ liệu} & \textbf{Kiểu} & \textbf{Mô tả} \\
\midrule
\endfirsthead
\multicolumn{5}{@{}l}{\small\textit{(tiếp theo trang trước)}}\\
\toprule
\textbf{Tên đặc trưng} & \textbf{Nhóm} & \textbf{Nguồn dữ liệu} & \textbf{Kiểu} & \textbf{Mô tả} \\
\midrule
\endhead
\bottomrule
\multicolumn{5}{r@{}}{\small\textit{(xem tiếp trang sau)}}\\
\endfoot
\bottomrule
\endlastfoot
age & Thông tin cá nhân & HỒ SƠ KH & num & Tuổi khách hàng \\
gender & Thông tin cá nhân & HỒ SƠ KH & cat & Giới tính \\
marital\_\allowbreak status & Thông tin cá nhân & HỒ SƠ KH & cat & Tình trạng hôn nhân \\
num\_\allowbreak dependents & Thông tin cá nhân & HỒ SƠ KH & num & Số người phụ thuộc \\
education\_\allowbreak code & Thông tin cá nhân & HỒ SƠ KH & num & Trình độ học vấn (1-5) \\
occupation & Thông tin cá nhân & HỒ SƠ KH & cat & Nghề nghiệp \\
employer\_\allowbreak type & Thông tin cá nhân & HỒ SƠ KH & cat & Loại hình tổ chức làm việc \\
job\_\allowbreak tenure\_\allowbreak months & Thông tin cá nhân & HỒ SƠ KH & num & Thâm niên công việc (tháng) \\
monthly\_\allowbreak income & Thông tin cá nhân & HỒ SƠ KH & num & Thu nhập khai báo hàng tháng \\
province & Thông tin cá nhân & HỒ SƠ KH & cat & Tỉnh/thành cư trú \\
residence\_\allowbreak type & Thông tin cá nhân & HỒ SƠ KH & cat & Hình thức nhà ở \\
customer\_\allowbreak since\_\allowbreak months & Thông tin cá nhân & HỒ SƠ KH & num & Thời gian quan hệ với ngân hàng \\
has\_\allowbreak salary\_\allowbreak account & Thông tin cá nhân & HỒ SƠ KH & bin & Có tài khoản nhận lương \\
n\_accounts & Thông tin cá nhân & HỒ SƠ KH & num & Số tài khoản đang mở \\
total\_\allowbreak avg\_\allowbreak balance & Thông tin cá nhân & HỒ SƠ KH & num & Tổng số dư bình quân 3 tháng \\
product\_\allowbreak code & Tham số khoản vay & ĐƠN VAY & cat & Mã sản phẩm tín dụng \\
requested\_\allowbreak amount & Tham số khoản vay & ĐƠN VAY & num & Số tiền đề nghị vay \\
tenor\_\allowbreak months & Tham số khoản vay & ĐƠN VAY & num & Kỳ hạn vay (tháng) \\
interest\_\allowbreak rate & Tham số khoản vay & ĐƠN VAY & num & Lãi suất áp dụng (năm) \\
loan\_\allowbreak purpose & Tham số khoản vay & ĐƠN VAY & cat & Mục đích vay \\
channel & Tham số khoản vay & ĐƠN VAY & cat & Kênh tiếp nhận hồ sơ \\
has\_\allowbreak collateral & Tham số khoản vay & ĐƠN VAY & bin & Có tài sản bảo đảm \\
monthly\_\allowbreak installment & Tham số khoản vay & ĐƠN VAY & num & Khoản trả góp hàng tháng \\
ltv & Tham số khoản vay & ĐƠN VAY & num & Tỷ lệ vay trên giá trị TSBĐ \\
apply\_hour & Tham số khoản vay & ĐƠN VAY & num & Giờ nộp hồ sơ \\
apply\_\allowbreak dayofweek & Tham số khoản vay & ĐƠN VAY & num & Ngày trong tuần nộp hồ sơ \\
cic\_score & Lịch sử tín dụng (CIC) & BẢN TIN CIC & num & Điểm tín dụng CIC (300-850) \\
current\_\allowbreak group & Lịch sử tín dụng (CIC) & BẢN TIN CIC & num & Nhóm nợ hiện tại (1-5) \\
worst\_\allowbreak group\_\allowbreak 12m & Lịch sử tín dụng (CIC) & BẢN TIN CIC & num & Nhóm nợ xấu nhất 12 tháng \\
dpd\_\allowbreak max\_\allowbreak 12m & Lịch sử tín dụng (CIC) & BẢN TIN CIC & num & Số ngày quá hạn tối đa 12 tháng \\
num\_\allowbreak institutions & Lịch sử tín dụng (CIC) & BẢN TIN CIC & num & Số TCTD đang có quan hệ \\
num\_\allowbreak active\_\allowbreak contracts & Lịch sử tín dụng (CIC) & BẢN TIN CIC & num & Số hợp đồng còn hiệu lực \\
num\_\allowbreak closed\_\allowbreak contracts & Lịch sử tín dụng (CIC) & BẢN TIN CIC & num & Số hợp đồng đã tất toán \\
total\_\allowbreak outstanding & Lịch sử tín dụng (CIC) & BẢN TIN CIC & num & Tổng dư nợ tại các TCTD \\
credit\_\allowbreak utilization & Lịch sử tín dụng (CIC) & BẢN TIN CIC & num & Tỷ lệ sử dụng hạn mức \\
monthly\_\allowbreak debt\_\allowbreak payment & Lịch sử tín dụng (CIC) & BẢN TIN CIC & num & Nghĩa vụ trả nợ hàng tháng hiện hữu \\
num\_\allowbreak inquiries\_\allowbreak 3m & Lịch sử tín dụng (CIC) & BẢN TIN CIC & num & Số lần truy vấn tín dụng 3 tháng \\
has\_\allowbreak written\_\allowbreak off\_\allowbreak debt & Lịch sử tín dụng (CIC) & BẢN TIN CIC & bin & Có nợ đã xử lý rủi ro \\
oldest\_\allowbreak account\_\allowbreak months & Lịch sử tín dụng (CIC) & BẢN TIN CIC & num & Độ dài lịch sử tín dụng (tháng) \\
num\_\allowbreak bad\_\allowbreak debt\_\allowbreak institutions & Lịch sử tín dụng (CIC) & BẢN TIN CIC & num & Số TCTD từng phát sinh nợ xấu \\
is\_\allowbreak thin\_\allowbreak file & Lịch sử tín dụng (CIC) & BẢN TIN CIC & bin & Khách hàng thiếu lịch sử tín dụng \\
cic\_\allowbreak missing & Lịch sử tín dụng (CIC) & BẢN TIN CIC & bin & Không lấy được dữ liệu CIC \\
txn\_\allowbreak cnt\_\allowbreak \{W\}d & Hành vi giao dịch & GIAO DỊCH & num & Số giao dịch trong cửa sổ W ngày \\
inflow\_\allowbreak \{W\}d & Hành vi giao dịch & GIAO DỊCH & num & Tổng dòng tiền vào trong cửa sổ W ngày \\
outflow\_\allowbreak \{W\}d & Hành vi giao dịch & GIAO DỊCH & num & Tổng dòng tiền ra trong cửa sổ W ngày \\
net\_\allowbreak flow\_\allowbreak \{W\}d & Hành vi giao dịch & GIAO DỊCH & num & Dòng tiền thuần trong cửa sổ W ngày \\
avg\_\allowbreak txn\_\allowbreak amt\_\allowbreak \{W\}d & Hành vi giao dịch & GIAO DỊCH & num & Giá trị giao dịch trung bình trong cửa sổ W ngày \\
txn\_\allowbreak amt\_\allowbreak std\_\allowbreak \{W\}d & Hành vi giao dịch & GIAO DỊCH & num & Độ lệch chuẩn giá trị giao dịch trong cửa sổ W ngày \\
salary\_\allowbreak cnt\_\allowbreak \{W\}d & Hành vi giao dịch & GIAO DỊCH & num & Số lần nhận lương trong cửa sổ W ngày \\
salary\_\allowbreak amt\_\allowbreak \{W\}d & Hành vi giao dịch & GIAO DỊCH & num & Tổng lương nhận trong cửa sổ W ngày \\
repay\_\allowbreak amt\_\allowbreak \{W\}d & Hành vi giao dịch & GIAO DỊCH & num & Tổng tiền trả nợ trong cửa sổ W ngày \\
risky\_\allowbreak amt\_\allowbreak \{W\}d & Hành vi giao dịch & GIAO DỊCH & num & Chi tiêu nhóm rủi ro trong cửa sổ W ngày \\
cash\_\allowbreak wd\_\allowbreak cnt\_\allowbreak \{W\}d & Hành vi giao dịch & GIAO DỊCH & num & Số lần rút tiền mặt trong cửa sổ W ngày \\
cash\_\allowbreak wd\_\allowbreak amt\_\allowbreak \{W\}d & Hành vi giao dịch & GIAO DỊCH & num & Giá trị rút tiền mặt trong cửa sổ W ngày \\
ecom\_\allowbreak amt\_\allowbreak \{W\}d & Hành vi giao dịch & GIAO DỊCH & num & Chi tiêu thương mại điện tử trong cửa sổ W ngày \\
utility\_\allowbreak amt\_\allowbreak \{W\}d & Hành vi giao dịch & GIAO DỊCH & num & Chi tiêu hóa đơn tiện ích trong cửa sổ W ngày \\
night\_\allowbreak txn\_\allowbreak cnt\_\allowbreak \{W\}d & Hành vi giao dịch & GIAO DỊCH & num & Số giao dịch ban đêm (0h-5h) trong cửa sổ W ngày \\
large\_\allowbreak txn\_\allowbreak cnt\_\allowbreak \{W\}d & Hành vi giao dịch & GIAO DỊCH & num & Số giao dịch giá trị lớn (>50 triệu) trong cửa sổ W ngày \\
mobile\_\allowbreak txn\_\allowbreak cnt\_\allowbreak \{W\}d & Hành vi giao dịch & GIAO DỊCH & num & Số giao dịch qua mobile banking trong cửa sổ W ngày \\
outflow\_\allowbreak ratio\_\allowbreak \{W\}d & Hành vi giao dịch & GIAO DỊCH & num & Tỷ lệ dòng tiền ra trên dòng tiền vào trong cửa sổ W ngày \\
risky\_\allowbreak ratio\_\allowbreak \{W\}d & Hành vi giao dịch & GIAO DỊCH & num & Tỷ trọng chi tiêu rủi ro trong cửa sổ W ngày \\
cash\_\allowbreak ratio\_\allowbreak \{W\}d & Hành vi giao dịch & GIAO DỊCH & num & Tỷ trọng rút tiền mặt trong cửa sổ W ngày \\
night\_\allowbreak ratio\_\allowbreak \{W\}d & Hành vi giao dịch & GIAO DỊCH & num & Tỷ trọng giao dịch ban đêm trong cửa sổ W ngày \\
repay\_\allowbreak burden\_\allowbreak \{W\}d & Hành vi giao dịch & GIAO DỊCH & num & Gánh nặng trả nợ trên dòng tiền vào trong cửa sổ W ngày \\
txn\_\allowbreak amt\_\allowbreak cv\_\allowbreak \{W\}d & Hành vi giao dịch & GIAO DỊCH & num & Hệ số biến thiên giá trị giao dịch trong cửa sổ W ngày \\
days\_\allowbreak since\_\allowbreak last\_\allowbreak txn & Hành vi giao dịch & GIAO DỊCH & num & Số ngày từ giao dịch gần nhất \\
txn\_\allowbreak cnt\_\allowbreak trend\_\allowbreak 30\_\allowbreak 90 & Hành vi giao dịch & GIAO DỊCH & num & Xu thế số giao dịch: 30 ngày so với mức trung bình 90 ngày \\
inflow\_\allowbreak trend\_\allowbreak 30\_\allowbreak 90 & Hành vi giao dịch & GIAO DỊCH & num & Xu thế dòng tiền vào: 30 ngày so với mức trung bình 90 ngày \\
inflow\_\allowbreak trend\_\allowbreak 90\_\allowbreak 180 & Hành vi giao dịch & GIAO DỊCH & num & Xu thế dòng tiền vào: 90 ngày so với mức trung bình 180 ngày \\
income\_\allowbreak verification\_\allowbreak ratio & Hành vi giao dịch & GIAO DỊCH & num & Tỷ lệ dòng tiền vào thực tế so với thu nhập khai báo \\
salary\_\allowbreak regularity & Hành vi giao dịch & GIAO DỊCH & num & Mức độ đều đặn của việc nhận lương trong 180 ngày \\
income\_\allowbreak stability & Hành vi giao dịch & GIAO DỊCH & num & Độ ổn định thu nhập (1 - hệ số biến thiên dòng tiền vào theo tháng) \\
dti & Biến tổ hợp & TÍNH TOÁN & num & Debt-\allowbreak to-\allowbreak Income: (trả góp mới + nghĩa vụ hiện hữu) / thu nhập \\
pti & Biến tổ hợp & TÍNH TOÁN & num & Payment-\allowbreak to-\allowbreak Income: trả góp khoản vay mới / thu nhập \\
amount\_\allowbreak to\_\allowbreak income & Biến tổ hợp & TÍNH TOÁN & num & Số tiền vay / thu nhập hàng tháng \\
outstanding\_\allowbreak to\_\allowbreak income & Biến tổ hợp & TÍNH TOÁN & num & Tổng dư nợ hiện hữu / thu nhập hàng tháng \\
disposable\_\allowbreak income & Biến tổ hợp & TÍNH TOÁN & num & Thu nhập còn lại sau trả nợ và chi phí sinh hoạt tối thiểu \\
disposable\_\allowbreak ratio & Biến tổ hợp & TÍNH TOÁN & num & Tỷ lệ thu nhập còn lại trên thu nhập \\
income\_\allowbreak per\_\allowbreak dependent & Biến tổ hợp & TÍNH TOÁN & num & Thu nhập bình quân đầu người trong hộ \\
credit\_\allowbreak history\_\allowbreak ratio & Biến tổ hợp & TÍNH TOÁN & num & Độ dài lịch sử tín dụng so với tuổi trưởng thành \\
debt\_\allowbreak per\_\allowbreak institution & Biến tổ hợp & TÍNH TOÁN & num & Dư nợ bình quân trên mỗi tổ chức tín dụng \\
dpd\_\allowbreak x\_\allowbreak institutions & Biến tổ hợp & TÍNH TOÁN & num & Tương tác giữa số ngày quá hạn và số TCTD \\
high\_\allowbreak dti\_\allowbreak short\_\allowbreak tenure & Biến tổ hợp & TÍNH TOÁN & bin & Cờ rủi ro: DTI > 0,5 và thâm niên công việc < 12 tháng \\
balance\_\allowbreak to\_\allowbreak installment & Biến tổ hợp & TÍNH TOÁN & num & Số dư bình quân / khoản trả góp hàng tháng \\
inflow\_\allowbreak to\_\allowbreak installment & Biến tổ hợp & TÍNH TOÁN & num & Dòng tiền vào 30 ngày / khoản trả góp hàng tháng \\
\end{longtable}

{\footnotesize\itshape Kiểu dữ liệu: num -- định lượng; cat -- định tính; bin -- nhị phân. Mỗi đặc trưng có hậu tố \code{\_\{W\}d} được tính trên 4 cửa sổ thời gian $W \in \{7, 30, 90, 180\}$ ngày; 23 dòng dạng này tương ứng 92 đặc trưng, đưa tổng số đặc trưng lên 154.}
\endgroup



\end{document}
